%\documentclass[aps,preprint]{revtex4}%

\documentclass[titlepage]{article}%
\usepackage{amsfonts}
\usepackage{achemso}
\usepackage{amsfonts}
\usepackage{amsmath}
\usepackage{amssymb}
\usepackage{bm}
\usepackage{graphicx}%
\setcounter{MaxMatrixCols}{30}
%TCIDATA{OutputFilter=latex2.dll}
%TCIDATA{Version=4.10.0.2363}
%TCIDATA{CSTFile=article.cst}
%TCIDATA{Created=Friday, August 29, 2003 10:08:38}
%TCIDATA{LastRevised=Wednesday, November 03, 2004 12:44:46}
%TCIDATA{<META NAME="GraphicsSave" CONTENT="32">}
%TCIDATA{<META NAME="DocumentShell" CONTENT="Articles\SW\REVTeX 4">}
%TCIDATA{Language=American English}
\newtheorem{theorem}{Theorem}
\newtheorem{acknowledgement}[theorem]{Acknowledgement}
\newtheorem{algorithm}[theorem]{Algorithm}
\newtheorem{axiom}[theorem]{Axiom}
\newtheorem{claim}[theorem]{Claim}
\newtheorem{conclusion}[theorem]{Conclusion}
\newtheorem{condition}[theorem]{Condition}
\newtheorem{conjecture}[theorem]{Conjecture}
\newtheorem{corollary}[theorem]{Corollary}
\newtheorem{criterion}[theorem]{Criterion}
\newtheorem{definition}[theorem]{Definition}
\newtheorem{example}[theorem]{Example}
\newtheorem{exercise}[theorem]{Exercise}
\newtheorem{lemma}[theorem]{Lemma}
\newtheorem{notation}[theorem]{Notation}
\newtheorem{problem}[theorem]{Problem}
\newtheorem{proposition}[theorem]{Proposition}
\newtheorem{remark}[theorem]{Remark}
\newtheorem{solution}[theorem]{Solution}
\newtheorem{summary}[theorem]{Summary}
\newenvironment{proof}[1][Proof]{\noindent\textbf{#1.} }{\ \rule{0.5em}{0.5em}}
\begin{document}

\title{Correlated One Particle Wave Functions}
\author{B. Weiner\\Department of Physics, Pennsylvania State University, DuBois PA 15801
\and J. V. Ortiz\\Department of Chemistry, Kansas State
University, Manhattan, KS 66506-3701}

\begin{abstract}
We introduce a one electron theory of multi electron systems that explicitly
describes electron - electron correlation. It is based on the observation that
Antisymmetrized Geminal Power (AGP) states are completly characterized by the
occupation numbers of their First Order Reduced Density Operators (FORDO)\ and
a set of Canonical General Spin Orbitals (CGSO's), which are in general
distinct from the natural general spin orbitals. Noting that the FORDO alone
does not determine an AGP\ state we give an explicit example of a Density
Matrix Functional.

We derive a generalization of the Independent Particle Fock operator and use
it to express the total state energy in terms of ioniztion potentials, kinetic
energies and potential energies of the CGSO's, which can thus be used to
modelling chemical processes and analyze chemical structure.

\end{abstract}
\date{11/01/2004}
\maketitle
\tableofcontents

\section{Introduction \label{intro}}

Independent Particle States (IPS's) have proved a mainstay of molecular
multi-particle quantum mechanics both in physical modelling and numerical
applications for many decades. The major attraction of these states lies in
the conceptually concrete and chemically understandable pictures of molecular
structure and properties that they provide. Apart from an initial choice of
$s$ atomic orbitals and possibly a decision about the molecular point group
and spin symmetry, one can obtain an unbiased description of the electronic
states of a molecule by finding an IPS\ that minimizes the energy via the Self
Consistent Field (SCF)\ equations. These equations provide a set of one
electron states $\left\{  \left.  \left\vert \varphi_{j}\right\rangle
\right\vert 1\leq j\leq2s\right\}  $, that are eigenstates of a self
consistent Fock operator corresponding to the eigenvalues $\left\{  \left.
\varepsilon_{j}\right\vert 1\leq j\leq2s\right\}  $ and are unique up to
unitary transformations that mix degenerate eigenstates. If the IPS\ describes
$M$ electrons then in general it is determined by the one-electron eigenstates
corresponding to the $M$ lowest eigenvalues $\left\{  \left.  \varepsilon
_{j}\right\vert 1\leq j\leq M\right\}  $. The wave functions of each of these
one electron states can then be analyzed in terms of a basis of atomic
orbitals, producing the familiar picture of one electron Probability
Distributions (PD's) labelled by one electron energies, that has been
invaluable in the understanding of chemical processes. The physical meaning of
the energies $\left\{  \left.  \varepsilon_{j}\right\vert 1\leq j\leq
M\right\}  $ can be deduced by considering the electron propagator, or by
directly examining the form of the Fock operator matrix elements, which show
that they correspond to ionization potentials.

It is often the case, however, that the IPS\ model is not sufficiently
accurate to predict or explain the finer details of chemical reaction
mechanisms or spectroscopic observations. In order to achieve this one must
take into account electron-electron correlation. Traditionally this is
accomplished by perturbational or configurational interaction methods that
increase numerical accuracy at the cost of a clear, tangible, chemical model
of the molecular system. The clarity is lost as pictorial models need be based
on one particle properties, but in general one particle properties do not
determine correlated multi-electron states. In particular one cannot use
chemical intuition to rearrange the positions of atoms and bonds and be
certain that the new conformation actually corresponds to a multi-electron
state, which is in contrast to the case of states determined by one particle properties.

However, we have described \cite{Weiner&ortiz04} a class of correlated states
that generalize the IPS model, explicitly contain a description of correlation
effects, and are the \emph{most general} states that produce a \emph{one
particle} theory of electrons i.e. multi-particle states are determined by one
particle properties. These are the Antisymmetrized Geminal Power (AGP) states
for an even number, $2N$, of electrons and the Generalized Antisymmetrized
Geminal Power (GAGP) ones for an odd number, $2N+1$. In the IPS\ case the
$2N$-electron state is completely determined by the occupation numbers of the
First Order Reduced Density Operator (FORDO) and a set of one particle
\emph{Natural }General Spin Orbitals (NGSOs). The occupation numbers in this
case are equal to one or zero. In the AGP\ case the $2N$-electron state is
completely determined by a set of one particle \emph{Canonical }General Spin
Orbitals (CGSOs) and occupation numbers of the FORDO, but now the occupation
numbers can be any value in the range $\left[  0,1\right]  ,$ though they must
be at least doubly degenerate. This theory \emph{explicitly contains}
correlation and generalizes the one-particle theory based on uncorrelated
Independent Particle States (IPS), which it contains as a special case. In
this paper we further show that one can explicitly express the total energy of
the $2N$-electron state in terms of the ionization potentials, kinetic
energies and potential energies of these CGSO's. Thus the AGP\ model produces
a simple pictorial interpretation in terms of occupation numbers and orbitals,
which totally \emph{determine} the $2N$-electron state and the CGSO's.

In contrast to IPS's when using AGP's it is important to note the distinction between

\begin{itemize}
\item AGP-NGSO's, which together with a set of doubly degenerate occupation
numbers, determine an AGP-FORDO, that correspond to a \emph{set }of
AGP\ states and

\item AGP-CGSO's, which together with a set of doubly degenerate occupation
numbers, determine a \emph{unique} AGP state.
\end{itemize}

The CGSO's are the NGSO's of the FORDO of a given AGP state, but the NGSO's of
the FORDO\ of this AGP state are \emph{not }in general the CSGO's of this
AGP\ state. In this paper when we are explicitly referring to CGSO's we will
use the symbol $\bm{\psi}$ and when we refer to NGSO's we will use the symbol
$\bm{\varphi.}$

The difference between NGSO's and CGSO's allows us to partition the manifold
of $2N$-electron AGP states into equivalence classes, $\left[  D^{1}\right]
$, indexed by FORDO's $D^{1}$ that are at least doubly degenerate, where each
element of a class can be parameterized by a vector of angles
\begin{equation}
\bm{\theta\,}=\left(  0,\theta_{2},\cdots,\theta_{s}\right)  ,\quad
0<\theta_{j}\leq2\pi;\quad2\leq j\leq s,
\end{equation}
and a representative FORDO $D^{1}.$ As the occupation numbers $\left\{
\left.  N_{j}\right\vert \quad1\leq j\leq s\right\}  $ of the $2N$-AGP\ state
$\left\vert g^{N}\right\rangle $ are in $1-1$ correspondence \cite{Yukalov00}
with the occupation numbers $\left\{  \left.  n_{j}\right\vert 1\leq j\leq
s\right\}  $ of the $2$-electron geminal state $\left\vert g\right\rangle $
each representative $D^{1}$ can, in turn, be parameterized by a vector of
geminal occupation numbers%
\begin{equation}
\bm{n}=\left(  n_{1},\cdots,n_{s}\right)  ,\quad0\leq n_{j}\leq1;\quad1\leq
j\leq s;\quad%
%TCIMACRO{\tsum \limits_{1\leq j\leq s}}%
%BeginExpansion
{\textstyle\sum\limits_{1\leq j\leq s}}
%EndExpansion
n_{j}=1
\end{equation}
and a representative set of NGSO's $\left\vert \bm{\varphi}\right\rangle
=\left\{  \left.  \left\vert \varphi_{j}\right\rangle \right\vert 1\leq
j\leq2s\right\}  $.

The energy functional $\mathcal{E}\left(  g^{N}\right)  $ over the manifold of
AGP\ states can thus be formulated as a functional $\mathcal{E}\left(
\bm{n},\bm{\theta,\varphi}\right)  $ specifically describing the phase and
FORDO $D^{1}\left(  \bm{n},\bm{\varphi}\right)  $ dependency. This leads to a
Density Matrix Functional (DMF), $\mathcal{F}$, defined on the set of
equivalence classes $\left\{  \left[  D^{1}\right]  \right\}  $ by\
\begin{equation}
\mathcal{F}\left(  D^{1}\right)  =\min_{\bm{\theta}}\left\{  \mathcal{E}%
\left(  \bm{n},\bm{\theta,\varphi}\right)  \right\}  ,
\end{equation}
which can be utilized in Density Matrix Functional Theory (DMFT) as a concrete
example of a DMF.

If we concentrate on the occupation and CGSO dependency of the energy we
denote the functional by $\mathcal{E}\left(  \bm{n},\bm{\psi}\right)  $, where
the geminal occupation numbers $\bm{n}$ must satisfy the constraints
\begin{subequations}
\begin{align}
n_{j}  &  \geq0;\quad1\leq j\leq s,\\
\sum_{1\leq j\leq s}n_{j}  &  =1,
\end{align}
and the CGSO's the orthonormality constraints%
\end{subequations}
\begin{equation}
\left\langle \left.  \psi_{j}\right\vert \psi_{k}\right\rangle =\delta
_{jk};\quad1\leq j,k\leq2s.
\end{equation}
Variation of the energy with respect to (wrt) orbital changes described by a
Taylors series expansion involve functional derivatives of $\mathcal{E}\left(
\bm{n},\bm{\psi}\right)  $ with respect to $\left\{  \left.  \psi
_{j}\right\vert 1\leq j\leq2s\right\}  .$ We show that these derivatives can
be expressed in terms of a generalized Fock operator $F\left(
\bm{n},\bm{\psi}\right)  $, which is \emph{only self adjoint} at stationary
points $\bm{\psi}_{s}$ of the orbital variation. Further if $\bm{n}_{s}$ is a
stationary point of the occupation number variation then the diagonal elements
of the matrix of $F\left(  \bm{n}_{s},\bm{\psi}_{s}\right)  $ in the CGSO
basis $\left\{  \left\vert \left\vert \psi_{j}\right\rangle \right\vert 1\leq
j\leq2s\right\}  $ are ionization potentials associated with these orbitals
weighted by their occupation numbers $\left\{  N_{j}\right\}  $.

Unlike the IPS case the Fock operator is in general \emph{not} diagonal in the
basis $\left\{  \left.  \left\vert \psi_{j}\right\rangle \right\vert 1\leq
j\leq2s\right\}  $. It is, however, possible to diagonalize $F\left(
\bm{n}_{s},\bm{\psi}_{s}\right)  $ in terms of eigenvectors $\left\{  \left.
\left\vert \chi_{j}\right\rangle \right\vert 1\leq j\leq2s\right\}  $ and
eigenvalues $\left\{  \left.  \varepsilon_{j}\right\vert 1\leq j\leq
2s\right\}  $ such that $\left\{  \varepsilon_{j}\right\}  $ are ionization
potentials associated with the orbitals $\left\{  \left\vert \chi
_{j}\right\rangle \right\}  $, but then the geminal $\left\vert g\left(
\bm{n}_{s},\bm{\psi}_{s}\right)  \right\rangle ,$ that determines the energy
optimized AGP\ $\left\vert g\left(  \bm{n}_{s},\bm{\psi}_{s}\right)
^{N}\right\rangle ,$ will \emph{not }be in a canonical form in terms of the
basis $\left\{  \left.  \left\vert \chi_{j}\right\rangle \right\vert 1\leq
j\leq2s\right\}  .$

One can avoid the use of constrained domains for the energy functional by
utilizing the unitary group of one electron Hilbert space $\mathcal{H}^{1},$
parameterized by $\left(  \bm{\lambda},\bm{\theta}\right)  ,$ to vary the
orbitals and a commutative nonunitary group of positive transformations,
parameterized by $\bm{\eta},$ to vary the occupation numbers. This allows us
to express the energy as a function $\mathcal{E}\left(
\bm{\eta },\bm{\lambda},\bm{\theta}\right)  $ where there are no constraints
on $\left(  \bm{\eta},\bm{\lambda},\bm{\theta}\right)  .$ Moreover the unitary
transformations produce NGSO's $\bm{\varphi=}$ $\bm{\varphi}\left(
\bm{\lambda}\right)  $ that belong to different equivalence classes $\left[
D^{1}\right]  $ and $\bm{\theta}$ produces transformations within an
equivalence class by altering the phases of $\bm{\varphi.}$

The one electron model that we obtain describes molecules by weighted PD's,
each determined by an occupation number and a phase optimized GSO associated
with an ionization potential, kinetic and potential energies in an analogous
fashion to the IPS case, but also contains correlation effects. Pictorial
manipulations of these PD's (taking into account relaxation of the ionization
energies, kinetic and potential energies and optimal phases) allows one to
analyze chemical processes and structures in a rigorously meaningful fashion,
as one is assured that the pictures produced corresponds to multi electron
states and thus can be used as an effective and rigorous chemical modelling tool.

The AGP\ state is a special case of a Multi Configurational (MC) state and we
identify many features optimized AGP\ and MC states have in common, such as
the decomposition of the total energy in terms of GSO quantities, while at the
same time focusing on the crucial distinctive property of AGP\ states that
they are totally determined by one particle properties.

The structure of this article is as follows: in section \ref{AGPstate} we
describe the canonical form of a geminal, and the AGP\ state, in terms of
geminal occupation numbers $\left\{  \left.  n_{j}\right\vert 1\leq j\leq
s\right\}  $ and CGSOs $\left\{  \left.  \left\vert \psi_{j}\right\rangle
\right\vert 1\leq j\leq2s\right\}  .$ Together with the AGP-FORDO occupation
numbers $\left\{  \left.  N_{j}\right\vert 1\leq j\leq s\right\}  ,$ which are
in $1-1$ correspondence with the geminal occupation numbers, these CGSOs
determine both the AGP\ state and the AGP-FORDO. In section \ref{transform} we
examine the $1-1$ transformation between geminal $\left\{  \left.
n_{j}\right\vert 1\leq j\leq s\right\}  $ and AGP $\left\{  \left.
N_{j}\right\vert 1\leq j\leq s\right\}  $ occupation numbers. In section
\ref{Energy} we discuss the energy expressions in terms of the occupation
numbers $\left\{  \left.  n_{j}\right\vert 1\leq j\leq s\right\}  $, CGSOs
$\left\{  \left.  \left\vert \psi_{j}\right\rangle \right\vert 1\leq
j\leq2s\right\}  $, NGSOs $\left\{  \left.  \left\vert \varphi_{j}%
\right\rangle \right\vert 1\leq j\leq2s\right\}  $ and phases $\left\{
\left.  \theta_{j}\right\vert 1\leq j\leq s\right\}  .$ In section
\ref{Evariation} we describe the necessary conditions needed to minimize the
energy over a manifold of AGP\ states. In section \ref{Fock} we derive a
generalized Fock operator that produces the derivative of the energy
functional with respect to NGSO's for fixed FORDO occupation numbers. In this
section we also analyze the total energy in terms of this Fock operator and
show that the total state energy can be explicitly expressed in terms of one
particle quantities. In section \ref{IPS} we examine how our generalized Fock
operator is related to the standard IPS\ Fock operator, by constraining the
AGP state to be an IPS and in section \ref{Discussion} we conclude with a
review of the material presented and a discussion of its implications.

\section{The $2N$-electron AGP state\label{AGPstate}}

In this section we review the canonical form of a geminal and AGP state and
their associated first and second order reduced density operators. We also
characterize equivalence classes of AGP states that have the same FORDO.

\subsection{Canonical Form}

The $2N$-electron AGP state, $\left\vert g^{N}\right\rangle $, is the $N^{th}$
antisymmetric power of the $2$-electron state described by the geminal
\begin{equation}
\left\vert g\right\rangle =\sum_{1\leq j\leq s}n_{j}^{\frac{1}{2}}\left\vert
\psi_{j}\psi_{j+s}\right\rangle , \label{geminal}%
\end{equation}
where $\left\{  \left.  n_{j}^{\frac{1}{2}}\right\vert 1\leq j\leq s\right\}
$ are the positive square roots of $\left\{  \left.  n_{j}\right\vert 1\leq
j\leq s\right\}  ,$ and is given by \cite{Weiner&ortiz04}
\begin{equation}
\left\vert g^{N}\right\rangle =\left(  S_{N}\right)  ^{-\frac{1}{2}}%
\sum_{1\leq j_{1}<\cdots<j_{N}\leq s}n_{j_{1}}^{\frac{1}{2}}\cdots n_{j_{s}%
}^{\frac{1}{2}}\left\vert \psi_{j_{1}}\psi_{j_{1}+s}\cdots\psi_{j_{N}}%
\psi_{j_{N}+s}\right\rangle . \label{agp}%
\end{equation}
The normalization factor, $\left(  S_{N}\right)  ^{-\frac{1}{2}}$, is given in
terms of the elementary symmetric function
\begin{equation}
S_{N}=\sum_{1\leq j_{1}<\cdots<j_{N}\leq s}n_{j_{1}}\cdots n_{j_{N}},
\end{equation}
and the geminals $\left\vert g\right\rangle $ represent normalized two
electron states so that
\begin{equation}
\sum_{1\leq j\leq s}n_{j}=1.
\end{equation}
It is evident from the expansion eq. $\left(  \ref{geminal}\right)  $ that the
geminal and geminal power state are at the least invariant to the group
\begin{equation}
SU\left(  2\right)  ^{s}=\overset{s\text{-factors}}{\overbrace{\,SU\left(
2\right)  \times\cdots\times SU\left(  2\right)  }},
\end{equation}
that is formed by $2\times2$ \emph{special} unitary transformations of the
paired CGSO's $\left\{  \left.  \left\vert \psi_{j}\right\rangle ,\left\vert
\psi_{j+s}\right\rangle \right\vert 1\leq j\leq s\right\}  ,$ while if
degeneracies exist amongst the $\left\{  \left.  n_{j}\right\vert 1\leq j\leq
s\right\}  $ this invariance group is larger \cite{AGPCoherent}. The canonical
CSGO's are unique up to transformations belonging this invariance group.

The CGSO's $\left\{  \left.  \left\vert \psi_{j}\right\rangle \right\vert
1\leq j\leq2s\right\}  $ can be expanded in terms of a spin orbitals basis
$\left\{  \left.  \left\vert \varsigma_{j\alpha}\right\rangle ,\left\vert
\varsigma_{j\beta}\right\rangle \right\vert 1\leq j\leq s\right\}  $ of one
electron state space $\mathcal{H}^{1}$ as
\begin{equation}
\left\vert \psi_{j}\right\rangle =\sum_{1\leq j\leq s}\left\{  a_{j}\left\vert
\varsigma_{j\alpha}\right\rangle +b_{j}\left\vert \varsigma_{j\beta
}\right\rangle \right\}  ,
\end{equation}
leading to one electron GSO wavefunctions defined on the spectrum of position
and spin operators as
\begin{equation}
\psi_{j}\left(  \bm{r,}\zeta\right)  =\left\langle \left.  \bm{r,}\zeta
\right\vert \psi_{j}\right\rangle .
\end{equation}
It should be noted that if expansion eq. $\left(  \ref{geminal}\right)  $ is
used to define the canonical expansion it is usually \emph{not possible} to
find a set of real CSGO's wavefunctions.

\subsection{First Order Reduced Density Operators\label{FORDO}}

The states $\left\vert g\right\rangle $ and $\left\vert g^{N}\right\rangle $
define FORDO's given by \cite{Coleman97}
\begin{equation}
D^{1}\left(  g\right)  =\sum_{1\leq j\leq s}n_{j}\left\{  \left\vert \psi
_{j}\right\rangle \left\langle \psi_{j}\right\vert +\left\vert \psi
_{j+s}\right\rangle \left\langle \psi_{j+s}\right\vert \right\}  \label{d1g}%
\end{equation}
and
\begin{equation}
D^{1}\left(  g^{N}\right)  =\sum_{1\leq j\leq s}N_{j}\left\{  \left\vert
\psi_{j}\right\rangle \left\langle \psi_{j}\right\vert +\left\vert \psi
_{j+s}\right\rangle \left\langle \psi_{j+s}\right\vert \right\}  \label{dng}%
\end{equation}
respectively. It is clear from the expansions eqs. $\left(  \ref{d1g}\right)
$ and $\left(  \ref{dng}\right)  $ that the CGSO's $\left\{  \left.
\left\vert \psi_{j}\right\rangle \right\vert 1\leq j\leq r\right\}  $ are
simultaneously the Natural General Spin Orbitals (NGSO's) of $D^{1}\left(
g\right)  $ and $D^{1}\left(  g^{N}\right)  $. However it is very important to
note that the NGSO's of $D^{1}\left(  g\right)  $ and $D^{1}\left(
g^{N}\right)  $ \emph{are not }in general the CGSO's of $\left\vert
g\right\rangle $. This is a ramification of

\begin{enumerate}
\item the general invariance of FORDO's to the multiplication of their NGSO's
by complex phase factors
\begin{equation}
D^{1}=\sum_{1\leq j\leq r}\nu_{j}\left\{  \left\vert \chi_{j}\right\rangle
\left\langle \chi_{j}\right\vert \right\}  =\sum_{1\leq j\leq r}\nu
_{j}\left\{  e^{i\theta_{j}}\left\vert \chi_{j}\right\rangle \left\langle
\chi_{j}\right\vert e^{-i\theta_{j}}\right\}  ,
\end{equation}
where $\left\{  \left.  \nu_{j}\right\vert 1\leq j\leq s\right\}  $ are
occupation numbers and

\item the observation that geminals \emph{are not }in general invariant to
such phase changes of their CGSO's.
\end{enumerate}

One can form equivalence classes $\left[  D^{1}\left(  g\left(  \bm{0}\right)
\right)  \right]  $ of geminals $\left\{  \left\vert g\left(
\bm{\theta}\right)  \right\rangle \right\}  $ that have the same FORDO
$D^{1}\left(  g\left(  \bm{0}\right)  \right)  $ and equivalence classes
$\left[  D^{1}\left(  g^{N}\left(  \bm{0}\right)  \right)  \right]  $ of
$2N$-electron AGP states $\left\{  \left\vert g^{N}\left(
\bm{\theta }\right)  \right\rangle \right\}  $, that have the same FORDO
$D^{1}\left(  g^{N}\left(  \bm{0}\right)  \right)  $. For a fixed $N$ these
classes correspond to each other in a bijective manner \cite{Yukalov00} and
can be parameterized by the angles $\bm{\theta}$, where
\begin{align}
\left\vert g\left(  \bm{\theta}\right)  \right\rangle  &  =\sum_{1\leq j\leq
s}n_{j}^{\frac{1}{2}}e^{i\theta_{j}}\left\vert \varphi_{j}\varphi
_{j+s}\right\rangle =\sum_{1\leq j\leq s}n_{j}^{\frac{1}{2}}\left\vert
e^{i\alpha_{j}}\varphi_{j}e^{i\alpha_{j+s}}\varphi_{j+s}\right\rangle
;\nonumber\\
&  \qquad\qquad\qquad\qquad\qquad\qquad\qquad\theta_{1}=0,\,0<\theta_{j}%
\leq2\pi,\,2\leq j\leq s,
\end{align}
where $\bm{\theta=}\left(  0,\theta_{2},\cdots,\theta_{s}\right)  $ and
$\alpha_{j}+\alpha_{j+s}=\theta_{j}$. Choosing $\theta_{1}=0$ eliminates the
overall phase factor indeterminacy of $\left\vert g\left(
\bm{\theta }\right)  \right\rangle $. The individual phase factors $\left\{
\left.  \alpha_{j}\right\vert 1\leq j\leq s\right\}  $ are not unique, though
their sums $\left\{  \left.  \theta_{j}\right\vert 1\leq j\leq s\right\}  $
are, allowing one to obtain a unique representative set $\left\{  \left.
\alpha_{j}\right\vert 1\leq j\leq s\right\}  $ by defining
\begin{equation}
\alpha_{j}=\frac{\theta_{j}}{2},\alpha_{j+s}=\frac{\theta_{j}}{2};\quad1\leq
j\leq s.
\end{equation}
Thus one has the correspondences
\begin{equation}
\left\{  \left\vert g\left(  \bm{\theta}\right)  \right\rangle
\longleftrightarrow\left\vert g^{N}\left(  \bm{\theta}\right)  \right\rangle
;\quad\theta_{1}=0;\,0<\theta_{j}\leq2\pi,2\leq j\leq s\right\}
\longleftrightarrow\left[  D^{1}\left(  g^{N}\left(  \bm{0}\right)  \right)
\right]  .
\end{equation}
Each of the elements of the class $\left[  D^{1}\left(  g^{N}\left(
\bm{0}\right)  \right)  \right]  $ can also be expressed in canonical form as
\begin{equation}
\sum_{1\leq j\leq s}n_{j}^{\frac{1}{2}}\left\vert \chi_{j}\chi_{j+s}%
\right\rangle ,
\end{equation}
where $\left\vert \chi_{j}\right\rangle =\left\vert e^{i\varphi_{j}}%
\varphi_{j}\right\rangle $.

\subsection{Second Order Reduced Density Operators\label{SORDO}}

The AGP\ states $\left\vert g^{N}\left(  \bm{\theta}\right)  \right\rangle $
are clearly distinct as they have different Second Order Reduced Density
Operators (SORDO's), $D^{2}\left(  g^{N}\left(  \bm{\theta}\right)  \right)
$, given by \cite{Yukalov00}\cite{Coleman97}
\begin{align}
&  D^{2}\left(  g^{N}\left(  \bm{\theta}\right)  \right)  =\frac{1}{S_{N}%
}\left\{  \sum_{1\leq j\leq s}n_{j}\left\vert \varphi_{j}\varphi
_{j+s}\right\rangle \left\langle \varphi_{j}\varphi_{j+s}\right\vert \right.
+\nonumber\\
&  \left.  \sum_{1\leq j,k\leq s}n_{j}^{\frac{1}{2}}n_{k}^{\frac{1}{2}%
}e^{i\left(  \theta_{j}-\theta_{k}\right)  }S_{N-1}\left(  \jmath k\right)
\left\vert \varphi_{j}\varphi_{j+s}\right\rangle \left\langle \varphi
_{k}\varphi_{k+s}\right\vert +\sum_{1\leq j<k\leq2s}n_{j}n_{k}S_{N-2}\left(
\jmath k\right)  \left\vert \varphi_{j}\varphi_{k}\right\rangle \left\langle
\varphi_{j}\varphi_{k}\right\vert \right\}  , \label{D2}%
\end{align}
where
\begin{equation}
S_{N-1}\left(  \jmath k\right)  =\sum_{\substack{1\leq j_{1}<\cdots
<j_{N-1}\leq s\\j,k\notin\left\{  j_{1},\cdots,j_{N-1}\right\}  }}n_{j_{1}%
}\cdots n_{j_{N-1}}%
\end{equation}
and
\begin{equation}
S_{N-2}\left(  \jmath k\right)  =\sum_{\substack{1\leq j_{1}<\cdots
<j_{N-2}\leq s\\j,k\notin\left\{  j_{1},\cdots,j_{N-2}\right\}  }}n_{j_{1}%
}\cdots n_{j_{N-2}}.
\end{equation}
Note that we define
\begin{align}
S_{N-1}\left(  \jmath k+s\right)   &  =S_{N-1}\left(  \jmath+sk\right)
=S_{N-1}\left(  \jmath+sk+s\right)  =S_{N-1}\left(  kj\right)  =S_{N-1}\left(
\jmath k\right) \nonumber\\
S_{N-2}\left(  \jmath k+s\right)   &  =S_{N-2}\left(  \jmath+sk\right)
=S_{N-2}\left(  \jmath+sk+s\right)  =S_{N-2}\left(  kj\right)  =S_{N-2}\left(
\jmath k\right) \nonumber\\
S_{N-1}\left(  \jmath\jmath\right)   &  =S_{N-2}\left(  \jmath\jmath\right)
=0,
\end{align}
for $1\leq j,k\leq s$. The SORDO eq. $\left(  \ref{D2}\right)  $ can be
expanded in terms of $\bm{\theta}$ dependent NGSO's $\left\{  \left.
\left\vert e^{i\alpha_{j}}\varphi_{j}\right\rangle ,\left\vert e^{i\alpha
_{j+s}}\varphi_{j+s}\right\rangle \right\vert 1\leq j\leq s\right\}  ,$ where
$\left\{  \left.  \alpha_{j}+\alpha_{j+s}=\theta_{j}\right\vert 1\leq j\leq
s\right\}  $ giving the explicitly orbital phase dependent expression
\begin{align}
&  D^{2}\left(  g^{N}\left(  \bm{\theta}\right)  \right)  =\frac{1}{S_{N}%
}\left\{  \sum_{1\leq j\leq s}n_{j}S_{N-1}\left(  \jmath\right)  \left\vert
e^{i\alpha_{j}}\varphi_{j}e^{i\alpha_{j+s}}\varphi_{j+s}\right\rangle
\left\langle e^{i\alpha_{j}}\varphi_{j}e^{i\alpha_{j+s}}\varphi_{j+s}%
\right\vert \right. \nonumber\\
&  \qquad\qquad\qquad\qquad+\sum_{1\leq j,k\leq s}n_{j}^{\frac{1}{2}}%
n_{k}^{\frac{1}{2}}S_{N-1}\left(  \jmath k\right)  \left\vert e^{i\alpha_{j}%
}\varphi_{j}e^{i\alpha_{j+s}}\varphi_{j+s}\right\rangle \left\langle
e^{i\alpha_{k}}\varphi_{k}e^{i\alpha_{k+s}}\varphi_{k+s}\right\vert
\nonumber\\
&  \qquad\qquad\qquad\qquad\left.  +\sum_{1\leq j<k\leq2s}n_{j}n_{k}%
S_{N-2}\left(  \jmath k\right)  \left\vert e^{i\alpha_{j}}\varphi
_{j}e^{i\alpha_{k}}\varphi_{k}\right\rangle \left\langle e^{i\alpha_{j}%
}\varphi_{j}e^{i\alpha_{k}}\varphi_{k}\right\vert \right\}  . \label{en2}%
\end{align}
In terms of CGSO's $\left\{  \psi_{j}\right\}  $ the SORDO can be expressed in
phase free form as
\begin{align}
&  D^{2}\left(  g^{N}\right)  =\frac{1}{S_{N}}\left\{  \sum_{1\leq j\leq
s}n_{j}\left\vert \psi_{j}\psi_{j+s}\right\rangle \left\langle \psi_{j}%
\psi_{j+s}\right\vert \right.  +\nonumber\\
&  \left.  \sum_{1\leq j,k\leq s}n_{j}^{\frac{1}{2}}n_{k}^{\frac{1}{2}}%
S_{N-1}\left(  \jmath k\right)  \left\vert \psi_{j}\psi_{j+s}\right\rangle
\left\langle \psi_{k}\psi_{k+s}\right\vert +\sum_{1\leq j<k\leq2s}n_{j}%
n_{k}S_{N-2}\left(  \jmath k\right)  \left\vert \psi_{j}\psi_{k}\right\rangle
\left\langle \psi_{j}\psi_{k}\right\vert \right\}  . \label{D2C}%
\end{align}


\section{The correspondence between AGP and Geminal Occupation
Numbers\label{transform}}

Any set of $s$ positive numbers $\left\{  \left.  N_{j}\right\vert 1\leq j\leq
s\right\}  $ that have values between $0$ and $1$ and add up to $N,$ can be
interpreted as the occupation numbers of a FORDO that is at least doubly
degenerate and can be associated with an AGP\ state. These numbers are in
$1-1$ correspondence with another set of positive numbers $\left\{  \left.
n_{j}\right\vert 1\leq j\leq s\right\}  ,$ that add up to $1$, which can be
interpreted as the occupation numbers of a FORDO\ of a geminal. The geminal
occupation numbers $\left\{  \left.  n_{j}\right\vert 1\leq j\leq s\right\}  $
together with an arbitrary set of $2s$ orthonormal one electron Canonical
General Spin Orbitals (CGSO's) then \emph{uniquely }determine a $2N$ electron
AGP state. Any state that has a FORDO that has at least doubly degenerate
occupation numbers is covered by a set of AGP\ states i.e. there exists
AGP\ states that have the same FORDO. In most applications e.g. calculating
the expectation values of observables, it is sufficient to know the $2N$
electron state occupation numbers in terms of the geminal occupation numbers,
which serve as a parameterization of the state occupation numbers.

The vector of occupation numbers $\mathbf{N}=\left(  N_{1},\cdots
,N_{s}\right)  $ of an AGP\ state $\left\vert g^{N}\right\rangle $ can be
expressed as a bijective function, $\mathcal{N}\left(  \mathbf{n}\right)  $,
of the geminal occupation numbers $\mathbf{n}=\left(  n_{1},\cdots
,n_{s}\right)  ,$ where
\begin{equation}
\mathcal{N}:\mathfrak{D}_{N}\subset\mathbb{R}^{s}\rightarrow\mathfrak{D}%
_{n}\subset\mathbb{R}^{s}.
\end{equation}
The domain $\mathfrak{D}_{N}$ and range $\mathfrak{D}_{n}$ are given by%
\begin{align}
\mathfrak{D}_{N}  &  =\left(  N_{1},\cdots,N_{s}\right)  ;\quad0\leq N_{j}%
\leq1,\quad1\leq j\leq s\nonumber\\
\sum_{1\leq j\leq s}N_{j}  &  =N
\end{align}
and
\begin{align}
\mathfrak{D}_{n}  &  =\left(  n_{1},\cdots,n_{s}\right)  ;\quad n_{j}%
\geq0,\quad1\leq j\leq s\nonumber\\
\sum_{1\leq j\leq s}n_{j}  &  =1.
\end{align}
The function $\mathcal{N}$ is defined as
\begin{equation}
N_{j}=\mathcal{N}_{j}\left(  \mathbf{n}\right)  =n_{j}\frac{S_{N-1}\left(
\jmath\right)  }{S_{N}};\quad1\leq j\leq s,
\end{equation}
where
\begin{equation}
S_{N-1}\left(  \jmath\right)  =\sum_{\substack{1\leq j_{1}<\cdots<j_{N-1}\leq
s\\j\notin\left\{  j_{1},\cdots,j_{N-1}\right\}  }}n_{j_{1}}\cdots n_{j_{N-1}%
}.
\end{equation}


\section{Energy Functionals \label{Energy}}

In this section we describe the energy of a $2N$-electron AGP state in three
different ways in terms of:

\begin{enumerate}
\item CGSO's and geminal occupation numbers.

\item A density matrix functional $\mathcal{F}\left(  D^{1}\right)  $ defined
on the equivalence classes $\left[  D^{1}\left(  g^{N}\left(  \bm{0}\right)
\right)  \right]  .$ The functional $\mathcal{F}\left(  D^{1}\right)  $ is
significant in two aspects, it can serve as a rigorous explicitly known model
functional in Density Matrix Functional Theory (DMFT) \cite{Cioslowski01}%
\cite{Staroverov02}\cite{Cioslowski02} and it clearly demonstrates the
difference between CGSO's and NGSO's.

\item Group parameters discussed in \cite{AGPCoherent}that
describe the manifold of AGP\ states. This parameterization allows
one to constructively separate the variation of the phases
$\left\{  \theta_{j}\right\}  $, NGSO's $\left\{  \left\vert
\varphi_{j}\right\rangle \right\}  $ and occupation numbers
$\left\{  n_{j}\right\}  $, which is necessary in order to
construct the functionals $\mathcal{F}\left(  D^{1}\right)  $. It
is also important in the analysis of the properties of the
generalized Fock operator associated with optimized AGP\ states
and the derivation of a generalized Brillouin condition.
\end{enumerate}

The supporting details can be found in Appendix \ref{En_Funcs}, which reviews
the form of the total energy expression for any $2N$-electron state in terms
of one electron unitary group parameters, GSO's and SORDM's.

\subsection{Energy in terms of Geminal Occupation Numbers and
CGSO's\label{EnON_CGSO}}

The energy, $\mathcal{E}\left(  \bm{n,\psi}\right)  ,$ of a $2N$-electron
AGP\ state for a molecule at a fixed nuclear configuration $\left(
\bm{R}_{1},\cdots,\bm{R}_{N_{nuc}}\right)  $ in terms of the CGSO's amplitudes
$\left\{  \left.  \left\langle \left.  \bm{r,\xi}\right\vert \psi
_{j}\right\rangle \right\vert 1\leq j\leq r\right\}  $ and geminal occupation
numbers $\bm{n}$ is
\begin{align}
\mathcal{E}\left(  \bm{n,\psi}\right)   &  =\frac{\mathcal{H}\left(
\bm{n,\psi}\right)  }{S_{N}\left(  \bm{n}\right)  }\nonumber\\
&  =\frac{1}{S_{N}\left(  \bm{n}\right)  }\left\{  \sum_{1\leq j\leq s}%
n_{j}S_{N-1}\left(  \jmath\right)  \left\{  \left\langle \psi_{j}\left\vert
h_{1}\psi_{j}\right.  \right\rangle +\left\langle \psi_{j+s}\left\vert
h_{1}\psi_{j+s}\right.  \right\rangle +\left\langle \psi_{j}\psi_{j+s}%
||\psi_{j}\psi_{j+s}\right\rangle \right\}  \right. \nonumber\\
&  +\sum_{1\leq j,k\leq s}n_{j}^{\frac{1}{2}}n_{k}^{\frac{1}{2}}S_{N-1}\left(
\jmath k\right)  \left\langle \psi_{j}\psi_{j+s}||\psi_{k}\psi_{k+s}%
\right\rangle \left.  +\sum_{1\leq j<k\leq s}n_{i}n_{j}S_{N-2}\left(  \jmath
k\right)  \left\langle \psi_{j}\psi_{k}|||\psi_{j}\psi_{k}\right\rangle
\right\}  . \label{energy}%
\end{align}
$\lceil$Note: Mathematically the energy functional should really be denoted as
$\mathcal{E}\left(  \bm{n,\psi}^{\ast}\bm{,\psi}\right)  $ as $\bm{\psi}^{\ast
}$ and $\bm{\psi}$ are treated as independent variables even though
$\bm{\psi}$ determines $\bm{\psi}^{\ast}$.$\rfloor$ The one-electron matrix
elements $\left\{  \left\langle \psi_{j}\left\vert h_{1}\psi_{k}\right.
\right\rangle \right\}  $ wrt the CGSO's are
\begin{equation}
\left\langle \psi_{j}\left\vert h_{1}\psi_{k}\right.  \right\rangle =-\int
\psi_{j}^{\ast}\left(  \bm{r,\xi}\right)  \left\{  \sum_{_{1\leq l\leq
N_{nuc}}}\frac{Z_{l}}{\left\Vert \bm{R}_{l}-\bm{r}\right\Vert }+\nabla
_{\bm{r}}^{2}\right\}  \psi_{k}\left(  \bm{r,\xi}\right)  d\bm{r}d\xi
;\quad1\leq j,k\leq2s, \label{h1}%
\end{equation}
the antisymmetric two-electron matrix elements $\left\{  \left\langle
jk||lm\right\rangle \right\}  $
\begin{align}
\left\langle jk||lm\right\rangle  &  =\int\left\{  \psi_{j}^{\ast}\left(
\bm{r}_{1}\bm{,\xi}_{1}\right)  \psi_{k}^{\ast}\left(  \bm{r}_{2}%
\bm{,\xi}_{2}\right)  \frac{1}{\left\Vert \bm{r}_{1}-\bm{r}_{2}\right\Vert
}\psi_{l}\left(  \bm{r}_{1}\bm{,\xi}_{1}\right)  \psi_{m}\left(
\bm{r}_{2}\bm{,\xi}_{2}\right)  \right. \nonumber\\
&  -\left.  \psi_{j}^{\ast}\left(  \bm{r}_{1}\bm{,\xi}_{1}\right)  \psi
_{k}^{\ast}\left(  \bm{r}_{2}\bm{,\xi}_{2}\right)  \frac{1}{\left\Vert
\bm{r}_{1}-\bm{r}_{2}\right\Vert }\psi_{l}\left(  \bm{r}_{2}\bm{,\xi}_{2}%
\right)  \psi_{m}\left(  \bm{r}_{1}\bm{,\xi}_{1}\right)  \right\}
d\bm{r}_{1}d\bm{r}_{2}d\xi_{1}d\xi_{2}\nonumber\\
&  \left.  {}\right.  \qquad\qquad\qquad\qquad\qquad\qquad\qquad\qquad
\qquad1\leq j,k\leq2s;1\leq l,m\leq2s \label{h2}%
\end{align}
and $\left\{  \left\langle jk|||jk\right\rangle \right\}  $
\begin{align}
\left\langle \psi_{j}\psi_{k}|||\psi_{j}\psi_{k}\right\rangle  &
=\left\langle jk||jk\right\rangle +\left\langle j\,k+s||j\,k+s\right\rangle
+\left\langle j+s\,k||j+s\,k\right\rangle +\left\langle
j+s\,k+s||j+s\,k+s\right\rangle \nonumber\\
&  \left.  {}\right.  \qquad\qquad\qquad\qquad\qquad\qquad\qquad\qquad
\qquad\qquad\qquad\qquad\qquad\qquad1\leq j,k\leq s.
\end{align}
We shall often express $\left\langle jk||lm\right\rangle $ in terms of
nonantisymmetric matrix elements $\left\{  \left(  \left.  jk\right\vert
h_{2}lm\right)  \right\}  $ as
\begin{equation}
\left\langle jk||lm\right\rangle =\left(  \left.  jk\right\vert h_{2}%
lm\right)  -\left(  \left.  jk\right\vert h_{2}ml\right)  ,
\end{equation}
where
\begin{equation}
\left(  \left.  jk\right\vert h_{2}lm\right)  =\int\psi_{j}^{\ast}\left(
\bm{x}_{1}\right)  \psi_{k}^{\ast}\left(  \bm{x}_{2}\right)  \frac
{1}{\left\Vert \bm{r}_{1}-\bm{r}_{2}\right\Vert }\psi_{l}\left(
\bm{x}_{1}\right)  \psi_{m}\left(  \bm{x}_{2}\right)  d\bm{x}_{1}d\bm{x}_{2}%
\end{equation}
and use the notation
\begin{align}
V_{2jklm}  &  =\left\langle jk||lm\right\rangle \nonumber\\
h_{2jklm}  &  =\left(  \left.  jk\right\vert h_{2}lm\right)  .
\end{align}


\subsection{Energy in terms of NGSO's, Phases and FORDO's\label{EnON_NGSO}}

In terms of $\theta$ dependent NGSO's $\left\{  \left\vert e^{i\alpha_{j}%
}\varphi_{j}\right\rangle ;1\leq j\leq2s\right\}  $ defined in eq. $\left(
\ref{en2}\right)  $ we obtain an energy expression factored through the
equivalence classes $\left[  D^{1}\left(  g^{N}\left(  \bm{0}\right)  \right)
\right]  $
\begin{align}
\mathcal{E}\left(  \bm{n},\bm{\varphi},\bm{\theta}\right)   &  =\frac
{\mathcal{H}\left(  \bm{n},\bm{\psi},\bm{\theta}\right)  }{S_{N}\left(
\bm{n}\right)  }=\frac{1}{S_{N}\left(  \bm{n}\right)  }\left\{  \sum_{1\leq
j\leq s}n_{j}S_{N-1}\left(  \jmath\right)  \left\{  \overset{}{\left\langle
\varphi_{j}\left\vert h_{1}\varphi_{j}\right.  \right\rangle +\left\langle
\varphi_{j+s}\left\vert h_{1}\varphi_{j+s}\right.  \right\rangle }\right.
\right. \nonumber\\
&  \left.  +\overset{}{\,\left\langle \varphi_{j}\varphi_{j+s}|h_{2}%
\varphi_{j}\varphi_{j+s}\right\rangle }\right\}  +\sum_{1\leq j,k\leq s}%
n_{j}^{\frac{1}{2}}n_{k}^{\frac{1}{2}}e^{i\left(  \theta_{j}-\theta
_{k}\right)  }S_{N-1}\left(  \jmath k\right)  \left\langle \varphi_{j}%
\varphi_{j+s}|h_{2}\varphi_{k}\varphi_{k+s}\right\rangle \nonumber\\
&  \qquad\qquad\qquad\qquad\qquad\left.  +\sum_{1\leq j<k\leq s}n_{i}%
n_{j}S_{N-2}\left(  \jmath k\right)  \left\langle \varphi_{j}\varphi
_{k}|||\varphi_{j}\varphi_{k}\right\rangle \right\}  , \label{thetaen}%
\end{align}
where $\left\vert \bm{\varphi}\right\rangle $ is a representative set of
NGSO's for the class $\left[  D^{1}\left(  g^{N}\left(  \bm{0}\right)
\right)  \right]  $. As $\left(  \bm{n},\bm{\varphi,\theta}\right)  $ uniquely
defines a $2N$-electron AGP-FORDO\ the expression eq. $\left(  \ref{thetaen}%
\right)  $ can be used to define a density matrix energy functional
\begin{equation}
\mathcal{F}\left(  D^{1}\right)  =\min_{\bm{\theta\in}\Theta}\mathcal{E}%
\left(  \bm{n},\bm{\varphi},\bm{\theta}\right)
\end{equation}
over the manifold of AGP\ states, where
\begin{equation}
\Theta=\left\{  \left.  \theta_{j}\right\vert \theta_{1}=0,\,0<\theta_{j}%
\leq2\pi,2\leq j\leq s\right\}  .
\end{equation}
Noting that each class is determined by $\left(  \bm{n},\bm{\varphi
}\right)  $ this functional can then be minimized to give
\begin{equation}
\mathcal{F}\left(  D_{\min}^{1}\right)  =\min_{D^{1}\in\mathcal{S}_{1}%
}\mathcal{F}\left(  D^{1}\right)  =\min_{\bm{n\in}\mathcal{D}_{n},\left\vert
\bm{\varphi}\right\rangle \in\mathcal{C}}\mathcal{F}\left(  D^{1}\left(
\bm{n},\bm{\varphi}\right)  \right)  =\min_{\left\vert g\right\rangle
\in\mathcal{H}^{2}}\mathcal{E}\left(  \left\vert g^{N}\right\rangle \right)  ,
\label{MinDMF}%
\end{equation}
where the domains $\mathcal{D}_{n}$, $\mathcal{C}$ and $\mathcal{S}_{1}$ are
given as%
\begin{align}
\mathcal{D}_{n}  &  =\left\{  \left.  n_{j}\right\vert ;\sum_{1\leq j\leq
s}n_{j}=1,n_{j}\geq0,\quad1\leq j\leq s\right\}  ,\\
\mathcal{C}  &  =\left\{  \left.  \left\vert \bm{\varphi}\right\rangle
\right\vert ;\left\vert \bm{\varphi}\right\rangle \neq e^{i\bm{\theta }}%
\left\vert \bm{\varphi}^{\prime}\right\rangle ,\text{ when }\left\vert
\bm{\varphi}^{\prime}\right\rangle \in\mathcal{C}\right\}  ,
\end{align}
where
\begin{equation}
e^{i\bm{\theta}}\left\vert \bm{\varphi}^{\prime}\right\rangle =\left(
e^{i\theta_{1}}\left\vert \varphi_{1}^{\prime}\right\rangle ,\cdots
,e^{i\theta_{2s}}\left\vert \varphi_{2s}^{\prime}\right\rangle \right)
\end{equation}
and
\begin{equation}
\mathcal{S}_{1}=\left\{  \left.  D^{1}\right\vert ;Tr\left\{  D^{1}\right\}
=2N,D^{1}\geq0,D^{1}\text{ at least doubly degenerate}\right\}  .
\end{equation}
One can see that the solution to eq. $\left(  \ref{MinDMF}\right)  $
characterizes the $2N$-electron AGP state that minimizes the total energy.

\subsection{Energy Functional in Terms of Group Parameters\label{Group_Param}}

The manifold of $2N$-electron AGP\ states can be parameterized in terms of a
nondegenerate ($c_{j}\neq c_{k}$ for $j\neq k$ ) reference geminal state
$\left\vert g_{0}\right\rangle $ and the subgroups, $U_{orb},$ $U_{phase}$ and
$G_{num}$ of the general linear group $GL\left(  \mathcal{H}^{1}\right)  $ of
one particle Hilbert space $\mathcal{H}^{1}$. The orbital variation
\emph{special} unitary subgroup $U_{orb}$ is defined as
\begin{equation}
U_{orb}=\left\{  \left.  u\left(  \bm{\alpha}\right)  =\exp\left\{
i\sum_{1\leq j\neq k\pm s\leq2s}\alpha_{jk}a_{k}^{\dagger}a_{j}\right\}
\right\vert \bm{\alpha}=\bm{\alpha}^{\dagger},\operatorname{Tr}%
\bm{\alpha}=0\right\}  , \label{orb}%
\end{equation}
where the second quantization operators $\left\{  a_{j}^{\dagger}\right\}  $
are defined by their action on the vacuum state $\left\vert \phi\right\rangle
$ by
\begin{equation}
a_{j}^{\dagger}\left\vert \phi\right\rangle =\left\vert \varphi_{j}%
\right\rangle ;\quad1\leq j\leq2s.
\end{equation}
The phase variation \emph{unitary} subgroup $U_{phase}$ is given by%
\begin{equation}
U_{phase}=\left\{  \left.  p\left(  \bm{\theta}\right)  =\exp\left\{
i\sum_{1\leq j\leq s}\theta_{j}\Omega_{j}\right\}  \right\vert 0\leq\theta
_{j}<2\pi\right\}  , \label{phase}%
\end{equation}
where
\begin{equation}
\Omega_{j}=\left(  a_{j}^{\dagger}a_{j}+a_{j+s}^{\dagger}a_{j+s}\right)
;\quad1\leq j\leq s
\end{equation}
and occupation number variation nonunitary subgroup $G_{num}$ by
\begin{equation}
G_{num}=\left\{  \left.  k\left(  \bm{\eta}\right)  =\exp\left\{  \sum_{1\leq
j\leq s}\eta_{j}\Omega_{j}\right\}  \right\vert \eta_{j}\in\mathbb{R}\right\}
. \label{occup}%
\end{equation}
The manifold of $2N$-electron AGP\ states $\mathcal{M}$ can then be expressed
as
\begin{align}
\mathcal{M}  &  \mathcal{=}\left\{  \left.  \overset{}{k\left(
\bm{\eta}\right)  p\left(  \bm{\theta}\right)  u\left(  \bm{\alpha}\right)
}\left\vert g_{0}^{N}\right\rangle \right\vert \quad0\leq\theta_{j}%
<2\pi;\,\eta_{j}\in\mathbb{R};\right. \nonumber\\
&  \qquad\qquad\qquad\qquad\qquad\left.  \bm{\alpha}=\bm{\alpha }^{\dagger
},\operatorname{Tr}\bm{\alpha}=0,1\leq j\neq k\pm s\leq2s\in\mathbb{C}%
\right\}
\end{align}
Explicitly the action on the reference geminal is given by
\begin{equation}
\left\vert g\left(  \bm{\eta},\bm{\alpha},\bm{\theta}\right)  \right\rangle
=k\left(  \bm{\eta}\right)  p\left(  \bm{\theta}\right)  u\left(
\bm{\alpha}\right)  \left\vert g_{0}\right\rangle =\sum_{1\leq l\leq s}%
e^{\eta_{l}}e^{i\theta_{l}}c_{0l}\left\vert \varphi_{l}\left(
\bm{\alpha}\right)  \varphi_{l+s}\left(  \bm{\alpha}\right)  \right\rangle ,
\end{equation}
where
\begin{equation}
\left\vert \varphi_{l}\left(  \bm{\alpha}\right)  \right\rangle =\exp\left\{
i\sum_{1\leq j\neq k\pm s\leq2s}\alpha_{jk}a_{j}^{\dagger}a_{k}\right\}
\left\vert \varphi_{0l}\right\rangle .
\end{equation}
It should be noted that the special unitary group, $U_{inv}$,
\begin{equation}
U_{inv}=\left\{  \left.  \exp i\left(  \sum_{1\leq j\leq s}\kappa_{jj}\left(
a_{j}^{\dagger}a_{j}-a_{j+s}^{\dagger}a_{j+s}\right)  +\kappa_{jj+s}%
a_{j}^{\dagger}a_{j+s}+\kappa_{jj+s}^{\ast}a_{j+s}^{\dagger}a_{j}\right)
\right\vert \kappa_{jj}\in\mathbb{R},\kappa_{jj+s}\in\mathbb{C}\right\}
\label{iso}%
\end{equation}
leaves the geminal $\left\vert g_{0}\right\rangle $ and AGP $\left\vert
g_{0}^{N}\right\rangle $ invariant as
\begin{equation}
\left(  a_{j}^{\dagger}a_{j}-a_{j+s}^{\dagger}a_{j+s}\right)  \left\vert
g_{0}\right\rangle =a_{j}^{\dagger}a_{j+s}\left\vert g_{0}\right\rangle
=a_{j+s}^{\dagger}a_{j}\left\vert g_{0}\right\rangle =0;\,1\leq j\leq s.
\label{iso1}%
\end{equation}


In terms of all these group parameters the energy functional can be expressed
as
\begin{equation}
\mathcal{E}\left(  \bm{\eta},\bm{\alpha},\bm{\theta}\right)  =\frac
{\mathcal{H}\left(  \bm{\eta},\bm{\alpha},\bm{\theta }\right)  }{S_{N}\left(
\bm{\eta}\right)  }=\frac{\left\langle \left.  g\left(
\bm{\eta},\bm{\alpha},\bm{\theta}\right)  ^{N}\right\vert Hg\left(
\bm{\eta},\bm{\alpha},\bm{\theta}\right)  ^{N}\right\rangle }{S_{N}\left(
\bm{\eta}\right)  }, \label{gp_funct}%
\end{equation}
where
\begin{equation}
\mathcal{H}\left(  \bm{\eta},\bm{\lambda},\bm{\theta}\right)  =\left\langle
\left.  k\left(  \bm{\eta}\right)  p\left(  \bm{\theta }\right)  u\left(
\bm{\alpha}\right)  g_{0}^{N}\right\vert Hk\left(  \bm{\eta}\right)  p\left(
\bm{\theta}\right)  u\left(  \bm{\alpha }\right)  g_{0}^{N}\right\rangle
\end{equation}
and
\begin{equation}
S_{N}\left(  \bm{\eta}\right)  =\left\langle \left.  k\left(
\bm{\eta}\right)  g_{0}^{N}\right\vert k\left(  \bm{\eta}\right)  g_{0}%
^{N}\right\rangle =\sum_{1\leq j_{1}<\cdots<j_{N}\leq s}e^{\eta_{j_{1}}%
+\cdots+\eta_{j_{N}}}n_{j_{1}}\cdots n_{j_{N}}.
\end{equation}
The parameters $\left(  \bm{\eta,\alpha}\right)  $ uniquely characterize the
FORDO $D^{1}\left(  g^{N}\left(  \bm{\eta},\bm{\alpha}\right)  \right)  $
i.e.
\begin{equation}
D^{1}\left(  \bm{\eta},\bm{\alpha}\right)  =D^{1}\left(  \bm{n}\left(
\bm{\eta}\right)  ,\bm{\varphi}\left(  \bm{\alpha}\right)  \right)
\end{equation}
so that a DMF can be unambiguously expressed as
\begin{equation}
\mathcal{F}\left(  \left(  \bm{\eta},\bm{\alpha}\right)  \right)
=\mathcal{F}\left(  D^{1}\left(  \bm{\eta},\bm{\alpha}\right)  \right)
=\min_{\bm{\theta\in}\Theta}\mathcal{E}\left(
\bm{\eta },\bm{\alpha},\bm{\theta}\right)
\end{equation}
and $\mathcal{F}\left(  \left(  \bm{\eta},\bm{\alpha}\right)  \right)  $ can
be minimized in an unconstrained fashion to give the optimized AGP energy
\begin{equation}
\mathcal{F}\left(  D_{\min}^{1}\right)  =\mathcal{F}\left(  \left(
\bm{\eta}_{\min},\bm{\alpha}_{\min}\right)  \right)  =\min
_{\bm{\eta},\bm{\alpha}}\mathcal{F}\left(  \left(
\bm{\eta },\bm{\alpha}\right)  \right)  .
\end{equation}


\section{Necessary Conditions for the Minimization of the Energy
\label{Evariation}}

We first consider the energy functional in terms of $\left(
\bm{n,\psi }\right)  ,$ which produces a constrained optimization problem,
then in terms of the unconstrained variables $\left(
\bm{\eta},\bm{\alpha },\bm{\theta}\right)  $. The first approach introduces a
generalized Fock operator, while the second approach shows that the Fock
operator $F\left(  \bm{n,\psi}\right)  $ must be self adjoint when based on an
optimized state. The details of the derivations are contained in Appendix
\ref{Deriv_Fock}, \ref{VarGroupPar} and \ref{StatCon}, where expressions based
on are general $2N$-electron state are obtained.

\subsection{Constrained Variational Problem\label{CVP}}

The solution to the optimization problem
\begin{equation}
\min_{\left(  \bm{n,\psi}\right)  }\mathcal{E}\left(  \bm{n,\psi }\right)
\label{min}%
\end{equation}
subject to the orthonormality constraints%
\begin{equation}
\left\langle \left.  \psi_{j}\right\vert \psi_{k}\right\rangle =\delta
_{jk};\quad1\leq j,k\leq2s \label{orthonorm}%
\end{equation}
and the constraints of positivity and normalization on the geminal occupation
numbers
\begin{subequations}
\begin{align}
n_{j}  &  \geq0;\quad1\leq j\leq s,\label{pos1}\\
\sum_{1\leq j\leq s}n_{j}  &  =1 \label{gemnorm}%
\end{align}
characterize an optimal $2N$-electron AGP\ state.

The occupation number constraints eqs. $\left(  \ref{pos1}\right)  $, $\left(
\ref{gemnorm}\right)  $ and the orthonormality constraints can be enforced by
the method of Lagrange multipliers and the introduction of a Lagrangian
\end{subequations}
\begin{equation}
\mathcal{L}_{T}\left(  \bm{n,\psi,\varepsilon,\xi,}\eta\right)  =\mathcal{E}%
\left(  \bm{n,\psi}\right)  -\sum_{1\leq j,k\leq2s}\varepsilon_{kj}\left(
\left\langle \left.  \psi_{j}\right\vert \psi_{k}\right\rangle -\delta
_{jk}\right)  -\eta\left(  \sum_{1\leq j\leq s}n_{j}-1\right)  +\sum_{1\leq
j\leq s}\xi_{j}n_{j}. \label{Lar}%
\end{equation}
The occupation number constraints are most effectively invoked and analyzed by
expressing the geminal occupation numbers as unconstrained functions of group
parameters introduced in section (\ref{Group_Param}). Thus we will not
consider the occupation constraint part of the Lagrangian further. The orbital
constraint part
\begin{equation}
\mathcal{L}\left(  \bm{n,\psi}\right)  =\mathcal{E}\left(
\bm{n,\psi }\right)  -\sum_{1\leq j,k\leq2s}\varepsilon_{kj}\left(
\left\langle \left.  \psi_{j}\right\vert \psi_{k}\right\rangle -\delta
_{jk}\right)  . \label{larE}%
\end{equation}
of the Lagrangian eq. $\left(  \ref{Lar}\right)  $ reveals some properties of
a generalized Fock operator $F\left(  \bm{n,\psi}\right)  $. The stationarity
conditions
\begin{equation}
\frac{\delta\mathcal{L}\left(  \bm{n,\psi}\right)  }{\delta\psi_{j}^{\ast}%
}=0;\quad1\leq j\leq2s,
\end{equation}
which involve the derivatives
\begin{equation}
\frac{\delta\mathcal{E}\left(  \bm{n,\psi}\right)  }{\delta\psi_{j}^{\ast}%
}=\frac{1}{S_{N}\left(  \bm{n}\right)  }\left\{  \frac{\delta\mathcal{H}%
\left(  \bm{n,\psi}\right)  }{\delta\psi_{j}^{\ast}}-\mathcal{E}\left(
\bm{n,\psi}\right)  \frac{\delta S_{N}\left(  \bm{n}\right)  }{\delta\psi
_{j}^{\ast}}\right\}  =\frac{1}{S_{N}\left(  \bm{n}\right)  }\frac
{\delta\mathcal{H}\left(  \bm{n,\psi}\right)  }{\delta\psi_{j}^{\ast}}%
\end{equation}
produce
\begin{equation}
\frac{\delta\mathcal{E}\left(  \bm{n,\psi}\right)  }{\delta\psi_{j}^{\ast}%
}=\sum_{1\leq j,k\leq2s}\left\vert \psi_{k}\right\rangle \varepsilon_{kj}.
\label{statcon}%
\end{equation}
We show in section \ref{Fock} that
\begin{equation}
\frac{\delta\mathcal{E}\left(  \bm{n,\psi}\right)  }{\delta\psi_{j}^{\ast}%
}=F\left(  \bm{n,\psi}\right)  \left\vert \psi_{j}\right\rangle
\end{equation}
so eq. $\left(  \ref{statcon}\right)  $ becomes
\begin{equation}
F\left(  \bm{n,\psi}\right)  \left\vert \psi_{j}\right\rangle =\sum_{1\leq
j,k\leq2s}\left\vert \psi_{k}\right\rangle \varepsilon_{kj}. \label{Fockeq}%
\end{equation}
As (see section \ref{Fock})%
\begin{equation}
\frac{\delta\mathcal{E}\left(  \bm{n,\psi}\right)  }{\delta\psi_{j}}=\left(
\frac{\delta\mathcal{E}\left(  \bm{n,\psi}\right)  }{\delta\psi_{j}^{\ast}%
}\right)  ^{\ast}%
\end{equation}
it is sufficient to consider derivatives wrt $\left\{  \psi_{j}^{\ast
}\right\}  $.

In \emph{contrast} to the IPS\ case, but in \emph{common} with Multi
Configurational Self Consistent Field (MCSCF)\ procedures, eq. $\left(
\ref{Fockeq}\right)  $ is not an eigenvalue equation. However by considering
the stationarity conditions from the unitary group parameterization one can
observe that at a stationary point of the orbital variation $F\left(
\bm{n,\psi}\right)  $ becomes self adjoint and hence
\begin{equation}
\varepsilon_{jk}=\varepsilon_{kj}^{\ast}%
\end{equation}
and the matrix $\bm{\varepsilon}$ of Lagrange parameters is hermitian.

\subsection{Unconstrained Variational Problem\label{UCVP}}

In this case we consider the problem
\begin{equation}
\min_{\left(  \bm{\eta,\alpha,\theta}\right)  }\mathcal{E}\left(
\bm{\eta,\alpha,\theta}\right)  ,
\end{equation}
where the energy functional is given in eq. $\left(  \ref{gp_funct}\right)  $
in terms of the groups eqs. $\left(  \ref{orb}\right)  ,\left(  \ref{phase}%
\right)  $ and $\left(  \ref{occup}\right)  $ parametrized by $\left(
\bm{\eta,\alpha,\theta}\right)  $. Setting
\begin{equation}
g_{0}\equiv g\left(  \bm{0,0,0}\right)
\end{equation}
this leads to (see Appendix \ref{variational}) the orbital variation
derivatives
\begin{equation}
\partial_{\bm{\lambda}_{jk}}^{1}\mathcal{E}\left(  g_{0}^{N}\right)  =\frac
{i}{S_{N}\left(  g_{0}\right)  }\left\langle g_{0}^{N}\left\vert \left[
H,a_{j}^{\dagger}a_{k}\right]  g_{0}^{N}\right.  \right\rangle ;\quad1\leq
j\neq k\pm s\leq2s,
\end{equation}
the phase variation derivatives
\begin{equation}
\partial_{\bm{\theta}_{j}}^{1}\mathcal{E}\left(  g_{0}^{N}\right)  =\frac
{i}{S_{N}\left(  g_{0}\right)  }\left\langle g_{0}^{N}\left\vert \left[
H,\left(  a_{j}^{\dagger}a_{j}+a_{j+s}^{\dagger}a_{j+s}\right)  \right]
g_{0}^{N}\right.  \right\rangle ;\quad1\leq j\leq s,
\end{equation}
and the occupation number variations derivatives
\begin{align}
\partial_{\bm{\eta}_{j}}^{1}\mathcal{E}\left(  g_{0}^{N}\right)   &  =\frac
{1}{S_{N}\left(  g_{0}\right)  }\left\{  \left\langle g_{0}^{N}\left\vert
\left[  H,\left(  a_{j}^{\dagger}a_{j}+a_{j+s}^{\dagger}a_{j+s}\right)
\right]  _{+}^{N}g_{0}\right.  \right\rangle \right. \nonumber\\
&  -\left.  2\mathcal{E}\left(  g_{0}^{N}\right)  \left\langle g_{0}%
^{N}\left\vert \left(  a_{j}^{\dagger}a_{j}+a_{j+s}^{\dagger}a_{j+s}\right)
g_{0}^{N}\right.  \right\rangle \right\}  ;\quad1\leq j\leq s.
\end{align}
All the derivatives can be evaluated at $\left(  \bm{0},\bm{0},\bm{0}\right)
$ by updating the reference geminal $\left\vert g_{0}\right\rangle
\equiv\left\vert g_{old}\right\rangle $ by a transformation
\begin{equation}
\left\vert g_{new}\right\rangle =k\left(  \bm{\eta}\right)  p\left(
\bm{\theta}\right)  u\left(  \bm{\alpha}\right)  \left\vert g_{old}%
\right\rangle
\end{equation}
to produce a new reference $\left\vert g_{0}\right\rangle \equiv\left\vert
g_{new}\right\rangle .$ The stationarity conditions (see Appendix
\ref{variational}) for orbital variation gives
\begin{equation}
\left\langle g_{0}^{N}\left\vert \left[  H,a_{j}^{\dagger}a_{k}\right]
g_{0}^{N}\right.  \right\rangle =0;\quad1\leq j\neq k\pm s\leq2s,
\label{ostat}%
\end{equation}
while phase stationarity gives%
\begin{equation}
\left\langle g_{0}^{N}\left\vert \left[  H,\left(  a_{j}^{\dagger}%
a_{j}+a_{j+s}^{\dagger}a_{j+s}\right)  \right]  g_{0}^{N}\right.
\right\rangle =0;\quad1\leq j\leq s.
\end{equation}
Taking into account the invariance group relationship eq. $\left(
\ref{iso1}\right)  $ this gives
\begin{equation}
\left\langle g_{0}^{N}\left\vert \left[  H,a_{j}^{\dagger}a_{j}\right]
g_{0}^{N}\right.  \right\rangle =0;\quad1\leq j\leq2s. \label{phasestat}%
\end{equation}
Eqs. $\left(  \ref{ostat}\right)  $ and $\left(  \ref{phasestat}\right)  $
together with the invariance group relationships eq. $\left(  \ref{iso1}%
\right)  $ give
\begin{equation}
\left\langle g_{0}^{N}\left\vert \left[  H,a_{j}^{\dagger}a_{k}\right]
g_{0}^{N}\right.  \right\rangle =0;\quad1\leq j,k\leq2s. \label{orbstat}%
\end{equation}


The preceding commutator relationships eq. $\left(  \ref{orbstat}\right)  $
are the same as are produced by stationary states, $\left\vert \Psi
\right\rangle ,$ in general MCSCF procedures, i.e.
\begin{equation}
\left\langle \Psi\left\vert \left[  H,a_{j}^{\dagger}a_{k}\right]
\Psi\right.  \right\rangle =0;\quad1\leq j,k\leq2s.
\end{equation}
Optimized AGP states are also stationary to occupation number variations
giving
\begin{equation}
\left\langle g_{0}^{N}\left\vert \left[  H,\left(  a_{j}^{\dagger}%
a_{j}+a_{j+s}^{\dagger}a_{j+s}\right)  \right]  _{+}g_{0}^{N}\right.
\right\rangle -2\mathcal{E}\left(  g_{0}^{N}\right)  \left\langle g_{0}%
^{N}\left\vert \left(  a_{j}^{\dagger}a_{j}+a_{j+s}^{\dagger}a_{j+s}\right)
g_{0}^{N}\right.  \right\rangle =0;\quad1\leq j\leq s,
\end{equation}
which using the special case of eq. $\left(  \ref{orbstat}\right)  $
\begin{equation}
\left\langle g_{0}^{N}\left\vert Ha_{j}^{\dagger}a_{j}g_{0}^{N}\right.
\right\rangle =\left\langle g_{0}^{N}\left\vert a_{j}^{\dagger}a_{j}Hg_{0}%
^{N}\right.  \right\rangle ;\quad1\leq j\leq2s
\end{equation}
leads to the relationship
\begin{equation}
\left\langle g_{0}^{N}\left\vert Ha_{j}^{\dagger}a_{j}g_{0}^{N}\right.
\right\rangle =\mathcal{E}\left(  g_{0}^{N}\right)  \left\langle g_{0}%
^{N}\left\vert a_{j}^{\dagger}a_{j}g_{0}^{N}\right.  \right\rangle ;\quad1\leq
j\leq2s.
\end{equation}
This can be recast in the form
\begin{equation}
\frac{1}{S_{N}\left(  g_{0}\right)  }\left\langle g_{0}^{N}\left\vert
Ha_{j}^{\dagger}a_{j}g_{0}^{N}\right.  \right\rangle =\mathcal{E}\left(
g_{0}^{N}\right)  N_{j};\quad1\leq j\leq2s. \label{occnumstat}%
\end{equation}
In the context of general optimized multi configurational states the
generators $\left\{  a_{j}^{\dagger}a_{j}\right\}  $ produce CI\ coefficients
variations, hence even for fully optimized Multi Configurational (MC) states
one has
\begin{equation}
\frac{1}{\left\langle \left.  \Psi\right\vert \Psi\right\rangle }\left\langle
\Psi\left\vert Ha_{j}^{\dagger}a_{j}\Psi\right.  \right\rangle =\mathcal{E}%
\left(  \Psi\right)  N_{j};\quad1\leq j\leq2s \label{MCocc}%
\end{equation}
as
\begin{equation}
\left\langle \Psi\left\vert \left[  H,a_{j}^{\dagger}a_{j}\right]  _{+}%
\Psi\right.  \right\rangle -\mathcal{E}\left(  \Psi\right)  \left\langle
\Psi\left\vert a_{j}^{\dagger}a_{j}\Psi\right.  \right\rangle =0;\quad1\leq
j\leq2s.
\end{equation}


\subsection{Generalized Brillouin Condition for Optimized AGP
States.\label{Brillouin}}

If an AGP\ state has nondegenerate canonical coefficients i.e. $c_{j}\neq
c_{k}$ for $1\leq j\neq k\leq s$ one can \cite{AGPCoherent}\cite{Weiner83a}%
\cite{Kurtz83} always construct operators (see Appendix \ref{q_op}) $\left\{
\left.  q_{jk}^{\dagger}\right\vert \quad1\leq j\neq k\pm s\leq2s\right\}  $
that have the property
\begin{equation}
q_{jk}^{\dagger}\left\vert g_{0}^{N}\right\rangle \neq0,\quad q_{jk}\left\vert
g_{0}^{N}\right\rangle =0;\quad1\leq j\neq k\pm s\leq2s,
\end{equation}
such that the orbital transformation group $U_{orb}$ can be expressed as
\begin{equation}
U_{orb}=\left\{  u\left(  \bm{\beta}\right)  =\exp\left\{  i\sum
_{1\leq\left\vert j\right\vert <\left\vert k\right\vert \leq s}\left\{
\beta_{jk}q_{jk}^{\dagger}+\beta_{jk}^{\ast}q_{jk}\right\}  \right\}
\right\}  ,
\end{equation}
where
\begin{equation}
\left\vert j\right\vert =\left\{
\begin{array}
[c]{c}%
j;\qquad1\leq j\leq s\\
j-s;\qquad s+1\leq j\leq2s
\end{array}
\right.  .
\end{equation}
This parameterization of $U_{orb}$ leads to the stationarity conditions
\begin{equation}
\left\langle \left.  g_{0}^{N}\right\vert \left[  H,q_{jk}^{\dagger}\right]
g_{0}^{N}\right\rangle =\left\langle \left.  g_{0}^{N}\right\vert \left[
H,q_{jk}\right]  g_{0}^{N}\right\rangle =0;\quad1\leq\left\vert j\right\vert
<\left\vert k\right\vert \leq s,
\end{equation}
that give the "Brillouin Conditions"
\begin{equation}
\left\langle \left.  g_{0}^{N}\right\vert Hq_{jk}^{\dagger}g_{0}%
^{N}\right\rangle =\left\langle \left.  g_{0}^{N}\right\vert q_{jk}Hg_{0}%
^{N}\right\rangle =0;\quad1\leq\left\vert j\right\vert <\left\vert
k\right\vert \leq s.\label{BC}%
\end{equation}
In Appendix \ref{q_op} we describe how the $\left\{  q_{jk}^{\dagger}%
,q_{jk}\right\}  $ basis can be transformed into the $\left\{  a_{j}^{\dagger
}a_{k}\right\}  $ basis of the Lie algebra of $U_{orb}$ by a block diagonal
nonsingular transformation, $\mathbb{T},$ such that the blocks $\mathbb{T}%
\left(  jk\right)  $ have the effect
\begin{align}
&  \left(  q_{jk}^{\dagger},q_{j+sk}^{\dagger},q_{jk+s}^{\dagger}%
,q_{j+sk+s}^{\dagger},q_{jk},q_{j+sk},q_{jk+s},q_{j+sk+s}\right)
\mathbb{T}\left(  jk\right)  \nonumber\\
&  \qquad\qquad\qquad\qquad\qquad\left.  =\right.  \left(  a_{j}^{\dagger
}a_{k},a_{j}^{\dagger}a_{k+s},a_{j+s}^{\dagger}a_{k},a_{j+s}^{\dagger}%
a_{k+s},a_{k}^{\dagger}a_{j},a_{k}^{\dagger}a_{j+s},a_{k+s}^{\dagger}%
a_{j},a_{k+s}^{\dagger}a_{j+s}\right)  ,
\end{align}
where
\begin{equation}
\mathbb{T}\left(  jk\right)  =\left(
\begin{array}
[c]{cc}%
\mathbb{T}_{1}\left(  jk\right)   & \mathbb{T}_{2}\left(  jk\right)  \\
\mathbb{T}_{2}\left(  jk\right)  ^{\ast} & \mathbb{T}_{1}\left(  jk\right)
^{\ast}%
\end{array}
\right)  ,
\end{equation}
giving
\begin{equation}
a_{l}^{\dagger}a_{m}=\sum_{p,q=0,s}\left\{  q_{j+pk+q}^{\dagger}T_{1}\left(
jk\right)  _{pqlm}+q_{j+pk+q}T_{2}\left(  jk\right)  _{pqlm}^{\ast}\right\}
.\label{Trans}%
\end{equation}
The Brillouin conditions eq. $\left(  \ref{BC}\right)  $ combined with eq.
$\left(  \ref{Trans}\right)  $ results in the generalized Brillouin
conditions
\begin{equation}
\left\langle g_{0}^{N}\left\vert Ha_{l}^{\dagger}a_{m}\right.  g_{0}%
^{N}\right\rangle =0;\quad1\leq l\neq m\pm s\leq2s,\label{GBC}%
\end{equation}
as $a_{l}^{\dagger}a_{m}$ can be expanded by eq. $\left(  \ref{Trans}\right)
$ to give%
\begin{equation}
\left\langle g_{0}^{N}\left\vert H\sum_{p,q=0,s}q_{j+pk+q}^{\dagger}%
T_{1}\left(  jk\right)  _{pqlm}-\sum_{p,q=0,s}q_{j+pk+q}T_{2}\left(
jk\right)  _{pqlm}^{\ast}Hg_{0}^{N}\right.  \right\rangle =0;\quad1\leq l\neq
m\pm s\leq2s.
\end{equation}
The remaining off diagonal particle-hole operators $\left\{  \left.
a_{l}^{\dagger}a_{m}\right\vert 1\leq l=m\pm s\leq2s\right\}  $ automatically
satisfy a generalized Brillouin condition
\begin{equation}
\left\langle g_{0}^{N}\left\vert Ha_{l}^{\dagger}a_{m}g_{0}^{N}\right.
\right\rangle =0;\quad1\leq l=m\pm s\leq2s\label{ph}%
\end{equation}
by the invariance group relationships eq. $\left(  \ref{iso1}\right)  .$

One should note that the Brillouin conditions eqs. $\left(  \ref{GBC}\right)
$ and $\left(  \ref{ph}\right)  $ hold between GSO's no matter whether they
are occupied (in CI terminology referred to as active or inactive) i.e.
\begin{equation}
\left\langle g_{0}^{N}\left\vert a_{k}^{\dagger}a_{k}g_{0}^{N}\right.
\right\rangle \neq0
\end{equation}
or unoccupied (in CI terminology referred to as virtual) i.e.
\begin{equation}
\left\langle g_{0}^{N}\left\vert a_{k}^{\dagger}a_{k}g_{0}^{N}\right.
\right\rangle =0.
\end{equation}
This should be contrasted to the MCSCF Brillouin condition where
\begin{equation}
\left\langle \Psi\left\vert Ha_{l}^{\dagger}a_{m}\right.  \Psi\right\rangle
\left\{
\begin{array}
[c]{c}%
\quad=0\text{ for }m\text{ being an unoccupied GSO and }l\text{ being an
unoccupied GSO \quad\quad}\\
=0\text{ for }m\text{ being an unoccupied GSO and }l\text{ being an occupied
GSO \quad\quad}\\
\neq0\text{ for }m\text{ being an occupied GSO and }l\text{ being an occupied
GSO\quad\quad\quad\ }%
\end{array}
\right.  . \label{MBC}%
\end{equation}


If there are degeneracies of the geminal canonical coefficients i.e.
$c_{j}=c_{k}$ for some $1\leq j\neq k\leq s,$ then the $\left\{
q_{jk}^{\dagger}\right\}  $ and $\left\{  q_{jk}\right\}  $ associated with
those coefficients can be replaced \cite{Goscinski82a} with sets of self
adjoint operators $\left\{  u_{jk}\right\}  $ and $\left\{  v_{ij}\right\}  $
respectively, where
\begin{equation}
u_{jk}\left\vert g_{0}^{N}\right\rangle \neq0,\quad v_{jk}\left\vert g_{0}%
^{N}\right\rangle =0,
\end{equation}
and one can deduce that a generalized Brillouin condition of the form eq.
$\left(  \ref{GBC}\right)  $ \emph{does not} hold for $\left\{  a_{j}%
^{\dagger}a_{k}\right\}  $ associated with those indices and the situation
becomes similar to the MCSCF Brillouin condition eq. $\left(  \ref{MBC}%
\right)  $ i.e. the matrix elements $\left\langle \Psi\left\vert
Ha_{l}^{\dagger}a_{m}\right.  \Psi\right\rangle $ only vanish when $l$ and/or
$m$ are unoccupied.

\section{Generalized Fock Operator \label{Fock} for the Variation wrt
$\bm{\psi}$}

In this section we derive a generalization of the IPS\ Fock operator. This
allows us to obtain expressions relating observable properties of GSO's, i.e.
their ionization potentials, kinetic energy and potential energy, to the total
$2N$-electron state energy.

\subsection{Functional Derivatives\label{FD}}

The energy expression eq. $\left(  \ref{energy}\right)  $ can be
differentiated wrt $\bm{\psi}^{\ast},$ $\bm{\psi}$ to produce the functional
derivatives (Appendix \ref{variational})
\begin{equation}
\left\{  \frac{\delta\mathcal{E}\left(  \bm{n,\psi}\right)  }{\delta\psi
_{j}^{\ast}},\frac{\delta\mathcal{E}\left(  \bm{n,\psi}\right)  }{\delta
\psi_{j}};\quad1\leq j\leq s\right\}  ,
\end{equation}
which can be expressed as%
\begin{align}
\frac{\delta\mathcal{E}\left(  \bm{n,\psi}\right)  }{\delta\psi_{j}^{\ast}}
&  =\frac{1}{S_{N}\left(  \bm{n}\right)  }\frac{\delta\mathcal{H}\left(
\bm{n,\psi}\right)  }{\delta\psi_{j}^{\ast}}\\
\frac{\delta\mathcal{E}\left(  \bm{n,\psi}\right)  }{\delta\psi_{j}}  &
=\frac{1}{S_{N}\left(  \bm{n}\right)  }\frac{\delta\mathcal{H}\left(
\bm{n,\psi}\right)  }{\delta\psi_{j}}%
\end{align}
in terms of the derivatives of the unnormalized energy $\mathcal{H}\left(
\bm{n,\psi}\right)  ,$ as $S_{N}\left(  \bm{n}\right)  $ does not depend on
$\bm{\psi}$ when $\left\{  \left.  \left\vert \psi_{j}\right\rangle
\right\vert \quad1\leq j\leq2s\right\}  $ are orthonormal. It is easy to see
that
\begin{equation}
\frac{\delta\mathcal{E}\left(  \bm{n,\psi}\right)  }{\delta\psi_{j}^{\ast}%
}=\left(  \frac{\delta\mathcal{E}\left(  \bm{n,\psi}\right)  }{\delta\psi_{j}%
}\right)  ^{\ast}%
\end{equation}
so we only need consider $\left\{  \frac{\delta\mathcal{E}\left(
\bm{n,\psi}\right)  }{\delta\psi_{j}^{\ast}}\right\}  .$ These derivatives can
be expressed $\left(  \text{see Appendix \ref{variational}}\right)  $ in terms
of a generalized Fock operator $F\left(  \bm{n,\psi}\right)  $ as
\begin{equation}
\frac{\delta\mathcal{E}\left(  \bm{n,\psi}\right)  }{\delta\psi_{j}^{\ast}%
}=F\left(  \bm{n,\psi}\right)  \left\vert \psi_{j}\right\rangle ,\quad1\leq
j\leq2s, \label{H_deriv}%
\end{equation}
that depends on the occupation numbers $\bm{n,}$ and the GSO's $\bm{\psi}$ and
is given by
\begin{equation}
F\left(  \bm{n,\psi}\right)  =\sum_{1\leq m,n\leq2s}\left\{  \left(
\mathbb{H}_{1}\left(  \bm{\psi}\right)  \mathbb{D}^{1}\left(  \bm{n}\right)
\right)  _{mn}+\frac{1}{2}\left(  L_{2}^{1}\left(  \mathbb{V}_{2}\left(
\bm{\psi}\right)  \mathbb{D}^{2}\left(  \bm{n}\right)  \right)  \right)
_{mn}\right\}  \left\vert \psi_{m}\right\rangle \left\langle \psi
_{n}\right\vert . \label{genfock}%
\end{equation}
The contraction operator $L_{2}^{1}$ is defined as
\begin{equation}
L_{2}^{1}\left(  \mathbb{A}\right)  _{jk}=\sum_{1\leq l\leq2s}A_{jlkl},
\end{equation}
$\mathbb{H}_{1}\left(  \bm{\psi}\right)  $ is the matrix of the one particle
Hamiltonian $h_{1}$, $\mathbb{V}_{2}\left(  \bm{\psi}\right)  $ is the matrix
of the two particle Hamiltonian $h_{2}$, $\mathbb{D}^{1}\left(  \bm{n}\right)
$ the FORDM and $\mathbb{D}^{2}\left(  \bm{n}\right)  $ the SORDM of the
$2N$-elecron AGP\ state (see in eq. $\left(  \ref{D2C}\right)  ),$ wrt the
basis $\left\{  \left\vert \psi_{j}\right\rangle ;1\leq j\leq2s\right\}  .$
One should recall that in L\"{o}wdin normalization and using an unordered
basis of two electron space
\begin{equation}
\mathbb{D}^{1}\left(  \bm{n}\right)  =\left(  2N-1\right)  ^{-1}L_{2}%
^{1}\left(  \mathbb{D}^{2}\left(  \bm{n}\right)  \right)  .
\end{equation}


The form of the expression eq. $\left(  \ref{genfock}\right)  $ is in fact
valid for all $N$-representable SORDM $\mathbb{D}^{2}$'s (not just AGP
SORDM's) and their associated FORDM, $\mathbb{D}^{1}$'s giving
\begin{equation}
F\left(  \mathbb{D}^{2}\bm{,\psi}\right)  =\sum_{1\leq m,n\leq2s}\left\{
\left(  \mathbb{H}_{1}\left(  \bm{\psi}\right)  \mathbb{D}^{1}\right)
_{mn}+\frac{1}{2}\left(  L_{2}^{1}\left(  \mathbb{V}_{2}\left(
\bm{\psi }\right)  \mathbb{D}^{2}\right)  \right)  _{mn}\right\}  \left\vert
\psi_{m}\right\rangle \left\langle \psi_{n}\right\vert .
\end{equation}
Thus they define a generalized Fock operator for \emph{any }MC state and
always produces the orbital functional derivatives eq. $\left(  \ref{H_deriv}%
\right)  .$ In general this generalized Fock operator is \emph{not }self
adjoint\emph{ }unless the state that it is based on fully optimizes the energy functional.

\subsection{Properties of a Stationary Point Fock Operator\label{Stat_Fock}}

We show in Appendix \ref{VarGroupPar} that the derivatives wrt the group
parameters controlling the GSO variations are
\begin{equation}
\frac{\partial\mathcal{E}\left(  \mathbb{D}^{2},\bm{0}\right)  }%
{\partial\alpha_{pq}}=i\left(  \left[  \mathbb{D}^{1},\mathbb{H}_{1}\left(
\bm{\varphi}\left(  \bm{0}\right)  \right)  \right]  +\frac{1}{2}L_{2}%
^{1}\left(  \left[  \mathbb{D}^{2},\mathbb{V}_{2}\left(  \bm{\varphi }\left(
\bm{0}\right)  \right)  \right]  \right)  \right)  _{qp}%
\end{equation}
and in Appendix \ref{Deriv_Fock} that the generalized Fock operator matrix
elements are
\begin{equation}
F\left(  \mathbb{D}^{2},\bm{\varphi}\right)  _{qp}=\operatorname{Tr}\left\{
\left(  \mathbb{H}_{1}\left(  \bm{\varphi}\right)  \mathbb{D}^{1}\right)
_{qp}+\frac{1}{2}\left(  L_{2}^{1}\left(  \mathbb{V}_{2}\left(
\bm{\varphi}\right)  \mathbb{D}^{2}\right)  \right)  _{qp}\right\}
\end{equation}
for \emph{any} state described by a SORDM $\mathbb{D}^{2}$. These two
equations give that
\begin{equation}
\frac{\partial\mathcal{E}\left(  \mathbb{D}^{2},\bm{0}\right)  }%
{\partial\alpha_{pq}}=i\left(  F\left(  \mathbb{D}^{2},\bm{\varphi}\left(
\bm{0}\right)  \right)  ^{\dagger}-F\left(  \mathbb{D}^{2}%
,\bm{\varphi }\left(  \bm{0}\right)  \right)  \right)  _{qp},
\end{equation}
which show if the state is stationary wrt orbital transformations then the
generalized Fock operator is self adjoint as
\begin{equation}
\frac{\partial\mathcal{E}\left(  \mathbb{D}^{2},\bm{0}\right)  }%
{\partial\alpha_{pq}}=0=i\left(  F\left(  \mathbb{D}^{2},\bm{\varphi }\left(
\bm{0}\right)  \right)  ^{\dagger}-F\left(  \mathbb{D}^{2},\bm{\varphi}\left(
\bm{0}\right)  \right)  \right)  _{qp};\quad1\leq p,q\leq2s.
\end{equation}


The generalized Fock operator associated with an optimized MC state has
further significant properties, in addition to being self adjoint, originating from

\begin{enumerate}
\item the stationarity wrt occupation number variations shown in eq. $\left(
\ref{MCocc}\right)  $ as
\begin{equation}
\frac{1}{\left\langle \left.  \Psi\right\vert \Psi\right\rangle }\left\langle
\left.  \Psi\right\vert Ha_{p}^{\dagger}a_{p}\Psi\right\rangle =\mathcal{E}%
\left(  \Psi\right)  N_{p};\quad1\leq p\leq2s,
\end{equation}
and

\item the expression for the generalized Fock operator matrix elements
\cite{YellowBook}
\begin{equation}
F\left(  \mathbb{D}^{2},\bm{\varphi}\right)  _{qp}=\frac{1}{\operatorname{Tr}%
\left\{  D^{2N}\right\}  }\operatorname{Tr}\left\{  a_{p}^{\dagger}\left[
a_{q},H\right]  D^{2N}\right\}
\end{equation}
obtained in Appendix \ref{Deriv_Fock}, which are valid for any state described
by the positive operator $D^{2N}$. For a pure state $\left\vert \Psi
\right\rangle $ this becomes
\begin{equation}
F\left(  \bm{n},\bm{\varphi}\right)  _{qp}=\frac{\left\langle \left.
\Psi\right\vert \left\{  a_{p}^{\dagger}a_{q}H-a_{p}^{\dagger}Ha_{q}\right\}
\Psi\right\rangle }{\left\langle \left.  \Psi\right\vert \Psi\right\rangle }.
\end{equation}

\end{enumerate}

Combining items $\left(  1\right)  $ and $\left(  2\right)  $ together shows
that the diagonal elements of the optimized MCSCF\ Fock operator are given by
\begin{equation}
F\left(  \mathbb{D}^{2},\bm{\varphi}\right)  _{pp}=\frac{\left\langle \left.
\Psi\right\vert \left\{  a_{p}^{\dagger}a_{q}H-a_{p}^{\dagger}Ha_{q}\right\}
\Psi\right\rangle }{\left\langle \left.  \Psi\right\vert \Psi\right\rangle
}=\mathcal{E}\left(  \mathbb{D}^{2},\bm{\varphi}\right)  \Omega_{p}%
-\frac{\left\langle \left.  a_{p}\Psi\right\vert Ha_{p}\Psi\right\rangle
}{\left\langle \left.  \Psi\right\vert \Psi\right\rangle },
\end{equation}
where the occupation, $\Omega_{p}$, of the GSO $\left\vert \varphi
_{p}\right\rangle $ is given by
\begin{equation}
\Omega_{p}=\frac{\left\langle \left.  \Psi\right\vert a_{p}^{\dagger}a_{p}%
\Psi\right\rangle }{\left\langle \left.  \Psi\right\vert \Psi\right\rangle }.
\end{equation}
If $\Omega_{p}$ $\neq0$ this can be developed to give
\begin{equation}
F\left(  \mathbb{D}^{2},\bm{\varphi}\right)  _{pp}=\Omega_{p}\left\{
\mathcal{E}\left(  \mathbb{D}^{2},\bm{\varphi}\right)  -\frac{\left\langle
\left.  a_{p}\Psi\right\vert Ha_{p}\Psi\right\rangle }{\left\langle \left.
a_{p}\Psi\right\vert a_{p}\Psi\right\rangle }\right\}  =\Omega_{p}\left\{
\mathcal{E}\left(  \mathbb{D}^{2},\bm{\varphi}\right)  -\mathcal{E}%
_{2N-1}\left(  \mathbb{D}^{2},\bm{\varphi,}\varphi_{p}\right)  \right\}  ,
\end{equation}
where $\mathcal{E}_{2N-1}\left(  \mathbb{D}^{2},\bm{\varphi,}\varphi
_{p}\right)  $ is the energy of the normalized $2N-1$ electron state produced
from the state $\frac{1}{\sqrt{\left\langle \left.  \Psi\right\vert
\Psi\right\rangle }}\Psi$ by removing an electron described by the
GSO\ $\left\vert \varphi_{p}\right\rangle $. If $\Omega_{p}$ $=0$ then
$F\left(  \mathbb{D}^{2},\bm{\varphi}\right)  _{pp}=0$ as the GSO $\left\vert
\varphi_{p}\right\rangle $ is unoccupied. In Appendix \ref{Fock&Energy} we
show that the total energy can be expressed in terms of the Fock operator
$F\left(  \mathbb{D}^{2},\bm{\varphi}\right)  \ $as
\begin{equation}
\mathcal{E}\left(  \mathbb{D}^{2},\bm{\varphi}\right)  =\frac{1}{2}\left\{
\overset{}{\operatorname{Tr}\left\{  F\left(  \mathbb{D}^{2}%
,\bm{\varphi}\right)  \right\}  +\mathcal{E}_{1}\left(  \mathbb{D}%
^{2},\bm{\varphi}\right)  }\right\}  ,
\end{equation}
where $\mathcal{E}_{1}\left(  \mathbb{D}^{2},\bm{\varphi}\right)  $
expectation value of the one electron part of the Hamiltonian. Hence the
energy $\mathcal{E}\left(  \mathbb{D}^{2},\bm{\varphi}\right)  $ can be
expressed as
\begin{equation}
\mathcal{E}\left(  \mathbb{D}^{2},\bm{\varphi}\right)  =\frac{1}{2}\sum_{1\leq
p\leq2s}\Omega_{p}I\left(  \varphi_{p}\right)  +\mathcal{E}_{1}\left(
\mathbb{D}^{2},\bm{\varphi}\right)  , \label{TEexpec}%
\end{equation}
where the ionization potential, $I\left(  \varphi_{p}\right)  $, of the GSO
$\left\vert \varphi_{p}\right\rangle $ is given by
\begin{equation}
I\left(  \varphi_{p}\right)  =\mathcal{E}\left(  \mathbb{D}^{2}%
,\bm{\varphi}\right)  -\mathcal{E}_{2N-1}\left(  \mathbb{D}^{2}%
,\bm{\varphi},\varphi_{p}\right)  .
\end{equation}


Although the optimized generalized Fock operator is self adjoint its matrix,
which is hermitian, is in general nondiagonal in the GSO basis. The matrix
$\left\{  F\left(  \mathbb{D}^{2},\bm{\varphi}\right)  _{qp}\right\}  $ can be
diagonalized by a unitary transformation $u$ of the GSO basis $\left\{
\left\vert \varphi_{j}\right\rangle \right\}  $ to produce
\begin{equation}
\tilde{F}\left(  \mathbb{D}^{2},\bm{\varphi}\right)  _{qp}=\frac
{1}{\left\langle \left.  \Psi\right\vert \Psi\right\rangle }\left\langle
\left.  \Psi\right\vert \left\{  b_{p}^{\dagger}b_{q}H-b_{p}^{\dagger}%
Hb_{q}\right\}  \Psi\right\rangle \delta_{qp},
\end{equation}
where
\begin{equation}
b_{p}^{\dagger}=ua_{p}^{\dagger}u^{\dagger}=\sum_{1\leq j\leq2s}a_{j}%
^{\dagger}u_{jp}%
\end{equation}
and
\begin{equation}
\left\vert \varsigma_{p}\right\rangle =\sum_{1\leq j\leq2s}\left\vert \psi
_{j}\right\rangle u_{jp}.
\end{equation}
The eigenvalues of the matrix $F\left(  \mathbb{D}^{2},\bm{\varphi}\right)  $
(the diagonal elements of the matrix $\tilde{F}\left(  \mathbb{D}%
^{2},\bm{\varphi}\right)  $ )
\begin{align}
\varepsilon_{p}  &  =\tilde{F}\left(  \mathbb{D}^{2},\bm{\varphi}\right)
_{pp}=\frac{1}{\left\langle \left.  \Psi\right\vert \Psi\right\rangle
}\left\{  \left\langle \left.  \Psi\right\vert b_{p}^{\dagger}b_{p}%
H\Psi\right\rangle -\left\langle \left.  b_{p}\Psi\right\vert Hb_{p}%
\Psi\right\rangle \right\} \nonumber\\
&  =\mathcal{N}_{p}\left\{  \mathcal{E}\left(  \mathbb{D}^{2}%
,\bm{\varphi }\right)  -\mathcal{E}_{2N-1}\left(  \mathbb{D}^{2}%
,\bm{\varphi,}\varsigma_{p}\right)  \right\}  ,
\end{align}
where the occupation of the GSO's $\left\{  \left\vert \varsigma
_{p}\right\rangle \right\}  $
\begin{equation}
\mathcal{N}_{p}=\frac{\left\langle \left.  \Psi\right\vert b_{p}^{\dagger
}b_{p}\Psi\right\rangle }{\left\langle \left.  \Psi\right\vert \Psi
\right\rangle }.
\end{equation}
Hence
\begin{equation}
\varepsilon_{p}=\mathcal{N}_{p}I\left(  \varsigma_{p}\right)  ,
\end{equation}
where $\left\{  I\left(  \varsigma_{p}\right)  \right\}  $ are the ionization
potentials associated with the GSO's $\left\{  \left\vert \varsigma
_{p}\right\rangle \right\}  .$ The transformation $u$ clearly transforms
occupied GSO's to a mixture of occupied GSO's and unoccupied to a mixture of
unoccupied, thus the optimized MCSCF state $\left\vert \Psi\right\rangle $
$=\left\vert \Psi\left(  \mathbb{D}^{2},\bm{\varphi}\right)  \right\rangle $
can be reexpressed in terms of the GSO's $\left\{  \left\vert \varsigma
_{p}\right\rangle \right\}  $ as
\begin{equation}
\left\vert \Psi\right\rangle =\left\vert \Psi\left(  \mathbb{\tilde{D}}%
^{2},\bm{\varsigma}\right)  \right\rangle ,
\end{equation}
where $\mathbb{\tilde{D}}^{2}$ is the matrix given by
\begin{equation}
\tilde{D}_{klij}^{2}=\frac{\left\langle \Psi\left\vert b_{i}^{\dagger}%
b_{j}^{\dagger}b_{l}b_{k}\Psi\right.  \right\rangle }{\left\langle \left.
\Psi\right\vert \Psi\right\rangle };\quad1\leq i,j,k,l\leq2s.
\end{equation}
The state $\left\vert \Psi\right\rangle $ is thus completely characterized by
the SORDM\ $\mathbb{D}^{2}$ and the basis $\left\{  \left\vert \varphi
_{p}\right\rangle \right\}  $ or by the SORDM\ $\mathbb{\tilde{D}}^{2}$ and
the basis $\left\{  \left\vert \varsigma_{p}\right\rangle \right\}  $. In
order to produce a more flexible notation we set
\begin{equation}
\mathbb{D}^{2}\left(  \bm{\varphi}\right)  =\mathbb{D}^{2}%
\end{equation}
and
\begin{equation}
\mathbb{D}^{2}\left(  \bm{\varsigma}\right)  =\mathbb{\tilde{D}}^{2},
\end{equation}
which leads to
\begin{equation}
\tilde{F}\left(  \mathbb{D}^{2},\bm{\varphi}\right)  \equiv F\left(
\mathbb{D}^{2}\left(  \bm{\varsigma}\right)  ,\bm{\varsigma}\right)
\end{equation}
and
\begin{equation}
F\left(  \mathbb{D}^{2},\bm{\varphi}\right)  \equiv F\left(  \mathbb{D}%
^{2}\left(  \bm{\varphi}\right)  ,\bm{\varphi}\right)  .
\end{equation}
One should note that the FORDO $D^{1}\left(  \Psi\right)  \equiv$
$D^{1}\left(  \mathbb{D}^{2},\bm{\varphi}\right)  \equiv D^{1}\left(
\mathbb{D}^{2}\left(  \bm{\varsigma}\right)  \right)  $ is still not in
general diagonal in the basis $\left\{  \left\vert \varsigma_{p}\right\rangle
\right\}  $ i.e. these GSO's are not NGSO's of $\left\vert \Psi\right\rangle $.

One can, however, also express the generalized Fock operator $F\left(
\mathbb{D}^{2},\bm{\varphi}\right)  $ in terms of the NGSO's $\left\{
\left\vert \upsilon_{j}\right\rangle \right\}  $ producing the generalized
Fock matrix $\mathbb{F}\left(  \mathbb{D}^{2}\left(  \bm{\upsilon}\right)
,\bm{\upsilon}\right)  ,$ that is not in general diagonal, but has diagonal
elements given by
\begin{equation}
F\left(  \mathbb{D}^{2}\left(  \bm{\upsilon}\right)  ,\bm{\upsilon }\right)
_{pp}=N_{p}\left\{  \mathcal{E}\left(  \mathbb{D}^{2}\left(
\bm{\upsilon}\right)  ,\bm{\upsilon}\right)  -\mathcal{E}_{2N-1}\left(
\mathbb{D}^{2}\left(  \bm{\upsilon}\right)  ,\bm{\upsilon,}\upsilon
_{p}\right)  \right\}  ,
\end{equation}
where $N_{p}$ is the occupation number of the NGSO\ $\left\vert \upsilon
_{p}\right\rangle $. As $\left\{  \left\vert \upsilon_{j}\right\rangle
\right\}  $ are NGSO's the expression for the one electron energy
$\mathcal{E}_{1}\left(  \mathbb{D}^{2},\bm{\varphi}\right)  $ $\equiv
\mathcal{E}_{1}\left(  \mathbb{D}^{2}\left(  \bm{\upsilon}\right)
,\bm{\upsilon}\right)  $ becomes%
\begin{equation}
\mathcal{E}_{1}\left(  \mathbb{D}^{2}\left(  \bm{\upsilon}\right)
,\bm{\upsilon}\right)  =\sum_{1\leq p\leq2s}N_{p}\left\langle \left.
\upsilon_{p}\right\vert h_{1\upsilon}\upsilon_{p}\right\rangle
\end{equation}
leading to the total state energy eq. $\left(  \ref{TEexpec}\right)  $ to be
expressed as
\begin{equation}
\mathcal{E}\left(  \mathbb{D}^{2}\left(  \bm{\upsilon}\right)
,\bm{\upsilon}\right)  =\mathcal{E}\left(  \mathbb{D}^{2}\left(
\bm{\varphi}\right)  ,\bm{\varphi}\right)  =\frac{1}{2}\sum_{1\leq p\leq
2s}N_{p}\left\{  I\left(  \upsilon_{p}\right)  +\left\langle \left.
\upsilon_{p}\right\vert h_{1\upsilon}\upsilon_{p}\right\rangle \right\}
=\frac{1}{2}%
%TCIMACRO{\tsum \limits_{1\leq p\leq2s}}%
%BeginExpansion
{\textstyle\sum\limits_{1\leq p\leq2s}}
%EndExpansion
N_{l}\kappa\left(  \upsilon_{p}\right)  , \label{Enkappa}%
\end{equation}
where we have introduced the new quantity,
\begin{equation}
\kappa\left(  \upsilon_{p}\right)  =I\left(  \upsilon_{p}\right)
+\left\langle \left.  \upsilon_{p}\right\vert h_{1\upsilon}\upsilon
_{p}\right\rangle . \label{kappa}%
\end{equation}
This new quantity is the sum of the ionization potential, the kinetic energy,
the electron - nuclear potential energy of attraction and potential energy of
interaction due to any other external forces of an electron described by a
NGSO\ $\left\vert \upsilon_{p}\right\rangle $. One should note that although
the optimized $2N$-electron state energy can be explicitly expressed in terms
of the one electron quantities $\left\{  \kappa\left(  \upsilon_{p}\right)
\right\}  $ associated with the GSO's $\left\{  \left\vert \upsilon
_{j}\right\rangle \right\}  $ and their occupation numbers $\left\{
N_{j}\right\}  $, the $2N$-electron state is \emph{not }in general determined
by these quantities.

In Appendix \ref{FockID} we further show that one can express the
$2N$-electron energy in terms of ionized $\left(  2N-1\right)  $-electron
state energies as
\begin{equation}
\mathcal{E}\left(  \mathbb{D}^{2},\bm{\varphi}\right)  =\frac{1}{2\left(
N-1\right)  }\sum_{1\leq p\leq2s}N_{p}\left(  \mathbb{D}^{2}%
,\bm{\varphi }\right)  \iota_{p}\left(  \mathbb{D}^{2},\bm{\varphi}\right)  ,
\end{equation}
where we have introduced another set of new quantities $\iota_{p}\left(
\mathbb{D}^{2},\bm{\varphi}\right)  $defined as
\begin{equation}
\iota_{p}\left(  \mathbb{D}^{2},\bm{\varphi}\right)  =\mathcal{E}_{p}\left(
\mathbb{D}^{2},\bm{\varphi}\right)  -\left\langle \left.  \varphi
_{p}\right\vert h_{1}\varphi_{p}\right\rangle
\end{equation}
and $N_{p}\left(  \mathbb{D}^{2},\bm{\varphi }\right)  $are the occupation
number numbers of the GSO's $\left\{  \left\vert \varphi_{p}\right\rangle
\right\}  $.

\subsection{Properties of an Optimized AGP\ Fock Operator\label{AGP_Fock}}

For AGP\ states the matrix elements of the generalized Fock operator becomes
(see Appendix \ref{AGP_Fock})
\begin{equation}
F\left(  \bm{n},\bm{\psi}\right)  _{qp}=\frac{1}{S_{N}\left(  g^{N}\right)
}\left\langle \left.  g^{N}\right\vert a_{p}^{\dagger}a_{q}H-a_{p}^{\dagger
}Ha_{q}g^{N}\right\rangle ,
\end{equation}
which has diagonal elements
\begin{equation}
F\left(  \bm{n},\bm{\psi}\right)  _{pp}=N_{p}\left\{  \mathcal{E}\left(
g^{N}\right)  -\frac{\left\langle \left.  a_{p}g^{N}\right\vert Ha_{p}%
g^{N}\right\rangle }{\left\langle \left.  a_{p}g^{N}\right\vert a_{p}%
g^{N}\right\rangle }\right\}  =N_{p}\left\{  \mathcal{E}\left(  g^{N}\right)
-\mathcal{E}_{2N-1}\left(  \psi_{p}\right)  \right\}  ,
\end{equation}
where $\mathcal{E}_{2N-1}\left(  \psi_{p}\right)  $ is the energy of the
normalized $2N-1$ electron state produced from the state $S_{N}\left(
g^{N}\right)  ^{-\frac{1}{2}}\left\vert g^{N}\right\rangle $ by removing an
electron described by the CGSO\ $\left\vert \psi_{p}\right\rangle $. As
$\left\{  \psi_{p}\right\}  $ are also NGSO's the energy $\mathcal{E}\left(
\bm{n},\bm{\psi}\right)  $ can be expressed as
\begin{equation}
\mathcal{E}\left(  \bm{n},\bm{\psi}\right)  =\frac{1}{2}\sum_{1\leq p\leq
2s}\left\{  N_{p}I\left(  \psi_{p}\right)  +N_{p}\left\langle \left.  \psi
_{p}\right\vert h_{1}\psi_{p}\right\rangle \right\}  \equiv\frac{1}{2}%
%TCIMACRO{\tsum \limits_{1\leq p\leq2s}}%
%BeginExpansion
{\textstyle\sum\limits_{1\leq p\leq2s}}
%EndExpansion
N_{p}\kappa\left(  \psi_{p}\right)  ,
\end{equation}
where GSO quantities $\left\{  \kappa\left(  \psi_{p}\right)  \right\}  $ are
given by
\begin{equation}
\kappa\left(  \psi_{p}\right)  =I\left(  \psi_{p}\right)  +\left\langle
\left.  \psi_{p}\right\vert h_{1}\psi_{p}\right\rangle .
\end{equation}
In contrast to the general MCSCF case the GSO's $\left\{  \left\vert \psi
_{p}\right\rangle \right\}  $ and their occupation $\left\{  N_{p}\right\}  $
numbers do determine the the $2N$-electron state AGP state.

\section{Reduction to the IPS \label{IPS}}

The standard basic model in chemistry is the IPS description, as we have noted
in the preceding, which is a special case of the AGP state model which becomes
a IPS\ in the case of integer occupation numbers i.e.
\begin{align}
N_{j}  &  =N_{j+s}=1,\quad1\leq j\leq N\nonumber\\
N_{j}  &  =N_{j+s}=0,\quad N+1\leq j\leq s
\end{align}
and%
\begin{align}
n_{j}  &  =n_{j+s}=1,\quad1\leq j\leq N\nonumber\\
n_{j}  &  =n_{j+s}=0,\quad N+1\leq j\leq s,
\end{align}
where the FORDM\ is given by
\begin{equation}
D_{jk}^{1}=N_{j}\delta_{jk};\quad1\leq j,k\leq2s.
\end{equation}
The SORDM for an IPS is given by
\begin{equation}
\mathbb{D}^{2}=\mathbb{D}^{1}\wedge\mathbb{D}^{1},
\end{equation}
where
\begin{equation}
\left(  \mathbb{D}^{1}\wedge\mathbb{D}^{1}\right)  _{ijkl}=D_{ik}^{1}%
D_{jl}^{1}-D_{il}^{1}D_{jk}^{1}.
\end{equation}
The generalized Fock operator that we have described in this paper
\begin{equation}
F\left(  \mathbb{D}^{2},\bm{\varphi}\right)  =\sum_{1\leq m,n\leq2s}\left\{
\left(  \mathbb{H}_{1}\left(  \bm{\varphi}\right)  \mathbb{D}^{1}\right)
_{mn}+\left(  L_{2}^{1}\left(  \mathbb{H}_{2}\left(  \bm{\varphi}\right)
\mathbb{D}^{2}\right)  \right)  _{mn}\right\}  \left\vert \varphi
_{m}\right\rangle \left\langle \varphi_{n}\right\vert , \label{F1}%
\end{equation}
can then be simplified by noting
\begin{align}
&  \left(  L_{2}^{1}\left(  \mathbb{H}_{2}\left(  \bm{\varphi}\right)
\mathbb{D}^{2}\right)  \right)  _{mn}=\sum_{1\leq k\leq2s}\left(
\mathbb{H}_{2}\left(  \bm{\varphi}\right)  \mathbb{D}^{2}\right)  _{mknk}%
=\sum_{1\leq p,q,k\leq2s}h_{2mkpq}\left(  \bm{\varphi}\right)  D_{pqnk}%
^{2}\nonumber\\
&  =\sum_{1\leq p,q,k\leq2s}h_{2mkpq}\left(  \bm{\varphi}\right)  \left(
D_{pn}^{1}D_{qk}^{1}-D_{pk}^{1}D_{qn}^{1}\right)  =\sum_{1\leq p,q,k\leq
2s}\left(  h_{2mkpq}\left(  \bm{\varphi}\right)  D_{qk}^{1}-h_{2mkqp}\left(
\bm{\varphi}\right)  D_{qk}^{1}\right)  D_{pn}^{1}.
\end{align}
We can then define the coulomb matrix as
\begin{equation}
J_{mp}\left(  \bm{\varphi}\right)  =\sum_{1\leq q,k\leq2s}h_{2mkpq}\left(
\bm{\varphi}\right)  D_{qk}^{1}%
\end{equation}
and exchange matrix as
\begin{equation}
K_{mp}\left(  \bm{\varphi}\right)  =\sum_{1\leq q,k\leq2s}h_{2mkqp}\left(
\bm{\varphi}\right)  D_{qk}^{1}%
\end{equation}
and obtain
\begin{equation}
\left(  L_{2}^{1}\left(  \mathbb{H}_{2}\left(  \bm{\varphi}\right)
\mathbb{D}^{2}\right)  \right)  _{mn}=\sum_{1\leq p\leq2s}\left\{
J_{mp}\left(  \bm{\varphi}\right)  -K_{mp}\left(  \bm{\varphi}\right)
\right\}  D_{pn}^{1}.
\end{equation}
The generalized Fock operator eq. $\left(  \ref{F1}\right)  $ can then be
expressed as
\begin{align}
F\left(  \bm{\varphi}\right)   &  =\sum_{1\leq m,n\leq2s}\left\{  \sum_{1\leq
q\leq2s}\left\{  h_{1mq}\left(  \bm{\varphi}\right)  +J_{mq}\left(
\bm{\varphi}\right)  -K_{mq}\left(  \bm{\varphi }\right)  \right\}  D_{qn}%
^{1}\right\}  \left\vert \varphi_{m}\right\rangle \left\langle \varphi
_{n}\right\vert \nonumber\\
&  =\sum_{1\leq m,n\leq2s}\left(  \bm{f}\left(  \bm{\varphi}\right)
\mathbb{D}^{1}\right)  _{mn}\left\vert \varphi_{m}\right\rangle \left\langle
\varphi_{n}\right\vert ,
\end{align}
where we have introduced the standard Fock matrix $\bm{f}\left(
\mathbb{D}^{1},\bm{\varphi}\right)  $ with matrix elements
\begin{equation}
f\left(  \bm{\varphi}\right)  _{mq}=h_{1mq}\left(  \bm{\varphi }\right)
+J_{mq}\left(  \bm{\varphi}\right)  -K_{mq}\left(  \bm{\varphi}\right)
\end{equation}
corresponding to the Fock operator
\begin{equation}
\mathfrak{f}\left(  \bm{\varphi}\right)  =\sum_{1\leq m,n\leq2s}f\left(
\bm{\varphi}\right)  _{mn}\left\vert \varphi_{m}\right\rangle \left\langle
\varphi_{n}\right\vert =\mathfrak{f}\left(  \bm{\varphi}\right)  ^{\dagger},
\end{equation}
which is always self adjoint. In the IPS case the generalized Fock operator
thus corresponds to the product of the conventional Fock operator and the
FORDO i.e.
\begin{equation}
F\left(  \bm{\varphi}\right)  =\mathfrak{f}\left(  \bm{\varphi }\right)
D^{1}.
\end{equation}


The stationarity condition
\begin{equation}
F\left(  \bm{\varphi}\right)  =F\left(  \bm{\varphi}\right)  ^{\dagger}%
\end{equation}
implies that
\begin{equation}
\left[  \mathfrak{f}\left(  \bm{\varphi}\right)  ,D^{1}\right]  =0. \label{HF}%
\end{equation}
As $D^{1}$ is a an orthogonal projection onto the subspace $\mathcal{H}_{o}$
generated by the occupied GSO's it defines the direct sum decomposition
\begin{equation}
\mathcal{H=H}_{o}\oplus\mathcal{H}_{u}, \label{DS}%
\end{equation}
where $\mathcal{H}_{u}$ is the subspace generated by the unoccupied GSO's.
Together with eq. $\left(  \ref{HF}\right)  $ this shows that the matrix
representation of $\mathfrak{f}\left(  \bm{\varphi}\right)  $ wrt the basis
$\left\{  \left\vert \varphi_{j}\right\rangle \right\}  $ is given by
\begin{equation}
\mathfrak{f}\left(  \bm{\varphi}\right)  \left\vert \bm{\varphi }\right\rangle
=\left\vert \bm{\varphi}\right\rangle \left(
\begin{array}
[c]{cc}%
\mathbb{F}_{o}\left(  \bm{\varphi}\right)  & \mathbb{O}\\
\mathbb{O} & \mathbb{F}_{u}\left(  \bm{\varphi}\right)
\end{array}
\right)  .
\end{equation}
Noting that the invariance group of an IPS always contains $U\left(
\mathcal{H}_{o}\right)  \otimes U\left(  \mathcal{H}_{u}\right)  $ leads to
usual Hartree-Fock eigenvalue equation
\begin{equation}
\mathfrak{f}\left(  \bm{\varphi}\right)  \left\vert \varphi_{j}\right\rangle
=\left\vert \varphi_{j}\right\rangle \varepsilon_{j};\quad1\leq j\leq2s.
\end{equation}


The standard aufbau principle is based on the eigenvalues $\left\{  \left.
\varepsilon_{j}\right\vert 1\leq j\leq2s\right\}  ,$ where the GSO's
corresponding to the lowest $2N$ eigenvalues are chosen to form an IPS\ state.
However, as is well known, this does not always lead to the lowest energy IPS.
A more reliable aufbau principle can be constructed by basing the choice of
GSO's on the one electron quantities, $\kappa\left(  \psi_{p}\right)  $,
introduced in eq. $\left(  \ref{kappa}\right)  ,$ and choosing $2N$ GSO's that
have the $2N$ lowest values of $\kappa\left(  \psi_{p}\right)  $. One is then
assured through eq. $\left(  \ref{Enkappa}\right)  $ to produce an IPS of
lowest energy.

\section{Discussion\label{Discussion}}

We have extended the one particle model of molecular quantum mechanics
described by IPS's to a more accurate correlated one particle model described
by AGP\ states for an even number of electrons and GAGP states for an odd
number. All of these states are determined by a set of one particle occupation
numbers and a set of CGSO's. In the IPS case the occupation numbers are one or
zero, while in the AGP case they must be at least doubly degenerate and for
the GAGP case there must be at most one nondegenerate integer occupation
number and the rest at least doubly degenerate.

The occupation numbers and CGSO's totally determine the many electron state, a
necessary criterion for a one particle model. In the AGP\ case the GSO's
produce weighted Probability Distributions (PD's)\ that produce pictorial
models of molecules that generalize the standard ones used in chemistry in
that they contain the effects of correlation. It is true that multi
configurational states also produce weighted Probability Distributions (PD's)
but they are \emph{not} determined by them, in contrast to the AGP\ case.

Another important aspect of the IPS model is the role of energies of the one
particle orbitals, which are given by the eigenvalues of the Fock operator.
When these are combined with the shape of PD's one obtains a description of
chemical properties i.e. reactivity and reaction mechanisms. Physically these
energies, are the poles of a first order approximation to the electron
propagator and correspond to the ionization potentials of the electrons from
the occupied orbitals. Mathematically they reflect the effect of the
orthonormality constraints in the IP-SCF\ procedure i.e. they give the rate of
change of the total energy with respect to changes in the normalization of the
corresponding orbitals. These properties are generalized by the optimized
MCSCF-Fock operator that we have described with two important differences: the
generalized Fock operator is not diagonal in the basis of GSO's and the
diagonal elements give weighted ionization potentials. This is a reflection of
our contention that the proper entity to generalize to correlated states in
order to obtain orbital derivatives is the product $\mathfrak{f}\left(
\bm{\varphi}\right)  D^{1}$ of the IPS Fock operator $\mathfrak{f}\left(
\bm{\varphi}\right)  $ and the FORDO $D^{1}$. In the IPS\ case the diagonal
elements of $\mathfrak{f}\left(  \bm{\varphi}\right)  D^{1}$ and
$\mathfrak{f}\left(  \bm{\varphi}\right)  $ in the basis $\left\{  \left\vert
\varphi_{j}\right\rangle \right\}  $, referring to the occupied GSO's, are
identical but this is not the case for AGP and MCSCF\ states. Moreover the
Fock operator $\mathfrak{f}\left(  \bm{\varphi}\right)  $ is always self
adjoint and both $\mathfrak{f}\left(  \bm{\varphi}\right)  D^{1}$ and
$\mathfrak{f}\left(  \bm{\varphi}\right)  $ are diagonal in this basis, thus
the diagonal elements are the eigenvalues of these operators. In the AGP and
MCSCF cases the generalized Fock operator $F\left(  \mathbb{D}^{2}%
,\bm{\varphi}\right)  $ is only self adjoint at stationary points of the GSO
variation and is not in general diagonal.

We have displayed how by analyzing the generalized Fock operator in the
NGSO\ basis, $\left\{  \varphi_{j}\right\}  $, for both the optimized MCSCF
and AGP states the total $2N$-electron state energy can be expressed as a sum
of occupation number, $\left\{  N_{j}\right\}  $, weighted observable
quantities $\left\{  \kappa\left(  \varphi_{j}\right)  \right\}  $, which are
the sum of the ionization potentials $\left\{  I\left(  \varphi_{j}\right)
\right\}  ,$ kinetic energy $\left\{  \left\langle \left.  \varphi
_{j}\right\vert -\nabla^{2}\varphi_{j}\right\rangle \right\}  $ and potential
energy $\left\{  \left\langle \left.  \varphi_{j}\right\vert V_{1}\varphi
_{j}\right\rangle \right\}  $ of electrons described by the NGSO's $\left\{
\varphi_{j}\right\}  $. The CGSO's, which are also NGSO's, of an optimized AGP
state together with the occupation numbers of the FORDO determine the
AGP\ state. Thus we have shown that an optimized~$2N$-electron AGP state and
its energy can be completely and explicitly expressed in terms of one electron
quantities and have produced a correlated one electron picture that can be
used in chemical modelling.

Another aspect of the IPS theory that is generalized by optimized AGP's is the
Brillouin condition
\begin{equation}
\left\langle \left.  g^{N}\right\vert Ha_{j}^{\dagger}a_{k}g^{N}\right\rangle
=0 \label{GBCD}%
\end{equation}
for $j=k\pm s$ and for all $j\neq k$ such that $c_{j}\neq c_{k}$. For
nondegenerate AGP's i.e. where all the canonical coefficients $\left\{
c_{j}\right\}  $ are distinct the condition eq. $\left(  \ref{GBCD}\right)  $
holds for all $j\neq k$. In particular eq. $\left(  \ref{GBCD}\right)  $ holds
for all occupied GSO's and should be contrasted to the optimized
MCSCF\ Brillouin condition
\begin{equation}
\left\langle \left.  \Psi\right\vert Ha_{p}^{\dagger}a_{h}\Psi\right\rangle
=0,
\end{equation}
where the subscript $h$ ranges over occupied GSO's, while the subscript $p$
over unoccupied ones.

By using subgroups of the general linear group of one electron Hilbert space
we have shown how to separate the variation over AGP states into occupation
number, phase and NGSO variations. The phase dependence is a reflection of the
fact the AGP FORDO (determined by the occupation numbers and NGSO's) does not
determine a unique AGP\ state, but only an equivalence class of states. By
optimizing over the phase for fixed reference FORDO's one obtains an explicit
Density Matrix Function (DMF)\ that can be used as a model functional for DMFT.

\section{Acknowledgements}

\appendix

\section{Energy Functionals\label{En_Funcs}}

In this appendix we obtain expressions for the total energy in terms of a
SORDM\ $\mathbb{D}^{2}$ and $\mathcal{U}\left(  \mathcal{H}^{1}\right)  $
group parameters $\bm{\alpha }$ that are constrained to form a hermitian
matrix. This condition can be easily enforced by expressing the $r\times r$
matrix $\bm{\alpha }$ in terms of $r^{2}$ unconstrained real parameters. We
show that this formulation of the energy is equivalent to one in terms of a
SORDM\ $\mathbb{D}^{2}$ and GSO's that are constrained to be orthonormalized.

The energy of a general $2N$-electron state, represented by the multi electron
positive operator, $D^{2N}$, is given by
\begin{equation}
\mathcal{E}\left(  D^{2N}\right)  =\frac{1}{\operatorname{Tr}\left\{
D^{2N}\right\}  }\operatorname{Tr}\left\{  \left\{  \sum_{1\leq i,k\leq
2s}h_{1ik}a_{i}^{\dagger}a_{k}+\frac{1}{4}\sum_{1\leq i,j,k,l\leq2s}%
V_{2ijkl}a_{i}^{\dagger}a_{j}^{\dagger}a_{l}a_{k}\right\}  D^{2N}\right\}  ,
\label{en1}%
\end{equation}
where
\begin{align}
h_{1ik}  &  =\left\langle \varphi_{i}\left\vert h_{1}\varphi_{k}\right.
\right\rangle \nonumber\\
V_{2ijkl}  &  =\left\langle ij||kl\right\rangle =\left\langle \varphi
_{i}\varphi_{j}\left\vert h_{2}\varphi_{k}\varphi_{l}\right.  \right\rangle
\end{align}
and the second quantized operators $\left\{  \left.  a_{j}^{\dagger
}\right\vert 1\leq j\leq2s\right\}  $ refer to the orthonormal reference basis
$\left\{  \left.  \left\vert \varphi_{j}\right\rangle \right\vert 1\leq
j\leq2s\right\}  $. For pure states
\begin{equation}
D^{2N}=\frac{\left\vert \Psi\right\rangle \left\langle \Psi\right\vert
}{\left\langle \Psi\left\vert \Psi\right.  \right\rangle }\equiv
\frac{\left\vert \Psi\right\rangle \left\langle \Psi\right\vert }%
{\mathcal{S}\left(  \Psi\right)  },
\end{equation}
where $\left\vert \Psi\right\rangle $ is in general an \emph{unnormalized
}vector\emph{ }state belonging to the $2N$-electron Hilbert space
$\mathcal{H}^{2N}.$ The expression eq. $\left(  \ref{en1}\right)  $ can be
reformulated as
\begin{equation}
\mathcal{E}\left(  \mathbb{D}^{2}\right)  =\operatorname{Tr}\left\{
\mathbb{H}_{1}\mathbb{D}^{1}\right\}  +\frac{1}{4}\operatorname{Tr}\left\{
\mathbb{V}_{2}\mathbb{D}^{2}\right\}
\end{equation}
in terms of the first and second order reduced density matrices $\mathbb{D}%
^{1}$ and $\mathbb{D}^{2}$ defined by
\begin{equation}
D_{klij}^{2}=\frac{1}{\operatorname{Tr}\left\{  D^{2N}\right\}  }%
\operatorname{Tr}\left\{  a_{i}^{\dagger}a_{j}^{\dagger}a_{l}a_{k}%
D^{2N}\right\}  ;\quad1\leq i,j,k,l\leq2s
\end{equation}
and
\begin{equation}
\mathbb{D}^{1}=\left(  2N-1\right)  ^{-1}L_{2}^{1}\left(  \mathbb{D}%
^{2}\right)  .
\end{equation}
The matrices $\mathbb{H}_{1}$ and $\mathbb{V}_{2}$ have the matrix elements
$\left\{  h_{1ik}\right\}  $ and $\left\{  \left\langle ij||kl\right\rangle
\right\}  $ respectively. The energy expression eq. $\left(  \ref{en1}\right)
$ can be used to define a functional, $\mathcal{E}\left(  \bm{\alpha }\right)
$, of orthonormal GSO's by considering a fixed reference state $D_{0}^{2N}$
and one electron unitary transformations $U\left(  \bm{\alpha
}\right)  $ by setting%
\begin{equation}
\mathcal{E}\left(  \bm{\alpha}\right)  =\frac{1}{\operatorname{Tr}\left\{
D_{0}^{2N}\right\}  }\operatorname{Tr}\left\{  \left\{  \sum_{1\leq i,k\leq
2s}h_{1ik}a_{i}^{\dagger}a_{k}+\frac{1}{4}\sum_{1\leq i,j,k,l\leq2s}%
V_{2ijkl}a_{i}^{\dagger}a_{j}^{\dagger}a_{l}a_{k}\right\}  U\left(
\bm{\alpha}\right)  ^{\dagger}D_{0}^{2N}U\left(  \bm{\alpha}\right)  \right\}
,
\end{equation}
where
\begin{equation}
U\left(  \bm{\alpha}\right)  =\exp\left\{  i\sum_{1\leq p,q\leq2s}\alpha
_{pq}a_{p}^{\dagger}a_{q}\right\}  ;\qquad\alpha_{pq}=\alpha_{qp}^{\ast}%
\quad1\leq p,q\leq2s.
\end{equation}
This leads to
\begin{align}
\mathcal{E}\left(  \bm{\alpha}\right)   &  =\frac{1}{\operatorname{Tr}\left\{
D_{0}^{2N}\right\}  }\operatorname{Tr}\left\{  \left\{  \sum_{1\leq i,k\leq
2s}h_{1ik}U\left(  \bm{\alpha}\right)  a_{i}^{\dagger}U\left(
\bm{\alpha}\right)  ^{\dagger}U\left(  \bm{\alpha}\right)  a_{k}U\left(
\bm{\alpha}\right)  ^{\dagger}\right.  \right. \nonumber\\
&  +\left.  \left.  \frac{1}{4}\sum_{1\leq i,j,k,l\leq2s}V_{2ijkl}U\left(
\bm{\alpha}\right)  a_{i}^{\dagger}U\left(  \bm{\alpha}\right)  ^{\dagger
}U\left(  \bm{\alpha}\right)  a_{j}^{\dagger}U\left(  \bm{\alpha}\right)
^{\dagger}U\left(  \bm{\alpha}\right)  a_{l}U\left(  \bm{\alpha}\right)
^{\dagger}U\left(  \bm{\alpha }\right)  a_{k}U\left(  \bm{\alpha}\right)
^{\dagger}D_{0}^{2N}\right\}  \right\}  ,
\end{align}
which can be rearranged to give
\begin{align}
\mathcal{E}\left(  \bm{\alpha}\right)   &  =\frac{1}{\operatorname{Tr}\left\{
D_{0}^{2N}\right\}  }\operatorname{Tr}\left\{  \left\{  \sum_{1\leq i,k\leq
2s}\left\langle \left.  \varphi_{i}\left(  \bm{\alpha}\right)  \right\vert
h_{1}\varphi_{k}\left(  \bm{\alpha}\right)  \right\rangle a_{i}^{\dagger}%
a_{k}\right.  \right. \nonumber\\
&  \qquad\qquad\qquad+\left.  \left.  \frac{1}{4}\sum_{1\leq i,j,k,l\leq
2s}\left\langle \left.  \varphi_{i}\left(  \bm{\alpha}\right)  \varphi
_{j}\left(  \bm{\alpha}\right)  \right\vert h_{2}\varphi_{k}\left(
\bm{\alpha}\right)  \varphi_{l}\left(  \bm{\alpha}\right)  \right\rangle
a_{i}^{\dagger}a_{j}^{\dagger}a_{l}a_{k}\right\}  D_{0}^{2N}\right\}
\end{align}
and
\begin{equation}
\mathcal{E}\left(  \bm{\alpha}\right)  =\operatorname{Tr}\left\{
\mathbb{H}_{1}\left(  \bm{\varphi}\left(  \bm{\alpha}\right)  \right)
\mathbb{D}_{0}^{1}\right\}  +\frac{1}{4}\operatorname{Tr}\left\{
\mathbb{V}_{2}\left(  \bm{\varphi}\left(  \bm{\alpha}\right)  \right)
\mathbb{D}_{0}^{2}\right\}  , \label{enmat}%
\end{equation}
where
\begin{equation}
\left\vert \varphi_{i}\left(  \bm{\alpha}\right)  \right\rangle =U\left(
\bm{\alpha}\right)  \left\vert \varphi_{i}\right\rangle .
\end{equation}
The expression eq. $\left(  \ref{enmat}\right)  $ utilizes the GSO dependent
matrices $\mathbb{H}_{1}\left(  \bm{\varphi}\left(  \bm{\alpha }\right)
\right)  $ and $\mathbb{V}_{2}\left(  \bm{\varphi}\left(  \bm{\alpha}\right)
\right)  $ whose matrix elements are
\begin{subequations}
\begin{align}
\mathbb{H}_{1}\left(  \bm{\varphi}\left(  \bm{\alpha}\right)  \right)   &
\leftrightarrow\left\{  \left\langle \left.  \varphi_{i}\left(
\bm{\alpha}\right)  \right\vert h_{1}\varphi_{k}\left(  \bm{\alpha }\right)
\right\rangle \right\} \\
\mathbb{V}_{2}\left(  \bm{\varphi}\left(  \bm{\alpha}\right)  \right)   &
\leftrightarrow\left\{  \left\langle \left.  \varphi_{i}\left(
\bm{\alpha}\right)  \varphi_{j}\left(  \bm{\alpha}\right)  \right\vert
h_{2}\varphi_{k}\left(  \bm{\alpha}\right)  \varphi_{l}\left(
\bm{\alpha}\right)  \right\rangle \right\}  .
\end{align}
Each new nonzero value of $\bm{\alpha}$ determines a new state $D^{2N}\left(
\bm{\alpha}\right)  =U\left(  \bm{\alpha}\right)  ^{\dagger}D_{0}^{2N}U\left(
\bm{\alpha}\right)  $ and new first and second order reduced density matrices
given by
\end{subequations}
\begin{align}
D^{2}\left(  \bm{\alpha}\right)  _{klij}  &  =\frac{1}{\operatorname{Tr}%
\left\{  D_{0}^{2N}\right\}  }\operatorname{Tr}\left\{  U\left(
\bm{\alpha}\right)  a_{i}^{\dagger}U\left(  \bm{\alpha}\right)  ^{\dagger
}U\left(  \bm{\alpha}\right)  a_{j}^{\dagger}U\left(  \bm{\alpha}\right)
^{\dagger}U\left(  \bm{\alpha}\right)  a_{l}U\left(  \bm{\alpha}\right)
^{\dagger}U\left(  \bm{\alpha }\right)  a_{k}U\left(  \bm{\alpha}\right)
^{\dagger}D_{0}^{2N}\right\} \nonumber\\
&  =\frac{1}{\operatorname{Tr}\left\{  D_{0}^{2N}\right\}  }\operatorname{Tr}%
\left\{  a_{i}^{\dagger}a_{j}^{\dagger}a_{l}a_{k}U\left(  \bm{\alpha }\right)
^{\dagger}D_{0}^{2N}U\left(  \bm{\alpha}\right)  \right\}  ;\quad1\leq
i,j,k,l\leq2s
\end{align}
and
\begin{align}
D^{1}\left(  \bm{\alpha}\right)  _{ji}  &  =\frac{1}{\operatorname{Tr}\left\{
D_{0}^{2N}\right\}  }\operatorname{Tr}\left\{  U\left(  \bm{\alpha}\right)
a_{i}^{\dagger}U\left(  \bm{\alpha}\right)  ^{\dagger}U\left(
\bm{\alpha}\right)  a_{j}U\left(  \bm{\alpha }\right)  ^{\dagger}D_{0}%
^{2N}\right\} \nonumber\\
&  =\frac{1}{\operatorname{Tr}\left\{  D_{0}^{2N}\right\}  }\operatorname{Tr}%
\left\{  a_{i}^{\dagger}a_{j}U\left(  \bm{\alpha}\right)  ^{\dagger}D_{0}%
^{2N}U\left(  \bm{\alpha}\right)  \right\}  .
\end{align}


The energy functional eq. $\left(  \ref{enmat}\right)  $ could be extended to
a full nonredundent energy functional by varying the reference
SORDM\ $\mathbb{D}_{0}^{2}$ in such a way as not to duplicate GSO changes.
However as the $N$-representability problem is still unsolved the domain of
variation of $\mathbb{D}_{0}^{2}$ is in general unknown, except in the AGP
state case, (which we study in this paper), where $N$-representable SORDM's
can be explicitly parameterized in terms of the occupation numbers $\left\{
\left.  n_{j}\right\vert 1\leq j\leq s\right\}  $. The energy functional eq.
$\left(  \ref{enmat}\right)  $ is of course also valid for SORDM's explicitly
constructed from multi configurational states, i.e. in terms of CI\ coefficients.

\section{Orbital Derivatives and Derivation of the Fock
Operator\label{Deriv_Fock}}

In this appendix we consider the functional derivatives of the total energy
functional wrt GSO's and express them in terms of a generalized Fock operator.
We then review a second quantized formulation of this generalized Fock operator.

\subsection{Direct Variation of GSO's}

In order to obtain derivative of terms in the energy expression it helpful to
utilize matrix elements wrt nonantisymmetric states
\begin{equation}
\left\vert \varphi_{j}\varphi_{k}\right)  =\left\vert \varphi_{j}\right)
\otimes\left\vert \varphi_{k}\right)  ,
\end{equation}
whose amplitudes are
\begin{equation}
\left(  \left.  \bm{r}_{1}\bm{,\xi}_{1}\bm{r}_{2}\bm{,\xi}_{2}\right\vert
\varphi_{j}\varphi_{k}\right)  =\varphi_{j}\left(  \bm{r}_{1}\bm{,\xi}_{1}%
\right)  \varphi_{k}\left(  \bm{r}_{2}\bm{,\xi}_{2}\right)  ,
\end{equation}
where
\begin{equation}
\varphi_{j}\left(  \bm{r}_{1}\bm{,\xi}_{1}\right)  =\left\langle \left.
\bm{r}_{1}\bm{,\xi}_{1}\right\vert \varphi_{j}\right\rangle .
\end{equation}
We will consider the energy in terms of the SORDM\ $\mathbb{D}^{2}$ and GSO's
$\bm{\varphi}$ given by
\begin{equation}
\mathcal{E}\left(  \mathbb{D}^{2},\bm{\varphi}\right)  =\operatorname{Tr}%
\left\{  \mathbb{H}_{1}\left(  \bm{\varphi}\right)  \mathbb{D}^{1}\right\}
+\frac{1}{4}\operatorname{Tr}\left\{  \mathbb{V}_{2}\left(
\bm{\varphi}\right)  \mathbb{D}^{2}\right\}  , \label{e1}%
\end{equation}
and we will concentrate on the direct variation of the GSO's. The first order
change in the energy due to changes $\left\{  \delta\bm{\varphi}^{\ast}%
,\delta\bm{\varphi}\right\}  $ in the GSO's is given by the Taylors series as
\begin{align}
\delta^{1}\mathcal{E}\left(  \mathbb{D}^{2},\bm{\varphi}\right)   &  =\left(
\partial_{\bm{\varphi}^{\ast}}^{1}\mathcal{E}\left(  \mathbb{D}^{2}%
,\bm{\varphi}\right)  ,\partial_{\bm{\varphi}}^{1}\mathcal{E}\left(
\mathbb{D}^{2},\bm{\varphi}\right)  \right)  \left(
\begin{array}
[c]{c}%
\delta\bm{\varphi}^{\ast}\\
\delta\bm{\varphi}
\end{array}
\right) \nonumber\\
&  =\left(  \partial_{\bm{\varphi}^{\ast}}^{1}\mathcal{E}\left(
\mathbb{D}^{2},\bm{\varphi}\right)  ,\partial_{\bm{\varphi}^{\ast}}%
^{1}\mathcal{E}\left(  \mathbb{D}^{2},\bm{\varphi}\right)  ^{\ast}\right)
\left(
\begin{array}
[c]{c}%
\delta\bm{\varphi}^{\ast}\\
\delta\bm{\varphi}
\end{array}
\right) \nonumber\\
&  =2\operatorname{Re}\left\{  \partial_{\bm{\varphi}^{\ast}}^{1}%
\mathcal{E}\left(  \mathbb{D}^{2},\bm{\varphi}\right)  ^{\dagger}%
\delta\bm{\varphi}^{\ast}\right\}  ,
\end{align}
where
\begin{equation}
\partial_{\varphi_{\mu}^{\ast}}^{1}\mathcal{E}\left(  \mathbb{D}%
^{2},\bm{\varphi}\right)  =\frac{\delta\mathcal{E}\left(  \mathbb{D}%
^{2},\bm{\varphi}\right)  }{\delta\varphi_{\mu}^{\ast}}.
\end{equation}
It is clear that
\begin{equation}
\partial_{\varphi_{\mu}^{\ast}}^{1}\mathcal{E}\left(  \mathbb{D}%
^{2},\bm{\varphi}\right)  =\left[  \partial_{\varphi_{\mu}}^{1}\mathcal{E}%
\left(  \mathbb{D}^{2},\bm{\varphi}\right)  \right]  ^{\ast}%
\end{equation}
as $\delta^{1}\mathcal{E}\left(  \mathbb{D}^{2},\bm{\varphi}\right)  $ must be
real valued for all $\left(  \delta\bm{\varphi}^{\ast},\delta
\bm{\varphi}\right)  $. In order to find the explicit expressions for the
derivatives of the energy one observes that:

\begin{enumerate}
\item The one electron functional derivatives are given by
\begin{equation}
\frac{\delta}{\delta\varphi_{\mu}^{\ast}}\left\langle \left.  \varphi
_{a}\right\vert h_{1}\varphi_{b}\right\rangle =h_{1}\left\vert \varphi
_{b}\right\rangle \delta_{a\mu}=\left\{  h_{1}\left\vert \varphi
_{b}\right\rangle \left\langle \varphi_{a}\right\vert \right\}  \left\vert
\varphi_{\mu}\right\rangle . \label{1partfuncder_o}%
\end{equation}
After noting that
\begin{equation}
h_{1}\left\vert \varphi_{b}\right\rangle \left\langle \varphi_{a}\right\vert
=\sum_{1\leq j,k\leq2s}h_{1jk}\left\vert \varphi_{j}\right\rangle \left\langle
\left.  \varphi_{k}\right\vert \varphi_{b}\right\rangle \left\langle
\varphi_{a}\right\vert =\sum_{1\leq j\leq2s}h_{1jb}\left\vert \varphi
_{j}\right\rangle \left\langle \varphi_{a}\right\vert
\end{equation}
eq. $\left(  \ref{1partfuncder_o}\right)  $ can be expanded as
\begin{equation}
\frac{\delta}{\delta\varphi_{\mu}^{\ast}}\left\langle \left.  \varphi
_{a}\right\vert h_{1}\varphi_{b}\right\rangle =\left\{  \sum_{1\leq m,n\leq
2s}h_{1mb}\delta_{an}\left\vert \varphi_{m}\right\rangle \left\langle
\varphi_{n}\right\vert \right\}  \left\vert \varphi_{\mu}\right\rangle .
\label{1partfuncder_ex}%
\end{equation}


\item The two electron functional derivatives are
\begin{align}
&  \frac{\delta}{\delta\varphi_{\mu}^{\ast}}\left(  \left.  \varphi_{a}%
\varphi_{b}\right\vert h_{2}\varphi_{c}\varphi_{d}\right)  =\left(  \left.
\varphi_{b}\right\vert h_{2}\varphi_{d}\right)  \left\vert \varphi
_{c}\right\rangle \delta_{a\mu}+\left(  \left.  \varphi_{a}\right\vert
h_{2}\varphi_{c}\right)  \left\vert \varphi_{d}\right\rangle \delta_{b\mu
}\nonumber\\
&  \qquad\qquad\qquad\left.  =\right.  \left\{  \left(  \left.  \varphi
_{b}\right\vert h_{2}\varphi_{d}\right)  \left\vert \varphi_{c}\right\rangle
\left\langle \varphi_{a}\right\vert +\left(  \left.  \varphi_{a}\right\vert
h_{2}\varphi_{c}\right)  \left\vert \varphi_{d}\right\rangle \left\langle
\varphi_{b}\right\vert \right\}  \left\vert \varphi_{\mu}\right\rangle ,
\label{2partfuncder_o}%
\end{align}
where the operators $\left\{  \left(  \left.  \varphi_{b}\right\vert
h_{2}\varphi_{d}\right)  ;1\leq b,d\leq2s\right\}  $ are defined by
\begin{equation}
\left(  \left.  \varphi_{b}\right\vert h_{2}\varphi_{d}\right)  =\sum_{1\leq
m,n\leq2s}\left(  \left.  \varphi_{m}\varphi_{b}\right\vert h_{2}\varphi
_{n}\varphi_{d}\right)  \left\vert \varphi_{m}\right\rangle \left\langle
\varphi_{n}\right\vert .
\end{equation}

\end{enumerate}

Eq. $\left(  \ref{1partfuncder_ex}\right)  $ can be further expanded to give
\begin{align}
\frac{\delta}{\delta\varphi_{\mu}^{\ast}}\left(  \left.  \varphi_{a}%
\varphi_{b}\right\vert h_{2}\varphi_{c}\varphi_{d}\right)   &  =\left\{
\sum_{1\leq k\leq2s}\left\{  h_{2kbcd}\left\vert \varphi_{k}\right\rangle
\left\langle \varphi_{a}\right\vert +h_{2kadc}\left\vert \varphi
_{k}\right\rangle \left\langle \varphi_{b}\right\vert \right\}  \right\}
\left\vert \varphi_{\mu}\right\rangle \nonumber\\
&  =\left\{  \sum_{1\leq m,n\leq2s}\left\{  h_{2mbcd}\delta_{an}%
+h_{2madc}\delta_{bn}\right\}  \left\vert \varphi_{m}\right\rangle
\left\langle \varphi_{n}\right\vert \right\}  \left\vert \varphi_{\mu
}\right\rangle . \label{2partfuncder_ex}%
\end{align}


The derivatives $\left\{  \frac{\delta\mathcal{E}\left(  \mathbb{D}%
^{2},\bm{\varphi}\right)  }{\delta\varphi_{\mu}^{\ast}}\right\}  $ can be
evaluated by using eqs. $\left(  \ref{1partfuncder_ex}\right)  $ and $\left(
\ref{2partfuncder_ex}\right)  $ to give
\begin{align}
&  \frac{\delta}{\delta\varphi_{\mu}^{\ast}}\left\{  \sum_{1\leq i,k\leq
2s}\left\langle \left.  \varphi_{i}\right\vert h_{1}\varphi_{k}\right\rangle
D_{ki}^{1}+\frac{1}{4}\sum_{1\leq i,j,k,l\leq2s}\left\{  \left(  \left.
\varphi_{i}\varphi_{j}\right\vert h_{2}\varphi_{k}\varphi_{l}\right)  -\left(
\left.  \varphi_{i}\varphi_{j}\right\vert h_{2}\varphi_{l}\varphi_{k}\right)
\right\}  D_{klij}^{2}\right\} \nonumber\\
&  =\sum_{1\leq m,n\leq2s}\left\{  \left\{  \sum_{1\leq k\leq2s}h_{1mk}%
D_{kn}^{1}\left\vert \varphi_{m}\right\rangle \left\langle \varphi
_{n}\right\vert \right.  \right. \nonumber\\
&  \qquad\qquad\qquad\qquad\qquad+\left.  \left.  \frac{1}{2}\sum_{1\leq
j,k,l\leq2s}\left\{  h_{2mjkl}D_{klnj}^{2}-h_{2mjlk}D_{klnj}^{2}\right\}
\right\}  \left\vert \varphi_{m}\right\rangle \left\langle \varphi
_{n}\right\vert \right\}  \left\vert \varphi_{\mu}\right\rangle ,
\end{align}
which can be written as
\begin{equation}
\frac{\delta\mathcal{E}\left(  \mathbb{D}^{2},\bm{\varphi}\right)  }%
{\delta\varphi_{\mu}^{\ast}}=\left\{  \sum_{1\leq m,n\leq2s}\left\{  \left(
\mathbb{H}_{1}\left(  \bm{\varphi}\right)  \mathbb{D}^{1}\right)
_{mn}+\left(  L_{2}^{1}\left(  \mathbb{H}_{2}\left(  \bm{\varphi}\right)
\mathbb{D}^{2}\right)  \right)  _{mn}\right\}  \left\vert \varphi
_{m}\right\rangle \left\langle \varphi_{n}\right\vert \right\}  \left\vert
\varphi_{\mu}\right\rangle ,
\end{equation}
where $\mathbb{H}_{2}\left(  \bm{\varphi}\right)  $ is the nonantisymmetric
matrix
\begin{equation}
\mathbb{H}_{2}\left(  \bm{\varphi}\right)  \leftrightarrow\left\{
h_{2ijkl}\right\}  ,
\end{equation}
that determine the antisymmetric matrix elements
\begin{equation}
V_{2ijkl}\left(  \bm{\varphi}\right)  =\left(  h_{2ijkl}\left(
\bm{\varphi}\right)  -h_{2ijlk}\left(  \bm{\varphi}\right)  \right)  .
\end{equation}
This allows us to define a generalized Fock operator
\begin{subequations}
\label{Fock_Op}%
\begin{align}
F\left(  \mathbb{D}^{2},\bm{\varphi}\right)   &  =\sum_{1\leq m,n\leq
2s}\left\{  \left(  \mathbb{H}_{1}\left(  \bm{\varphi}\right)  \mathbb{D}%
^{1}\right)  _{mn}+\frac{1}{2}\left(  L_{2}^{1}\left(  \mathbb{V}_{2}\left(
\bm{\varphi}\right)  \mathbb{D}^{2}\right)  \right)  _{mn}\right\}  \left\vert
\varphi_{m}\right\rangle \left\langle \varphi_{n}\right\vert \label{Fock_Op_a}%
\\
&  =\sum_{1\leq m,n\leq2s}\left\{  \left(  \mathbb{H}_{1}\left(
\bm{\varphi}\right)  \mathbb{D}^{1}\right)  _{mn}+\left(  L_{2}^{1}\left(
\mathbb{H}_{2}\left(  \bm{\varphi}\right)  \mathbb{D}^{2}\right)  \right)
_{mn}\right\}  \left\vert \varphi_{m}\right\rangle \left\langle \varphi
_{n}\right\vert , \label{Fock_Op_b}%
\end{align}
that produces the functional derivatives
\end{subequations}
\begin{equation}
\frac{\delta\mathcal{E}\left(  \mathbb{D}^{2},\bm{\varphi}\right)  }%
{\delta\varphi_{\mu}^{\ast}}=F\left(  \mathbb{D}^{2},\bm{\varphi}\right)
\left\vert \varphi_{\mu}\right\rangle ;\quad1\leq\mu\leq2s.
\label{FuncDerFock}%
\end{equation}
As
\begin{equation}
L_{2}^{1}\left(  \mathbb{A}^{\dagger}\right)  =L_{2}^{1}\left(  \mathbb{A}%
\right)  ^{\dagger}%
\end{equation}
the adjoint of the Fock operator eq. $\left(  \ref{Fock_Op}\right)  $ is given
by
\begin{equation}
F\left(  \mathbb{D}^{2},\bm{\varphi}\right)  ^{\dagger}=\sum_{1\leq m,n\leq
2s}\left\{  \left(  \mathbb{D}^{1}\mathbb{H}_{1}\left(  \bm{\varphi }\right)
\right)  _{mn}+\left(  L_{2}^{1}\left(  \mathbb{D}^{2}\mathbb{H}_{2}\left(
\bm{\varphi}\right)  \right)  \right)  _{mn}\right\}  \left\vert \varphi
_{m}\right\rangle \left\langle \varphi_{n}\right\vert . \label{Fock_Op_adj}%
\end{equation}
In general the generalized Fock operator is not self adjoint i.e.
\begin{equation}
F\left(  \mathbb{D}^{2},\bm{\varphi}\right)  ^{\dagger}\neq F\left(
\mathbb{D}^{2},\bm{\varphi}\right)  .
\end{equation}


\subsection{Matrix Elements of $F\left(  \bm{\varphi}\right)  $ in Terms of
Second Quantized Operators\label{SQ_Fock}}

In order to obtain an expression for the matrix elements $\left\{  \left.
F\left(  \mathbb{D}^{2},\bm{\varphi}\right)  _{jk}\right\vert 1\leq
j,k\leq2s\right\}  $ in terms of the second quantized operators $\left\{
\left.  a_{j}^{\dagger}\right\vert 1\leq j,k\leq2s\right\}  ,$ that refer to
the GSO's $\left\{  \left.  \left\vert \varphi_{j}\right\rangle \right\vert
1\leq j,k\leq2s\right\}  $, we consider the expectation value of the
commutators $\left\{  \left.  a_{p}^{\dagger}\left[  a_{q},H\right]
\right\vert 1\leq j,k\leq2s\right\}  ,$ where
\begin{equation}
H=h_{1}+\frac{1}{4}V_{2}=\sum_{1\leq i,k\leq2s}h_{1ik}a_{i}^{\dagger}%
a_{k}+\frac{1}{4}\sum_{1\leq i,j,k,l\leq2s}V_{2ijkl}a_{i}^{\dagger}%
a_{j}^{\dagger}a_{l}a_{k},
\end{equation}%
\begin{equation}
h_{1}=\sum_{1\leq i,k\leq2s}h_{1ik}a_{i}^{\dagger}a_{k}%
\end{equation}
and
\begin{equation}
V_{2}=\sum_{1\leq i,j,k,l\leq2s}V_{2ijkl}a_{i}^{\dagger}a_{j}^{\dagger}%
a_{l}a_{k}.
\end{equation}
The one electron part of $H$ gives
\begin{equation}
a_{p}^{\dagger}\sum_{1\leq i,k\leq2s}h_{1ik}\left[  a_{q},a_{i}^{\dagger}%
a_{k}\right]  =\sum_{1\leq k\leq2s}h_{1qk}a_{p}^{\dagger}a_{k},
\end{equation}
while the two electron interaction part gives
\begin{equation}
a_{p}^{\dagger}\sum_{1\leq i,j,k,l\leq2s}V_{2ijkl}\left[  a_{q},a_{i}%
^{\dagger}a_{j}^{\dagger}a_{l}a_{k}\right]  =2\sum_{1\leq j,k,l\leq
2s}V_{2qjkl}a_{p}^{\dagger}a_{j}^{\dagger}a_{l}a_{k}.
\end{equation}
The expectation values of the commutators can be expanded as
\begin{align}
&  \frac{1}{\operatorname{Tr}\left\{  D^{2N}\right\}  }\operatorname{Tr}%
\left\{  a_{p}^{\dagger}\left[  a_{q},H\right]  D^{2N}\right\}  =\nonumber\\
&  \qquad\qquad\qquad\frac{1}{\operatorname{Tr}\left\{  D^{2N}\right\}
}\operatorname{Tr}\left\{  \left\{  \sum_{1\leq k\leq2s}h_{1qk}a_{p}^{\dagger
}a_{k}+\frac{1}{2}\sum_{1\leq j,k,l\leq2s}V_{2qjkl}a_{p}^{\dagger}%
a_{j}^{\dagger}a_{l}a_{k}\right\}  D^{2N}\right\}  ,
\end{align}
giving the important identity
\begin{equation}
F\left(  \mathbb{D}^{2},\bm{\varphi}\right)  _{qp}=\frac{1}{\operatorname{Tr}%
\left\{  D^{2N}\right\}  }\operatorname{Tr}\left\{  a_{p}^{\dagger}\left[
a_{q},H\right]  D^{2N}\right\}  =\operatorname{Tr}\left\{  \left(
\mathbb{H}_{1}\left(  \bm{\varphi}\right)  \mathbb{D}^{1}\right)  _{qp}%
+\frac{1}{2}\left(  L_{2}^{1}\left(  \mathbb{V}_{2}\left(
\bm{\varphi}\right)  \mathbb{D}^{2}\right)  \right)  _{qp}\right\}  .
\label{Fockelem}%
\end{equation}


\section{Variation of Group Parameters\label{VarGroupPar}}

In this appendix we consider the GSO variations in terms of one electron
unitary group, $\mathcal{U}\left(  \mathcal{H}^{1}\right)  ,$ parameters and
obtain the derivatives of the total energy functional wrt to these parameters
and relate them to the generalized Fock operator obtained in Appendix
\ref{Deriv_Fock}.

The energy of a general $2N$-electron state can be expressed as
\begin{equation}
\mathcal{E}\left(  \mathbb{D}^{2},\bm{\alpha}\right)  =\operatorname{Tr}%
\left\{  \left\{  \sum_{1\leq i,j\leq2s}h_{1ij}a_{i}^{\dagger}a_{j}+\frac
{1}{4}\sum_{1\leq i,j,k,l\leq2s}V_{2ijkl}a_{i}^{\dagger}a_{j}^{\dagger}%
a_{l}a_{k}\right\}  U\left(  \bm{\alpha}\right)  ^{\dagger}D^{2N}U\left(
\bm{\alpha}\right)  \right\}  ,
\end{equation}
where
\begin{equation}
U\left(  \bm{\alpha}\right)  =\exp\left\{  i\sum_{1\leq p,q\leq2s}\alpha
_{pq}a_{p}^{\dagger}a_{q}\right\}  ;\qquad\alpha_{pq}=\alpha_{qp}^{\ast}%
\quad1\leq p,q\leq2s.
\end{equation}
in terms of the unitary group $U\left(  \mathcal{H}^{1}\right)  $ of
transformations $\left\{  U\left(  \bm{\alpha}\right)  \right\}  $ of the
Hilbert space of one electron states $\mathcal{H}^{1}$ and SORDM's
$\mathbb{D}^{2}$. The Taylor series expansion gives
\begin{align}
&  \mathcal{E}\left(  \mathbb{D}^{2},\bm{\alpha}\right)  \simeq
\operatorname{Tr}\left\{  \sum_{1\leq i,j\leq2s}h_{1ij}a_{i}^{\dagger}%
a_{j}+\frac{1}{4}\sum_{1\leq i,j,k,l\leq2s}V_{2ijkl}\operatorname{Tr}%
a_{i}^{\dagger}a_{j}^{\dagger}a_{l}a_{k}\right. \nonumber\\
&  \qquad\qquad\qquad\qquad\qquad\left.  \left\{  1-i\sum_{1\leq p,q\leq
2s}\alpha_{pq}a_{p}^{\dagger}a_{q}\right\}  D^{2N}\left\{  1+i\sum_{1\leq
p,q\leq2s}\alpha_{pq}a_{p}^{\dagger}a_{q}\right\}  \right\} \nonumber\\
&  \simeq\mathcal{E}\left(  \mathbb{D}^{2},\bm{0}\right)  -i\sum_{1\leq
p,q\leq2s}\operatorname{Tr}\left\{  \left\{  \sum_{1\leq i,j\leq2s}%
h_{1ij}\left[  a_{i}^{\dagger}a_{j},a_{p}^{\dagger}a_{q}\right]  \right.
\right. \nonumber\\
&  \qquad\qquad\qquad\qquad\qquad\qquad+\left.  \left.  \frac{1}{4}\sum_{1\leq
i,j,k,l\leq2s}V_{2ijkl}\left[  a_{i}^{\dagger}a_{j}^{\dagger}a_{l}a_{k,}%
a_{p}^{\dagger}a_{q}\right]  \right\}  D^{2N}\right\}  \alpha_{pq}.
\end{align}
Thus showing that the first order change of the energy functional
wrt \bm{\alpha}$ evaluated around \bm{\alpha}$=\bm{0}$ is
\begin{align}
\delta^{1}\mathcal{E}\left(  \mathbb{D}^{2},\bm{0}\right)   &  =-i\sum_{1\leq
p,q\leq2s}\operatorname{Tr}\left\{  \left\{  \sum_{1\leq i,j\leq2s}%
h_{1ij}\left[  a_{i}^{\dagger}a_{j},a_{p}^{\dagger}a_{q}\right]  \right.
\right. \nonumber\\
&  \qquad\qquad+\left.  \left.  \frac{1}{4}\sum_{1\leq i,j,k,l\leq2s}%
V_{2ijkl}\left[  a_{i}^{\dagger}a_{j}^{\dagger}a_{l}a_{k,}a_{p}^{\dagger}%
a_{q}\right]  \right\}  \left.  D^{2N}\right\}  \right\}  \alpha_{pq},
\end{align}
displaying that the derivatives wrt $\left\{  \left.  \alpha_{pq}\right\vert
1\leq p,q\leq2s\right\}  $ are
\begin{align}
\frac{\partial\mathcal{E}\left(  \mathbb{D}^{2},\bm{0}\right)  }%
{\partial\alpha_{pq}}  &  =-i\operatorname{Tr}\left\{  \left\{  \sum_{1\leq
i,j\leq2s}h_{1ij}\left[  a_{i}^{\dagger}a_{j},a_{p}^{\dagger}a_{q}\right]
\right.  \right. \nonumber\\
&  \qquad\qquad+\left.  \left.  \frac{1}{4}\sum_{1\leq i,j,k,l\leq2s}%
V_{2ijkl}\left[  a_{i}^{\dagger}a_{j}^{\dagger}a_{l}a_{k},a_{p}^{\dagger}%
a_{q}\right]  \right\}  D^{2N}\right\}  .
\end{align}
The commutators give
\begin{subequations}
\begin{align}
\left[  a_{i}^{\dagger}a_{j},a_{p}^{\dagger}a_{q}\right]   &  =a_{i}^{\dagger
}a_{q}\delta_{pj}-a_{p}^{\dagger}a_{j}\delta_{pj}\\
\left[  a_{i}^{\dagger}a_{j}^{\dagger}a_{l}a_{k},a_{p}^{\dagger}a_{q}\right]
&  =a_{i}^{\dagger}a_{j}^{\dagger}a_{q}a_{k}\delta_{pl}-a_{p}^{\dagger}%
a_{j}^{\dagger}a_{l}a_{k}\delta_{iq}-a_{i}^{\dagger}a_{p}^{\dagger}a_{l}%
a_{k}\delta_{jq}+a_{i}^{\dagger}a_{j}^{\dagger}a_{l}a_{q}\delta_{pk}%
\end{align}
leading to
\end{subequations}
\begin{align}
&  \frac{\partial\mathcal{E}\left(  \mathbb{D}^{2},\bm{0}\right)  }%
{\partial\alpha_{pq}}=-i\operatorname{Tr}\left\{  \left\{  \sum_{1\leq
i,j\leq2s}h_{1ij}\left(  a_{i}^{\dagger}a_{q}\delta_{pj}-a_{p}^{\dagger}%
a_{j}\delta_{iq}\right)  \right.  \right.  +\nonumber\\
&  \left.  \left.  \frac{1}{4}\sum_{1\leq i,j,k,l\leq2s}V_{2ijkl}\left(
a_{i}^{\dagger}a_{j}^{\dagger}a_{q}a_{k}\delta_{pl}-a_{p}^{\dagger}%
a_{j}^{\dagger}a_{l}a_{k}\delta_{iq}-a_{i}^{\dagger}a_{p}^{\dagger}a_{l}%
a_{k}\delta_{jq}+a_{i}^{\dagger}a_{j}^{\dagger}a_{l}a_{q}\delta_{pk}\right)
\right\}  D^{2N}\right\}  .
\end{align}
This produces%
\begin{align}
\frac{\partial\mathcal{E}\left(  \mathbb{D}^{2},\bm{0}\right)  }%
{\partial\alpha_{pq}}=-  &  i\operatorname{Tr}\left\{  \sum_{1\leq j\leq
2s}\left\{  \left\{  h_{1jp}a_{j}^{\dagger}a_{q}-h_{1qj}a_{p}^{\dagger}%
a_{j}\right\}  \right.  \right. \nonumber\\
&  +\left.  \left.  \frac{1}{2}\sum_{1\leq j,k,l\leq2s}\left\{  V_{2ijkl}%
a_{l}^{\dagger}a_{j}^{\dagger}a_{q}a_{k}-V_{2ijkl}a_{p}^{\dagger}%
a_{j}^{\dagger}a_{l}a_{k}\right\}  \right\}  D^{2N}\right\}  ,
\end{align}
giving%
\begin{equation}
\frac{\partial\mathcal{E}\left(  \mathbb{D}^{2},\bm{0}\right)  }%
{\partial\alpha_{pq}}=-i\left\{  \sum_{1\leq j\leq2s}\left\{  D_{qj}%
^{1}h_{1jp}-h_{1qj}D_{jp}^{1}\right\}  +\frac{1}{2}\sum_{1\leq i,j,k\leq
2s}\left\{  D_{kqlj}^{2}V_{2ijkl}-V_{2ijkl}D_{klpj}^{2}\right\}  \right\}
\end{equation}
and finally
\begin{equation}
\frac{\partial\mathcal{E}\left(  \mathbb{D}^{2},\bm{0}\right)  }%
{\partial\alpha_{pq}}=-i\left(  \left[  \mathbb{D}^{1},\mathbb{H}_{1}\left(
\bm{\varphi}\left(  \bm{\alpha}\right)  \right)  \right]  +\frac{1}{2}%
L_{2}^{1}\left(  \left[  \mathbb{D}^{2},\mathbb{V}_{2}\left(
\bm{\varphi}\left(  \bm{\alpha}\right)  \right)  \right]  \right)  \right)
_{qp}. \label{derpq}%
\end{equation}
The derivatives $\left\{  \frac{\partial\mathcal{E}\left(  \mathbb{D}%
^{2},\bm{0}\right)  }{\partial\alpha_{pq}}\right\}  $ can also be expressed in
terms of the nonantisymmetric matrix $\mathbb{H}_{2}\left(
\bm{\varphi}\left(  \bm{\alpha}\right)  \right)  $ as
\begin{equation}
\frac{\partial\mathcal{E}\left(  \mathbb{D}^{2},\bm{0}\right)  }%
{\partial\alpha_{pq}}=-i\left(  \left[  \mathbb{D}^{1},\mathbb{H}_{1}\left(
\bm{\varphi}\left(  \bm{\alpha}\right)  \right)  \right]  +L_{2}^{1}\left(
\left[  \mathbb{D}^{2},\mathbb{H}_{2}\left(  \bm{\varphi }\left(
\bm{\alpha}\right)  \right)  \right]  \right)  \right)  _{qp}. \label{derpq1}%
\end{equation}
If we compare eq. $\left(  \ref{derpq}\right)  $ with eqs. $\left(
\ref{Fock_Op}\right)  $ and $\left(  \ref{Fock_Op_adj}\right)  $ one observes
the important result that
\begin{equation}
\frac{\partial\mathcal{E}\left(  \mathbb{D}^{2},\bm{0}\right)  }%
{\partial\alpha_{pq}}=i\left(  F\left(  \bm{\varphi}\left(
\bm{\alpha }\right)  \right)  -F\left(  \bm{\varphi}\left(
\bm{\alpha}\right)  \right)  ^{\dagger}\right)  _{qp},
\end{equation}
where
\begin{equation}
\left\vert \varphi_{j}\left(  \bm{\alpha}\right)  \right\rangle =U\left(
\bm{\alpha}\right)  \left\vert \varphi_{j}\right\rangle ;\quad1\leq j\leq2s.
\end{equation}


\section{Stationarity Conditions\label{StatCon}}

In this appendix we relate the stationarity of the total energy functional wrt
to orthonormally constrained variations of the GSO's to the stationarity wrt
unconstrained variations of the one electron group parameters. As a result one
can observe that the representation of the generalized Fock operator wrt the
basis of optimal GSO's is hermitian and is in fact the matrix of Lagrange
parameters associated with these GSO's.

At a stationary point wrt unitary group parameters $\bm{\alpha}$ the
derivatives eq. $\left(  \ref{derpq}\right)  $ must equal zero hence
\begin{equation}
\left(  F\left(  \bm{\varphi}\left(  \bm{\alpha}\right)  \right)  -F\left(
\bm{\varphi}\left(  \bm{\alpha}\right)  \right)  ^{\dagger}\right)
_{qp}=0\quad1\leq p,q\leq2s
\end{equation}
and the generalized Fock operator is self adjoint i.e.
\begin{equation}
F\left(  \bm{\varphi}\left(  \bm{\alpha}\right)  \right)  =F\left(
\bm{\varphi}\left(  \bm{\alpha}\right)  \right)  ^{\dagger}. \label{adjcon}%
\end{equation}
In order to constrain the GSO's to be orthonormal in a direct nonunitary group
variational procedure one considers the Lagrangian
\begin{equation}
\mathcal{L}\left(  \mathbb{D}^{2},\bm{\varphi}\right)  =\mathcal{E}\left(
\mathbb{D}^{2},\bm{\varphi}\right)  -\sum_{1\leq j,k\leq2s}\varepsilon
_{kj}\left\{  \left\langle \left.  \varphi_{j}\right\vert \varphi
_{k}\right\rangle -\delta_{jk}\right\}  ,
\end{equation}
where $\left\{  \left.  \varepsilon_{jk}\right\vert 1\leq j,k\leq2s\right\}  $
are Lagrange multipliers, which produces the constrained stationarity
conditions
\begin{equation}
\frac{\delta\mathcal{L}\left(  \mathbb{D}^{2},\bm{\varphi}\right)  }%
{\delta\varphi_{j}^{\ast}}=\frac{\delta\mathcal{E}\left(  \mathbb{D}%
^{2},\bm{\varphi}\right)  }{\delta\varphi_{j}^{\ast}}-\sum_{1\leq k\leq
2s}\left\vert \varphi_{k}\right\rangle \varepsilon_{kj}=0;\quad1\leq j\leq2s.
\label{Lagran}%
\end{equation}
Using eq. $\left(  \ref{FuncDerFock}\right)  $ eq. $\left(  \ref{Lagran}%
\right)  $ can be expressed as
\begin{equation}
F\left(  \mathbb{D}^{2},\bm{\varphi}\right)  \left\vert \varphi_{j}%
\right\rangle =\sum_{1\leq k\leq2s}\left\vert \varphi_{k}\right\rangle
\varepsilon_{kj};\quad1\leq j\leq2s,
\end{equation}
which in combination with eq. $\left(  \ref{adjcon}\right)  $ shows that the
matrix of Lagrange multipliers is hermitian when $\bm{\varphi}$ is a
constrained stationary point.

\section{Derivatives of the AGP\ Energy Functional\label{variational}}

In this appendix we focus on the special case of the AGP energy functional and
consider its derivatives wrt two different sets of coordinates (a) the geminal
occupation numbers, NGSO's and phases and (b) the group parameters of
subgroups of the one electron general linear group and phases.

The $2N$-AGP\ energy functional in terms of the NGSO's $\left\{  \left.
\varphi_{j}\right\vert 1\leq j\leq2s\right\}  $, the geminal occupation
numbers $\left\{  \left.  n_{j}\right\vert 1\leq j\leq s\right\}  $ and the
phases $\left\{  \left.  \theta_{j}\right\vert 1\leq j\leq s\right\}  $ can be
expressed as the quotient%
\begin{align}
&  \mathcal{E}\left(  \bm{n},\bm{\varphi},\bm{\theta}\right)  =\frac
{\mathcal{H}\left(  \bm{n},\bm{\varphi},\bm{\theta}\right)  }{S_{N}\left(
\bm{n}\right)  }\nonumber\\
&  =\frac{1}{S_{N}\left(  \bm{n}\right)  }\left\{  \sum_{1\leq j\leq s}%
n_{j}\left\{  \left\langle \varphi_{j}\left\vert h_{1}\varphi_{j}\right.
\right\rangle +\left\langle \varphi_{j+s}\left\vert h_{1}\varphi_{j+s}\right.
\right\rangle +\left\langle \varphi_{j}\varphi_{j+s}\left\Vert \varphi
_{j}\varphi_{j+s}\right.  \right\rangle \right\}  \right. \nonumber\\
&  +\sum_{1\leq j,k\leq s}n_{j}^{\frac{1}{2}}n_{k}^{\frac{1}{2}}e^{i\left(
\theta_{j}-\theta_{k}\right)  }S_{N-1}\left(  \jmath k\right)  \varphi
_{j}\varphi_{j+s}|\left\vert \varphi_{k}\varphi_{k+s}\right\rangle \left.
+\frac{1}{2}\sum_{1\leq j,k\leq s}n_{j}n_{k}S_{N-2}\left(  \jmath k\right)
\left\langle \varphi_{j}\varphi_{k}|||\varphi_{j}\varphi_{k}\right\rangle
\right\}  .
\end{align}
The first order variation $\delta^{1}\mathcal{E}\left(
\bm{n},\bm{\varphi},\bm{\theta}\right)  $ of the energy is
\begin{equation}
\delta^{1}\mathcal{E}\left(  \bm{n},\bm{\varphi},\bm{\theta }\right)  =\left(
%
\begin{array}
[c]{cccc}%
\partial_{\bm{n}}^{1}\mathcal{E}\left(
\bm{n},\bm{\varphi },\bm{\theta}\right)  & \partial_{\bm{\varphi}^{\ast}}%
^{1}\mathcal{E}\left(  \bm{n},\bm{\varphi},\bm{\theta}\right)  &
\partial_{\bm{\varphi}}^{1}\mathcal{E}\left(
\bm{n},\bm{\varphi },\bm{\theta}\right)  & \partial_{\bm{\theta}}%
^{1}\mathcal{E}\left(  \bm{n},\bm{\varphi},\bm{\theta}\right)
\end{array}
\right)  \left(
\begin{array}
[c]{c}%
\delta\bm{n}\\
\delta\bm{\varphi}^{\ast}\\
\delta\bm{\varphi}\\
\delta\bm{\theta}
\end{array}
\right)  ,
\end{equation}
where
\begin{subequations}
\label{Fderiv}%
\begin{align}
\partial_{\bm{n}}^{1}\mathcal{E}\left(
\bm{n},\bm{\varphi },\bm{\theta}\right)   &  =\frac{1}{S_{N}\left(
\bm{n}\right)  }\left\{  \partial_{\bm{n}}^{1}\mathcal{H}\left(
\bm{n},\bm{\varphi},\bm{\theta}\right)  -\mathcal{E}\left(
\bm{n},\bm{\varphi},\bm{\theta}\right)  \partial_{\bm{n}}^{1}S_{N}\left(
\bm{n}\right)  \right\} \\
\partial_{\bm{\phi}^{\ast}}^{1}\mathcal{E}\left(
\bm{n},\bm{\varphi},\bm{\theta}\right)   &  =\frac{1}{S_{N}\left(
\bm{n}\right)  }\left\{  \partial_{\bm{\phi}^{\ast}}^{1}\mathcal{H}\left(
\bm{n},\bm{\varphi},\bm{\theta}\right)  -\mathcal{E}\left(
\bm{n},\bm{\varphi},\bm{\theta}\right)  \partial_{\bm{\phi}^{\ast}}^{1}%
S_{N}\left(  \bm{n}\right)  \right\}  =\left[  \partial_{\bm{\phi}}%
^{1}\mathcal{E}\left(  \bm{n},\bm{\varphi},\bm{\theta}\right)  \right]
^{\ast}\\
\partial_{\bm{\theta}}^{1}\mathcal{E}\left(
\bm{n},\bm{\varphi },\bm{\theta}\right)   &  =\frac{1}{S_{N}\left(
\bm{n}\right)  }\left\{  \partial_{\bm{\theta}}^{1}\mathcal{H}\left(
\bm{n},\bm{\varphi},\bm{\theta}\right)  -\mathcal{E}\left(
\bm{n},\bm{\varphi},\bm{\theta}\right)  \partial_{\bm{\theta}}^{1}S_{N}\left(
\bm{n}\right)  \right\}
\end{align}
and
\end{subequations}
\begin{subequations}
\begin{align}
\partial_{\bm{\varphi}^{\ast}}^{1}\mathcal{E}\left(
\bm{n},\bm{\varphi},\bm{\theta}\right)  \cdot\delta\bm{\varphi}^{\ast}  &
=\sum_{1\leq j\leq2s}\partial_{\varphi_{j}^{\ast}}^{1}\mathcal{E}\left(
\bm{n},\bm{\varphi},\bm{\theta}\right)  \left(  \delta\varphi_{j}^{\ast
}\right)  =\sum_{1\leq j\leq2s}\int\frac{\delta\mathcal{E}\left(
\bm{n},\bm{\varphi},\bm{\theta}\right)  }{\delta\varphi_{j}^{\ast}\left(
\bm{r,}\zeta\right)  }\delta\varphi_{j}^{\ast}\left(  \bm{r,}\zeta\right)
d\bm{r}d\zeta\label{funcder}\\
\partial_{\bm{n}}^{1}\mathcal{E}\left(
\bm{n},\bm{\varphi },\bm{\theta}\right)  \cdot\delta\bm{n}  &  =\sum_{1\leq
j\leq s}\frac{\partial\mathcal{E}\left(
\bm{n},\bm{\varphi},\bm{\theta }\right)  }{\partial n_{j}}\delta n_{j}\\
\partial_{\bm{\theta}}^{1}\mathcal{E}\left(
\bm{n},\bm{\varphi },\bm{\theta}\right)  \cdot\delta\bm{\theta}  &
=\sum_{1\leq j\leq s}\frac{\partial\mathcal{E}\left(
\bm{n},\bm{\varphi},\bm{\theta }\right)  }{\partial\theta_{j}}\delta\theta
_{j}.
\end{align}
If the variations $\left\{  \delta\bm{\varphi}^{\ast},\delta
\bm{\varphi}\right\}  $ maintain the orthonormality of the NGSO's i.e.
\end{subequations}
\begin{equation}
\left\langle \left.  \varphi_{j}\right\vert \varphi_{k}\right\rangle
=\delta_{jk}%
\end{equation}
then
\begin{subequations}
\label{fderiv2}%
\begin{align}
\partial_{\bm{\phi}^{\ast}}^{1}\mathcal{E}\left(
\bm{n},\bm{\varphi},\bm{\theta}\right)   &  =\frac{1}{S_{N}\left(
\bm{n}\right)  }\partial_{\bm{\phi}^{\ast}}^{1}\mathcal{H}\left(
\bm{n},\bm{\varphi},\bm{\theta}\right) \\
\partial_{\bm{\theta}}^{1}\mathcal{E}\left(
\bm{n},\bm{\varphi },\bm{\theta}\right)   &  =\frac{1}{S_{N}\left(
\bm{n}\right)  }\partial_{\bm{\theta}}^{1}\mathcal{H}\left(
\bm{n},\bm{\varphi },\bm{\theta}\right)
\end{align}
as $\partial_{\bm{\varphi}^{\ast}}^{1}S_{N}\left(  \bm{n}\right)
=\partial_{\bm{\varphi}}^{1}S_{N}\left(  \bm{n}\right)  =\partial
_{\bm{\theta}}^{1}S_{N}\left(  \bm{n}\right)  =\bm{0}$.

Variation in terms of the variables $\left\{  \bm{n},\bm{\varphi
},\bm{\theta}\right\}  $ leads to a constrained variational problem, which can
be avoided if one uses the group parameters $\left\{
\bm{\eta },\bm{\lambda},\bm{\theta}\right\}  $ and the energy functional
\end{subequations}
\begin{equation}
\mathcal{E}\left(  \bm{\eta},\bm{\lambda},\bm{\theta}\right)  =\frac
{\mathcal{H}\left(  \bm{\eta},\bm{\lambda},\bm{\theta }\right)  }{S_{N}\left(
\bm{\eta}\right)  }=\frac{\left\langle \left.  g\left(
\bm{\eta},\bm{\lambda},\bm{\theta}\right)  ^{N}\right\vert Hg\left(
\bm{\eta},\bm{\lambda},\bm{\theta}\right)  ^{N}\right\rangle }{S_{N}\left(
\bm{\eta}\right)  }%
\end{equation}
and considers the derivatives
\begin{subequations}
\label{GFderiv}%
\begin{align}
\partial_{\bm{\eta}}^{1}\mathcal{E}\left(
\bm{\eta},\bm{\lambda },\bm{\theta}\right)   &  =\frac{1}{S_{N}\left(
\bm{\eta}\right)  }\left\{  \partial_{\bm{\eta}}^{1}\mathcal{H}\left(
\bm{\eta },\bm{\lambda},\bm{\theta}\right)  -\mathcal{E}\left(
\bm{\eta },\bm{\lambda},\bm{\theta}\right)  \partial_{\bm{\eta}}^{1}%
S_{N}\left(  \bm{\eta}\right)  \right\} \\
\partial_{\bm{\lambda}}^{1}\mathcal{E}\left(
\bm{\eta},\bm{\lambda},\bm{\theta}\right)   &  =\frac{1}{S_{N}\left(
\bm{\eta}\right)  }\partial_{\bm{\lambda}}^{1}\mathcal{H}\left(
\bm{\eta},\bm{\lambda},\bm{\theta}\right) \\
\partial_{\bm{\theta}}^{1}\mathcal{E}\left(
\bm{\eta},\bm{\lambda },\bm{\theta}\right)   &  =\frac{1}{S_{N}\left(
\bm{\eta}\right)  }\partial_{\bm{\theta}}^{1}\mathcal{H}\left(
\bm{\eta},\bm{\lambda},\bm{\theta}\right)  .
\end{align}
The generators of the GSO and the phase variations produce
\end{subequations}
\begin{equation}
\partial_{\bm{\lambda}_{jk}}^{1}\mathcal{E}\left(
\bm{\eta },\bm{\lambda},\bm{\theta}\right)  =\frac{i}{S_{N}\left(
\bm{\eta}\right)  }\left\langle \left.  g\left(
\bm{\eta },\bm{\lambda},\bm{\theta}\right)  ^{N}\right\vert \left[
H,a_{j}^{\dagger}a_{k}\right]  g\left(
\bm{\eta},\bm{\lambda },\bm{\theta}\right)  ^{N}\right\rangle ;\quad1\leq
j\neq k\pm s\leq2s \label{Dlamda}%
\end{equation}
and
\begin{equation}
\partial_{\bm{\theta}_{j}}^{1}\mathcal{E}\left(
\bm{\eta },\bm{\lambda},\bm{\theta}\right)  =\frac{i}{S_{N}\left(
\bm{\eta}\right)  }\left\langle \left.  g\left(
\bm{\eta },\bm{\lambda},\bm{\theta}\right)  ^{N}\right\vert \left[  H,\left(
a_{j}^{\dagger}a_{j}+a_{j+s}^{\dagger}a_{j+s}\right)  \right]  g\left(
\bm{\eta},\bm{\lambda},\bm{\theta}\right)  ^{N}\right\rangle ;\quad1\leq j\leq
s, \label{Dtheta}%
\end{equation}
while the generators of the occupation number variations give the derivatives
\begin{align}
\partial_{\bm{\eta}_{j}}^{1}\mathcal{E}\left(
\bm{\eta},\bm{\lambda},\bm{\theta}\right)   &  =\frac{1}{S_{N}\left(
\bm{\eta}\right)  }\left\{  \left\langle \left.  g\left(
\bm{\eta },\bm{\lambda},\bm{\theta}\right)  ^{N}\right\vert \left[  H,\left(
a_{j}^{\dagger}a_{j}+a_{j+s}^{\dagger}a_{j+s}\right)  \right]  _{+}g\left(
\bm{\eta},\bm{\lambda},\bm{\theta}\right)  ^{N}\right\rangle \right.
\nonumber\\
&  -\left.  2\mathcal{E}\left(  \bm{\eta},\bm{\lambda},\bm{\theta }\right)
\left\langle \left.  g\left(  \bm{\eta},\bm{\lambda },\bm{\theta}\right)
^{N}\right\vert \left(  a_{j}^{\dagger}a_{j}+a_{j+s}^{\dagger}a_{j+s}\right)
g\left(  \bm{\eta},\bm{\lambda },\bm{\theta}\right)  ^{N}\right\rangle
\right\}  ;\quad1\leq j\leq s. \label{Deta}%
\end{align}


\section{Stationarity Conditions for an Optimized AGP\label{StatAGPCond}}

In this appendix we analyze the derivatives considered in Appendix
\ref{variational} at stationary points of the AGP\ energy functional and
obtain some important relationships between the total energy and the
occupation numbers of the AGP\ state.

At a stationary point of the unconstrained variation wrt $\left(
\bm{\eta},\bm{\lambda},\bm{\theta}\right)  $
\begin{equation}
\partial_{\bm{\theta}_{j}}^{1}\mathcal{E}\left(
\bm{\eta },\bm{\lambda},\bm{\theta}\right)  =\frac{i}{S_{N}\left(
\bm{\eta}\right)  }\left\langle \left.  g\left(
\bm{\eta },\bm{\lambda},\bm{\theta}\right)  ^{N}\right\vert \left[  H,\left(
a_{j}^{\dagger}a_{j}+a_{j+s}^{\dagger}a_{j+s}\right)  \right]  g\left(
\bm{\eta},\bm{\lambda},\bm{\theta}\right)  ^{N}\right\rangle =0;\quad1\leq
j\leq s, \label{s1}%
\end{equation}
and%
\begin{equation}
\partial_{\bm{\lambda}_{jk}}^{1}\mathcal{E}\left(
\bm{\eta },\bm{\lambda},\bm{\theta}\right)  =\frac{i}{S_{N}\left(
\bm{\eta}\right)  }\left\langle \left.  g\left(
\bm{\eta },\bm{\lambda},\bm{\theta}\right)  ^{N}\right\vert \left[
H,a_{j}^{\dagger}a_{k}\right]  g\left(
\bm{\eta},\bm{\lambda },\bm{\theta}\right)  ^{N}\right\rangle =0;\quad1\leq
j\neq k\pm s\leq2s. \label{s2}%
\end{equation}
Eqs. $\left(  \ref{s1}\right)  $ and $\left(  \ref{s2}\right)  $ combined with
the commutator relationships deduced from the invariance group (see eq.
$\left(  \ref{iso}\right)  $) i.e.
\begin{align}
&  \left\langle \left.  g\left(  \bm{\eta},\bm{\lambda},\bm{\theta }\right)
^{N}\right\vert \left[  H,\left(  a_{j}^{\dagger}a_{j}-a_{j+s}^{\dagger
}a_{j+s}\right)  \right]  g\left(  \bm{\eta },\bm{\lambda},\bm{\theta}\right)
^{N}\right\rangle \nonumber\\
&  \qquad\qquad\qquad\qquad\qquad\left.  =\right.  \left\langle \left.
g\left(  \bm{\eta},\bm{\lambda},\bm{\theta}\right)  ^{N}\right\vert \left[
H,a_{j}^{\dagger}a_{j+s}\right]  g\left(
\bm{\eta },\bm{\lambda},\bm{\theta}\right)  ^{N}\right\rangle \nonumber\\
&  \qquad\qquad\qquad\qquad\qquad\left.  =\right.  \left\langle \left.
g\left(  \bm{\eta},\bm{\lambda},\bm{\theta}\right)  ^{N}\right\vert \left[
H,a_{j+s}^{\dagger}a_{j}\right]  g\left(
\bm{\eta },\bm{\lambda},\bm{\theta}\right)  ^{N}\right\rangle \left.
=\right.  0
\end{align}
give
\begin{equation}
\left\langle \left.  g\left(  \bm{\eta},\bm{\lambda},\bm{\theta }\right)
^{N}\right\vert \left[  H,a_{j}^{\dagger}a_{k}\right]  g\left(
\bm{\eta},\bm{\lambda},\bm{\theta}\right)  ^{N}\right\rangle =0;\quad1\leq
j,k\leq2s. \label{s3}%
\end{equation}
The stationarity of the occupation number variations gives
\begin{align}
\partial_{\bm{\eta}_{j}}^{1}\mathcal{E}\left(
\bm{\eta},\bm{\lambda},\bm{\theta}\right)   &  =\frac{1}{S_{N}\left(
\bm{\eta}\right)  }\left\{  \left\langle \left.  g\left(
\bm{\eta },\bm{\lambda},\bm{\theta}\right)  ^{N}\right\vert \left[  H,\left(
a_{j}^{\dagger}a_{j}+a_{j+s}^{\dagger}a_{j+s}\right)  \right]  _{+}g\left(
\bm{\eta},\bm{\lambda},\bm{\theta}\right)  ^{N}\right\rangle \right.
\nonumber\\
&  -\left.  2\mathcal{E}\left(  \bm{\eta},\bm{\lambda},\bm{\theta }\right)
\left\langle \left.  g\left(  \bm{\eta},\bm{\lambda },\bm{\theta}\right)
^{N}\right\vert \left(  a_{j}^{\dagger}a_{j}+a_{j+s}^{\dagger}a_{j+s}\right)
g\left(  \bm{\eta},\bm{\lambda },\bm{\theta}\right)  ^{N}\right\rangle
\right\}  =0;\quad1\leq j\leq s,
\end{align}
which combined with eq. $\left(  \ref{s3}\right)  $ and noting that
\begin{equation}
\left(  a_{j}^{\dagger}a_{j}+a_{j+s}^{\dagger}a_{j+s}\right)  \left\vert
g\left(  \bm{\eta},\bm{\lambda},\bm{\theta}\right)  ^{N}\right\rangle
=2a_{j}^{\dagger}a_{j}\left\vert g\left(
\bm{\eta },\bm{\lambda},\bm{\theta}\right)  ^{N}\right\rangle
\end{equation}
leads to
\begin{equation}
\left\langle \left.  g\left(  \bm{\eta},\bm{\lambda},\bm{\theta }\right)
^{N}\right\vert Ha_{j}^{\dagger}a_{j}g\left(
\bm{\eta },\bm{\lambda},\bm{\theta}\right)  ^{N}\right\rangle =\mathcal{E}%
\left(  \bm{\eta},\bm{\lambda},\bm{\theta}\right)  \left\langle \left.
g\left(  \bm{\eta},\bm{\lambda},\bm{\theta}\right)  ^{N}\right\vert
a_{j}^{\dagger}a_{j}g\left(  \bm{\eta},\bm{\lambda },\bm{\theta}\right)
^{N}\right\rangle ;\quad1\leq j\leq2s,
\end{equation}
which can be written as
\begin{equation}
\left\langle \left.  g\left(  \bm{\eta},\bm{\lambda},\bm{\theta }\right)
^{N}\right\vert Ha_{j}^{\dagger}a_{j}g\left(
\bm{\eta },\bm{\lambda},\bm{\theta}\right)  ^{N}\right\rangle =\mathcal{E}%
\left(  \bm{\eta},\bm{\lambda},\bm{\theta}\right)  N_{j}.
\end{equation}


\section{AGP Excitation Operators\label{q_op}}

In this appendix we describe excitation operators for the AGP\ state that are
analogous to particle-hole operators for the IPS.

One can define operators $\left\{  \left.  q_{kj}^{\dagger}\right\vert 1\leq
j\neq k\pm s\leq2s\right\}  $ by
\begin{equation}
q_{jk}^{\dagger}\left(  \bm{\eta,\psi}\right)  =c_{k}\left(
\bm{\eta }\right)  a_{j}^{\dagger}a_{k}-\sigma\left(  j\right)  \sigma\left(
k\right)  c_{j}\left(  \bm{\eta}\right)  a_{\left\vert k\right\vert }%
^{\dagger}a_{\left\vert j\right\vert }\text{ for }n_{j}\neq n_{k},
\end{equation}
where%
\begin{equation}
\left\vert j\right\vert =\left\{
\begin{array}
[c]{c}%
j;\qquad1\leq j\leq s\\
j-s;\qquad s+1\leq j\leq2s
\end{array}
\right.  ,
\end{equation}%
\begin{equation}
\sigma\left(  j\right)  =\left\{
\begin{array}
[c]{c}%
-1;\qquad1\leq j\leq s\quad\quad\\
\quad1;\qquad s+1\leq j\leq2s
\end{array}
\right.  ,
\end{equation}
the operators $\left\{  a_{j}^{\dagger}\right\}  $ refer to basis formed by
the optimal CGSO's $\left\{  \left\vert \psi_{j}\right\rangle \right\}  $ and
$\left\{  c_{j}\right\}  $ are the canonical coefficients introduced in
section \ref{AGPstate}. The importance of these operators lies in the
observation that their adjoints $\left\{  \left.  q_{kj}\right\vert 1\leq
j\neq k\pm s\leq2s\right\}  $ have the special annihilation property%
\begin{equation}
q_{jk}\left\vert g\left(  \bm{\eta},\bm{\varphi}\right)  ^{N}\right\rangle
=0;\quad1\leq j\neq k\pm s\leq2s.
\end{equation}
One can show \cite{Goscinski82a} that
\begin{equation}
\left(
\begin{array}
[c]{c}%
q_{jk}^{\dagger}\\
q_{j+sk}^{\dagger}\\
q_{jk+s}^{\dagger}\\
q_{j+sk+s}^{\dagger}\\
q_{jk}\\
q_{j+sk}\\
q_{jk+s}\\
q_{j+sk+s}%
\end{array}
\right)  =\left(
\begin{array}
[c]{cccccccc}%
c_{k} & 0 & \cdots & \cdots & \cdots & \cdots & 0 & -c_{j}\\
0 & c_{k} & 0 & \cdots & \cdots & 0 & c_{j} & 0\\
\vdots & 0 & c_{k} & 0 & 0 & c_{j} & 0 & \vdots\\
\vdots & \vdots & 0 & c_{k} & -c_{j} & 0 & \vdots & \vdots\\
\vdots & \vdots & 0 & -c_{j} & c_{k} & 0 & \vdots & \vdots\\
\vdots & 0 & c_{j} & 0 & 0 & c_{k} & 0 & \vdots\\
0 & c_{j} & 0 & \cdots & \cdots & 0 & c_{k} & 0\\
-c_{j} & 0 & \cdots & \cdots & \cdots & \cdots & 0 & c_{k}%
\end{array}
\right)  \left(
\begin{array}
[c]{c}%
a_{j}^{\dagger}a_{k}\\
a_{j}^{\dagger}a_{k+s}\\
a_{j+s}^{\dagger}a_{k}\\
a_{j+s}^{\dagger}a_{k+s}\\
a_{k}^{\dagger}a_{j}\\
a_{k}^{\dagger}a_{j+s}\\
a_{k+s}^{\dagger}a_{j}\\
a_{k+s}^{\dagger}a_{j+s}%
\end{array}
\right)  ,
\end{equation}
leading to the transpose relationship
\begin{align}
&  \left(  q_{jk}^{\dagger},q_{j+sk}^{\dagger},q_{jk+s}^{\dagger}%
,q_{j+sk+s}^{\dagger},q_{jk},q_{j+sk},q_{jk+s},q_{j+sk+s}\right) \nonumber\\
&  \qquad\qquad\qquad\qquad\qquad\left.  =\right.  \left(  a_{j}^{\dagger
}a_{k},a_{j}^{\dagger}a_{k+s},a_{j+s}^{\dagger}a_{k},a_{j+s}^{\dagger}%
a_{k+s},a_{k}^{\dagger}a_{j},a_{k}^{\dagger}a_{j+s},a_{k+s}^{\dagger}%
a_{j},a_{k+s}^{\dagger}a_{j+s}\right)  \mathbb{S}\left(  jk\right)
\end{align}
and that the transformation matrix $\mathbb{S}\left(  jk\right)  $ is
nonsingular and has the form
\begin{equation}
\mathbb{S}\left(  jk\right)  =\left(
\begin{array}
[c]{cc}%
\mathbb{S}_{1}\left(  jk\right)  & \mathbb{S}_{2}\left(  jk\right) \\
\mathbb{S}_{2}\left(  jk\right)  ^{\ast} & \mathbb{S}_{1}\left(  jk\right)
^{\ast}%
\end{array}
\right)  .
\end{equation}
.

In the case of equal coefficients and geminal occupation numbers i.e.
$c_{j}=c_{k}\Leftrightarrow n_{j}=n_{k}$ one can use the following nonsingular
transformation
\begin{equation}
\left(
\begin{array}
[c]{c}%
u_{jk}\\
u_{j+sk}\\
u_{jk+s}\\
u_{j+sk+s}\\
v_{jk}\\
v_{j+sk}\\
v_{jk+s}\\
v_{j+sk+s}%
\end{array}
\right)  =\left(
\begin{array}
[c]{cccccccc}%
c_{k} & 0 & \cdots & \cdots & \cdots & \cdots & 0 & c_{j}\\
0 & c_{k} & 0 & \cdots & \cdots & 0 & -c_{j} & 0\\
\vdots & 0 & c_{k} & 0 & 0 & -c_{j} & 0 & \vdots\\
\vdots & \vdots & 0 & c_{k} & c_{j} & 0 & \vdots & \vdots\\
\vdots & \vdots & 0 & -c_{j} & c_{k} & 0 & \vdots & \vdots\\
\vdots & 0 & c_{j} & 0 & 0 & c_{k} & 0 & \vdots\\
0 & c_{j} & 0 & \cdots & \cdots & 0 & c_{k} & 0\\
-c_{j} & 0 & \cdots & \cdots & \cdots & \cdots & 0 & c_{k}%
\end{array}
\right)  \left(
\begin{array}
[c]{c}%
a_{j}^{\dagger}a_{k}\\
a_{j}^{\dagger}a_{k+s}\\
a_{j+s}^{\dagger}a_{k}\\
a_{j+s}^{\dagger}a_{k+s}\\
a_{k}^{\dagger}a_{j}\\
a_{k}^{\dagger}a_{j+s}\\
a_{k+s}^{\dagger}a_{j}\\
a_{k+s}^{\dagger}a_{j+s}%
\end{array}
\right)  ,
\end{equation}
leading to the transpose relationship
\begin{align}
&  \left(  u_{jk},u_{j+sk},u_{jk+s},u_{j+sk+s},v_{jk},v_{j+sk},v_{jk+s}%
,v_{j+sk+s}\right) \nonumber\\
&  \qquad\qquad\qquad\qquad\qquad\left.  =\right.  \left(  a_{j}^{\dagger
}a_{k},a_{j}^{\dagger}a_{k+s},a_{j+s}^{\dagger}a_{k},a_{j+s}^{\dagger}%
a_{k+s},a_{k}^{\dagger}a_{j},a_{k}^{\dagger}a_{j+s},a_{k+s}^{\dagger}%
a_{j},a_{k+s}^{\dagger}a_{j+s}\right)  \mathbb{\tilde{S}}\left(  jk\right)
\end{align}
and operators that have the properties
\begin{equation}
v_{ab}\left\vert g\left(  \bm{\eta},\bm{\varphi}\right)  ^{N}\right\rangle
=0,\quad u_{ab}\left\vert g\left(  \bm{\eta },\bm{\varphi}\right)
^{N}\right\rangle \neq0;\quad a=j,j+s;b=k,k+s.
\end{equation}


\section{\bigskip Fock Operator and the Total Energy\label{Fock&Energy}}

In this appendix we first describe a general property of the derivatives of
polynomial functions, then use it to obtain an expression for the total energy
in terms of one electron properties.

\subsection{Homogeneous Functions\label{HomFunc}}

Polynomial functions can be expanded as a sum of homogeneous mononomial
functions as%

\begin{equation}
f\left(  \bm{x}\right)  =f_{0}+f_{1}\left(  \bm{x}\right)  +\cdots
+f_{n}\left(  \bm{x}\right)  =%
%TCIMACRO{\tsum \limits_{0\leq m\leq n}}%
%BeginExpansion
{\textstyle\sum\limits_{0\leq m\leq n}}
%EndExpansion
f_{m}\left(  \bm{x}\right)  ,
\end{equation}
where
\begin{equation}
f_{m}\left(  \bm{x}\right)  =%
%TCIMACRO{\tsum \limits_{1\leq j_{1},\cdots,j_{m}\leq r}}%
%BeginExpansion
{\textstyle\sum\limits_{1\leq j_{1},\cdots,j_{m}\leq r}}
%EndExpansion
c_{j_{1}\cdots j_{m}}x_{j_{1}}\cdots x_{j_{m}}.
\end{equation}
The variables $\bm{x}$ can be complex and in particular $\bm{x=}\left(
\bm{z}^{\ast}\bm{,z}\right)  $, where we consider $\bm{z}^{\ast}\bm{,z}$ as
independent variables. Consider
\begin{equation}
\frac{\partial f_{m}}{\partial x_{l}}\left(  \bm{x}\right)  =%
%TCIMACRO{\tsum \limits_{1\leq j_{1},\cdots,j_{m-1}\neq l\leq r}}%
%BeginExpansion
{\textstyle\sum\limits_{1\leq j_{1},\cdots,j_{m-1}\neq l\leq r}}
%EndExpansion
c_{j_{1}\cdots j_{l}\cdots j_{m}}x_{j_{1}}\cdots x_{j_{m-1}}%
\end{equation}
then
\begin{equation}
x_{l}\frac{\partial f_{m}}{\partial x_{l}}\left(  \bm{x}\right)  =%
%TCIMACRO{\tsum \limits_{1\leq j_{1},\cdots,j_{m-1}\neq l\leq r}}%
%BeginExpansion
{\textstyle\sum\limits_{1\leq j_{1},\cdots,j_{m-1}\neq l\leq r}}
%EndExpansion
c_{j_{1}\cdots j_{l}\cdots j_{m}}x_{j_{1}}\cdots x_{j_{m-1}}x_{l},
\end{equation}
which gives
\begin{equation}%
%TCIMACRO{\tsum \limits_{1\leq l\leq r}}%
%BeginExpansion
{\textstyle\sum\limits_{1\leq l\leq r}}
%EndExpansion
x_{l}\frac{\partial f_{m}}{\partial x_{l}}\left(  \bm{x}\right)
=mf_{m}\left(  \bm{x}\right)  .
\end{equation}
Hence%
\begin{align}%
%TCIMACRO{\tsum \limits_{1\leq l\leq r}}%
%BeginExpansion
{\textstyle\sum\limits_{1\leq l\leq r}}
%EndExpansion
x_{l}\frac{\partial f}{\partial x_{l}}\left(  \bm{x}\right)   &  =%
%TCIMACRO{\tsum \limits_{1\leq l\leq r}}%
%BeginExpansion
{\textstyle\sum\limits_{1\leq l\leq r}}
%EndExpansion
x_{l}\frac{\partial}{\partial x_{l}}\left\{
%TCIMACRO{\tsum \limits_{0\leq m\leq n}}%
%BeginExpansion
{\textstyle\sum\limits_{0\leq m\leq n}}
%EndExpansion
f_{m}\left(  \bm{x}\right)  \right\}  =%
%TCIMACRO{\tsum \limits_{\substack{1\leq l\leq r\\0\leq m\leq n}}}%
%BeginExpansion
{\textstyle\sum\limits_{\substack{1\leq l\leq r\\0\leq m\leq n}}}
%EndExpansion
x_{l}\frac{\partial f_{m}\left(  \bm{x}\right)  }{\partial x_{l}}\nonumber\\
&  =%
%TCIMACRO{\tsum \limits_{\substack{0\leq m\leq n\\1\leq l\leq r}}}%
%BeginExpansion
{\textstyle\sum\limits_{\substack{0\leq m\leq n\\1\leq l\leq r}}}
%EndExpansion
x_{l}\frac{\partial f_{m}\left(  \bm{x}\right)  }{\partial x_{l}}=%
%TCIMACRO{\tsum \limits_{0\leq m\leq n}}%
%BeginExpansion
{\textstyle\sum\limits_{0\leq m\leq n}}
%EndExpansion
mf_{m}\left(  \bm{x}\right)  =f_{1}\left(  \bm{x}\right)  +2f_{2}\left(
\bm{x}\right)  +\cdots+nf_{n}\left(  \bm{x}\right)  . \label{homex}%
\end{align}


\subsection{Fock Operator and Energy}

The energy of \emph{any} $2N$-electron state is given by
\begin{equation}
\mathcal{E}\left(  \mathbb{D}^{2},\bm{\varphi}\right)  =\operatorname{Tr}%
\left\{  \mathbb{H}_{1}\left(  \bm{\varphi}\right)  \mathbb{D}^{1}\right\}
+\frac{1}{4}\operatorname{Tr}\left\{  \mathbb{V}_{2}\left(
\bm{\varphi}\right)  \mathbb{D}^{2}\right\}  \equiv\mathcal{E}_{1}\left(
\mathbb{D}^{1},\bm{\varphi}\right)  +\mathcal{E}_{2}\left(  \mathbb{D}%
^{2},\bm{\varphi}\right)  ,
\end{equation}
where
\begin{equation}
\mathcal{E}_{1}\left(  \mathbb{D}^{1},\bm{\varphi}\right)  =\operatorname{Tr}%
\left\{  \mathbb{H}_{1}\left(  \bm{\varphi}\right)  \mathbb{D}^{1}\right\}
\end{equation}
and
\begin{equation}
\mathcal{E}_{2}\left(  \mathbb{D}^{2},\bm{\varphi}\right)  =\frac{1}%
{4}\operatorname{Tr}\left\{  \mathbb{V}_{2}\left(  \bm{\varphi}\right)
\mathbb{D}^{2}\right\}  .
\end{equation}
As
\begin{equation}
\frac{\delta\mathcal{E}\left(  \mathbb{D}^{2},\bm{\varphi}\right)  }%
{\delta\varphi_{j}^{\ast}}=F\left(  \mathbb{D}^{2},\bm{\varphi}\right)
\left\vert \varphi_{j}\right\rangle
\end{equation}
one obtains
\begin{equation}%
%TCIMACRO{\tsum \limits_{1\leq j\leq2s}}%
%BeginExpansion
{\textstyle\sum\limits_{1\leq j\leq2s}}
%EndExpansion
\int\varphi_{j}^{\ast}\left(  \bm{r,\xi}\right)  \frac{\delta\mathcal{E}%
\left(  \mathbb{D}^{2},\bm{\varphi}\right)  }{\delta\varphi_{j}^{\ast}\left(
\bm{r,\xi}\right)  }d\bm{r}d\bm{\xi=}%
%TCIMACRO{\tsum \limits_{1\leq j\leq2s}}%
%BeginExpansion
{\textstyle\sum\limits_{1\leq j\leq2s}}
%EndExpansion
\left\langle \left.  \varphi_{j}\right\vert F\left(  \mathbb{D}^{2}%
,\bm{\varphi}\right)  \varphi_{j}\right\rangle =\operatorname{Tr}\left\{
F\left(  \mathbb{D}^{2},\bm{\varphi}\right)  \right\}  .
\end{equation}
giving%
\begin{align}
\operatorname{Tr}\left\{  F\left(  \mathbb{D}^{2},\bm{\varphi}\right)
\right\}   &  =%
%TCIMACRO{\tsum \limits_{1\leq j\leq2s}}%
%BeginExpansion
{\textstyle\sum\limits_{1\leq j\leq2s}}
%EndExpansion
\int\varphi_{j}^{\ast}\left(  \bm{r,\xi}\right)  \frac{\delta\mathcal{E}%
\left(  \mathbb{D}^{2},\bm{\varphi}\right)  }{\delta\varphi_{j}^{\ast}\left(
\bm{r,\xi}\right)  }d\bm{r}d\bm{\xi}\\
&  \bm{=}%
%TCIMACRO{\tsum \limits_{1\leq j\leq2s}}%
%BeginExpansion
{\textstyle\sum\limits_{1\leq j\leq2s}}
%EndExpansion
\int\varphi_{j}^{\ast}\left(  \bm{r,\xi}\right)  \left\{  \frac{\delta
\mathcal{E}_{1}\left(  \mathbb{D}^{2},\bm{\varphi}\right)  }{\delta\varphi
_{j}^{\ast}\left(  \bm{r,\xi}\right)  }+\frac{\delta\mathcal{E}_{2}\left(
\mathbb{D}^{2},\bm{\varphi}\right)  }{\delta\varphi_{j}^{\ast}\left(
\bm{r,\xi}\right)  }\right\}  d\bm{r}d\bm{\xi},
\end{align}
which by eq. $\left(  \ref{homex}\right)  $ produces
\begin{equation}
\operatorname{Tr}\left\{  F\left(  \mathbb{D}^{2},\bm{\varphi}\right)
\right\}  =\mathcal{E}_{1}\left(  \mathbb{D}^{2},\bm{\varphi}\right)
+2\mathcal{E}_{2}\left(  \mathbb{D}^{2},\bm{\varphi}\right)  .
\end{equation}
Noting that
\begin{equation}
\mathcal{E}\left(  \mathbb{D}^{2},\bm{\varphi}\right)  =\mathcal{E}_{1}\left(
\mathbb{D}^{2},\bm{\varphi}\right)  +\mathcal{E}_{2}\left(  \mathbb{D}%
^{2},\bm{\varphi}\right)
\end{equation}
one obtains
\begin{equation}
\operatorname{Tr}\left\{  F\left(  \mathbb{D}^{2},\bm{\varphi}\right)
\right\}  =\left\{  \mathcal{E}\left(  \mathbb{D}^{2},\bm{\varphi}\right)
+\mathcal{E}_{2}\left(  \mathbb{D}^{2},\bm{\varphi}\right)  \right\}
=\left\{  2\mathcal{E}\left(  \mathbb{D}^{2},\bm{\varphi}\right)
-\mathcal{E}_{1}\left(  \mathbb{D}^{2},\bm{\varphi}\right)  \right\}  .
\end{equation}
Hence the important relationships
\begin{subequations}
\label{FE}%
\begin{align}
\mathcal{E}\left(  \mathbb{D}^{2},\bm{\varphi}\right)   &  =\operatorname{Tr}%
\left\{  F\left(  \mathbb{D}^{2},\bm{\varphi}\right)  \right\}  -\mathcal{E}%
_{2}\left(  \mathbb{D}^{2},\bm{\varphi}\right) \label{FEa}\\
\mathcal{E}\left(  \mathbb{D}^{2},\bm{\varphi}\right)   &  =\frac{1}%
{2}\left\{  \overset{}{\operatorname{Tr}\left\{  F\left(  \mathbb{D}%
^{2},\bm{\varphi}\right)  \right\}  +\mathcal{E}_{1}\left(  \mathbb{D}%
^{2},\bm{\varphi}\right)  }\right\}  . \label{FEb}%
\end{align}


Recalling eq. $\left(  \ref{Fockelem}\right)  $
\end{subequations}
\begin{equation}
F\left(  \mathbb{D}^{2},\bm{\varphi}\right)  _{qp}=\frac{1}{\operatorname{Tr}%
\left\{  D^{2N}\right\}  }\operatorname{Tr}\left\{  \overset{}{\left\{
a_{p}^{\dagger}a_{q}H-a_{q}Ha_{p}^{\dagger}\right\}  }D^{2N}\right\}
\end{equation}
one observes that for \emph{any }state $D^{2N}$
\begin{align}
\operatorname{Tr}\left\{  F\left(  \mathbb{D}^{2},\bm{\varphi}\right)
\right\}   &  =\sum_{1\leq p\leq2s}\frac{1}{\operatorname{Tr}\left\{
D^{2N}\right\}  }\operatorname{Tr}\left\{  \overset{}{\left\{  a_{p}^{\dagger
}a_{p}H-a_{p}Ha_{p}^{\dagger}\right\}  }D^{2N}\right\} \nonumber\\
&  =\frac{1}{\operatorname{Tr}\left\{  D^{2N}\right\}  }\operatorname{Tr}%
\left\{  \mathfrak{N}HD^{2N}\right\}  -\sum_{1\leq p\leq2s}N_{p}\left(
\mathbb{D}^{2},\bm{\varphi}\right)  \frac{\operatorname{Tr}\left\{
a_{p}Ha_{p}^{\dagger}D^{2N}\right\}  }{\operatorname{Tr}\left\{
a_{p}^{\dagger}a_{p}D^{2N}\right\}  }, \label{TrF1}%
\end{align}
where
\begin{equation}
\mathfrak{N=}\sum_{1\leq p\leq2s}a_{p}^{\dagger}a_{p}%
\end{equation}
is the number operator and
\begin{equation}
N_{p}\left(  \mathbb{D}^{2},\bm{\varphi}\right)  =\frac{\operatorname{Tr}%
\left\{  a_{p}^{\dagger}a_{p}D^{2N}\right\}  }{\operatorname{Tr}\left\{
D^{2N}\right\}  }%
\end{equation}
is the diagonal element of the FORDM of $D^{2N}$ in basis $\left\{  \left.
\varphi_{j}\right\vert 1\leq j\leq2s\right\}  $ i.e. the occupation of the GSO
$\left\vert \varphi_{p}\right\rangle .$ Eq. $\left(  \ref{TrF1}\right)  $ can
be further developed to give
\begin{equation}
\operatorname{Tr}\left\{  F\left(  \mathbb{D}^{2},\bm{\varphi}\right)
\right\}  =2N\mathcal{E}\left(  \mathbb{D}^{2},\bm{\varphi}\right)
-\sum_{1\leq p\leq2s}N_{p}\left(  \mathbb{D}^{2},\bm{\varphi}\right)
\mathcal{E}_{p}\left(  D^{2N},\bm{\varphi}\right)  ,
\end{equation}
where the energy of the $2N$ electron state $D^{2N}$ is
\begin{equation}
\mathcal{E}\left(  \mathbb{D}^{2},\bm{\varphi}\right)  =\frac
{\operatorname{Tr}\left\{  HD^{2N}\right\}  }{\operatorname{Tr}\left\{
D^{2N}\right\}  }%
\end{equation}
and the energy of the $2N-1$ electron state $a_{p}^{\dagger}D^{2N}a_{p}$ is
\begin{equation}
\mathcal{E}_{p}\left(  D^{2N},\bm{\varphi}\right)  =\frac{\operatorname{Tr}%
\left\{  a_{p}Ha_{p}^{\dagger}D^{2N}\right\}  }{\operatorname{Tr}\left\{
a_{p}^{\dagger}a_{p}D^{2N}\right\}  }.
\end{equation}
This gives
\begin{equation}
\operatorname{Tr}\left\{  F\left(  \mathbb{D}^{2},\bm{\varphi}\right)
\right\}  =\sum_{1\leq p\leq2s}N_{p}\left(  \mathbb{D}^{2}%
,\bm{\varphi }\right)  \left\{  \mathcal{E}\left(  \mathbb{D}^{2}%
,\bm{\varphi}\right)  -\mathcal{E}_{p}\left(  \mathbb{D}^{2}%
,\bm{\varphi}\right)  \right\}  =\sum_{1\leq p\leq2s}N_{p}\left(
\mathbb{D}^{2},\bm{\varphi}\right)  I\left(  \mathbb{D}^{2},\varphi
_{p}\right)  , \label{TrF}%
\end{equation}
where the ionization potential associated with the GSO\ $\left\vert
\varphi_{p}\right\rangle $ from the state $D^{2N}$ is
\begin{equation}
I\left(  \mathbb{D}^{2},\varphi_{p}\right)  =\mathcal{E}\left(  \mathbb{D}%
^{2},\bm{\varphi}\right)  -\mathcal{E}_{p}\left(  \mathbb{D}^{2}%
,\bm{\varphi}\right)  .
\end{equation}
Eqs. $\left(  \ref{FE}\right)  $ then give
\begin{subequations}
\begin{align}
\mathcal{E}\left(  \mathbb{D}^{2},\bm{\varphi}\right)   &  =\sum_{1\leq
p\leq2s}N_{p}\left(  \mathbb{D}^{2},\bm{\varphi}\right)  I\left(
\mathbb{D}^{2},\varphi_{p}\right)  -\mathcal{E}_{2}\left(  \mathbb{D}%
^{2},\bm{\varphi}\right) \label{FE1a}\\
\mathcal{E}\left(  \mathbb{D}^{2},\bm{\varphi}\right)   &  =\frac{1}%
{2}\left\{  \sum_{1\leq p\leq2s}N_{p}\left(  \mathbb{D}^{2}%
,\bm{\varphi }\right)  I\left(  \mathbb{D}^{2},\varphi_{p}\right)
+\mathcal{E}_{1}\left(  \mathbb{D}^{2},\bm{\varphi}\right)  \right\}  .
\label{FE1b}%
\end{align}


If $\left\{  \left\vert \varphi_{p}\right\rangle \right\}  $ are the NGSO's of
the state $D^{2N}$ then the one electron energy $\mathcal{E}_{1}\left(
\mathbb{D}^{2},\bm{\varphi}\right)  $ can be expressed as
\end{subequations}
\begin{equation}
\mathcal{E}_{1}\left(  \mathbb{D}^{2},\bm{\varphi}\right)  =%
%TCIMACRO{\tsum \limits_{1\leq p\leq2s}}%
%BeginExpansion
{\textstyle\sum\limits_{1\leq p\leq2s}}
%EndExpansion
N_{p}\left(  \mathbb{D}^{2},\bm{\varphi}\right)  \left\langle \left.
\varphi_{p}\right\vert h_{1}\varphi_{p}\right\rangle
\end{equation}
and we obtain
\begin{equation}
\mathcal{E}\left(  \mathbb{D}^{2},\bm{\varphi}\right)  =\frac{1}{2}%
%TCIMACRO{\tsum \limits_{1\leq l\leq2s}}%
%BeginExpansion
{\textstyle\sum\limits_{1\leq l\leq2s}}
%EndExpansion
N_{p}\left(  \mathbb{D}^{2},\bm{\varphi}\right)  \left\{  I\left(
\mathbb{D}^{2},\varphi_{p}\right)  +\left\langle \left.  \varphi
_{p}\right\vert h_{1}\varphi_{p}\right\rangle \right\}  \equiv\frac{1}{2}%
%TCIMACRO{\tsum \limits_{1\leq p\leq2s}}%
%BeginExpansion
{\textstyle\sum\limits_{1\leq p\leq2s}}
%EndExpansion
N_{p}\left(  \mathbb{D}^{2},\bm{\varphi}\right)  \kappa\left(  \mathbb{D}%
^{2},\varphi_{p}\right)  ,
\end{equation}
where
\begin{equation}
\kappa\left(  \mathbb{D}^{2},\varphi_{p}\right)  =I\left(  \mathbb{D}%
^{2},\varphi_{p}\right)  +\left\langle \left.  \varphi_{p}\right\vert
h_{1}\varphi_{p}\right\rangle .
\end{equation}


\section{Stationary AGP Fock Operator\label{AGPFock}}

In this appendix we consider the generalized Fock operator for the special
case of an optimized AGP state.

The matrix elements $\left\{  \left.  F\left(  g\left(  \bm{\eta
},\bm{\psi}\right)  ^{N}\right)  _{qp}\right\vert 1\leq p,q\leq2s\right\}  $
of the generalized Fock operator are
\begin{align}
F\left(  g\left(  \bm{\eta},\bm{\psi}\right)  ^{N}\right)  _{qp}  &  =\frac
{1}{S_{N}\left(  \bm{\eta}\right)  }\left\langle \left.  g\left(
\bm{\eta},\bm{\psi}\right)  ^{N}\right\vert a_{p}^{\dagger}\left[
a_{q},H\right]  g\left(  \bm{\eta},\bm{\psi}\right)  ^{N}\right\rangle
\nonumber\\
&  =\frac{1}{S_{N}\left(  \bm{\eta}\right)  }\left\langle \left.  g\left(
\bm{\eta},\bm{\psi}\right)  ^{N}\right\vert \left\{  a_{p}^{\dagger}%
a_{q}H-a_{p}^{\dagger}Ha_{q}\right\}  g\left(  \bm{\eta},\bm{\psi }\right)
^{N}\right\rangle .
\end{align}
The diagonal elements of $F\left(  g\left(  \bm{\eta},\bm{\psi
}\right)  ^{N}\right)  $ for an optimized AGP\ have the form
\begin{align}
F\left(  g\left(  \bm{\eta},\bm{\psi}\right)  ^{N}\right)  _{qq}  &
=\frac{N_{q}}{S_{N}\left(  \bm{\eta}\right)  }\left\{  \mathcal{E}\left(
g\left(  \bm{\eta},\bm{\psi}\right)  ^{N}\right)  -\frac{\left\langle \left.
g\left(  \bm{\eta},\bm{\psi}\right)  ^{N}\right\vert a_{q}^{\dagger}%
Ha_{q}g\left(  \bm{\eta},\bm{\psi}\right)  ^{N}\right\rangle }{\left\langle
\left.  g\left(  \bm{\eta},\bm{\psi }\right)  ^{N}\right\vert a_{q}^{\dagger
}a_{q}g\left(  \bm{\eta },\bm{\psi}\right)  ^{N}\right\rangle }\right\}
\nonumber\\
&  =\frac{N_{q}}{S_{N}\left(  \bm{\eta}\right)  }\left\{  \mathcal{E}\left(
g\left(  \bm{\eta},\bm{\psi}\right)  ^{N}\right)  -\mathcal{E}_{q}\left(
g\left(  \bm{\eta},\bm{\psi}\right)  ^{N}\right)  \right\}  ,
\end{align}
where $\mathcal{E}_{q}\left(  g\left(  \bm{\eta},\bm{\psi}\right)
^{N}\right)  $ is the energy of the normalized $2N-1$ electron state formed by
removing an electron described by a CGSO $\left\vert \psi_{q}\right\rangle $
\begin{equation}
\mathcal{E}_{q}\left(  g\left(  \bm{\eta},\bm{\psi}\right)  ^{N}\right)
=\frac{\left\langle \left.  g\left(  \bm{\eta},\bm{\psi }\right)
^{N}\right\vert a_{q}^{\dagger}Ha_{q}g\left(  \bm{\eta },\bm{\psi}\right)
^{N}\right\rangle }{\left\langle \left.  a_{q}g\left(
\bm{\eta},\bm{\psi}\right)  ^{N}\right\vert a_{q}g\left(
\bm{\eta},\bm{\psi}\right)  ^{N}\right\rangle }.
\end{equation}
It is easy to observe that
\begin{equation}
\mathcal{E}_{q+s}\left(  g\left(  \bm{\eta},\bm{\psi}\right)  ^{N}\right)
=\mathcal{E}_{q}\left(  g\left(  \bm{\eta},\bm{\psi }\right)  ^{N}\right)
\end{equation}
and
\begin{equation}
F\left(  g\left(  \bm{\eta},\bm{\psi}\right)  ^{N}\right)  _{qq}=F\left(
g\left(  \bm{\eta},\bm{\psi}\right)  ^{N}\right)  _{q+sq+s};\quad1\leq q\leq
s.
\end{equation}
The off diagonal elements are given by
\begin{equation}
F\left(  g\left(  \bm{\eta},\bm{\psi}\right)  ^{N}\right)  _{qp}=\frac
{1}{S_{N}\left(  \bm{\eta}\right)  }\left\langle \left.  g\left(
\bm{\eta},\bm{\psi}\right)  ^{N}\right\vert a_{p}^{\dagger}Ha_{q}g\left(
\bm{\eta},\bm{\psi}\right)  ^{N}\right\rangle ,
\end{equation}
which are \emph{not }in general zero in the CGSO basis, but are distinct from
the expressions obtained in the general MCSCF case as a result of the
optimized AGP\ Brillouin condition.

\section{A Fock Identity\label{FockID}}

In this appendix we use expression that expresses the total energy in terms of
one electron quantities to derive an expression for the total energy in terms
of singly ionized states.

Using eq. $\left(  \ref{FE1b}\right)  $ we obtain
\begin{subequations}
\begin{equation}
\mathcal{E}\left(  \mathbb{D}^{2},\bm{\varphi}\right)  =N\mathcal{E}\left(
\mathbb{D}^{2},\bm{\varphi}\right)  -\frac{1}{2}\left\{  \sum_{1\leq p\leq
2s}N_{p}\left(  \mathbb{D}^{2},\bm{\varphi}\right)  \mathcal{E}_{p}\left(
\mathbb{D}^{2},\bm{\varphi}\right)  -\mathcal{E}_{1}\left(  \mathbb{D}%
^{2},\bm{\varphi}\right)  \right\}  ,
\end{equation}
which leads to
\end{subequations}
\begin{equation}
\mathcal{E}\left(  \mathbb{D}^{2},\bm{\varphi}\right)  =\frac{1}{2\left(
N-1\right)  }\left\{  \sum_{1\leq p\leq2s}N_{p}\left(  \mathbb{D}%
^{2},\bm{\varphi}\right)  \mathcal{E}_{p}\left(  \mathbb{D}^{2}%
,\bm{\varphi}\right)  -\mathcal{E}_{1}\left(  \mathbb{D}^{2}%
,\bm{\varphi}\right)  \right\}  . \label{EnFunc}%
\end{equation}
Eq. $\left(  \ref{EnFunc}\right)  $ can also be expressed as
\begin{equation}
\mathcal{E}\left(  \mathbb{D}^{2},\bm{\varphi}\right)  =\frac{1}{2\left(
N-1\right)  }\left\{  \sum_{1\leq p\leq2s}N_{p}\left(  \mathbb{D}%
^{2},\bm{\varphi}\right)  \left\{  \mathcal{E}_{p}\left(  \mathbb{D}%
^{2},\bm{\varphi}\right)  -\left\langle \left.  \varphi_{p}\right\vert
h_{1}\varphi_{p}\right\rangle \right\}  \right\}  , \label{EnFuncNew}%
\end{equation}
leading, in the case of NGSO's, to
\begin{equation}
\mathcal{E}\left(  \mathbb{D}^{2},\bm{\varphi}\right)  =\frac{1}{2\left(
N-1\right)  }\sum_{1\leq p\leq2s}N_{p}\left(  \mathbb{D}^{2}%
,\bm{\varphi }\right)  \iota_{p}\left(  \mathbb{D}^{2},\bm{\varphi}\right)  ,
\end{equation}
where
\begin{equation}
\iota_{p}\left(  \mathbb{D}^{2},\bm{\varphi}\right)  =\mathcal{E}_{p}\left(
\mathbb{D}^{2},\bm{\varphi}\right)  -\left\langle \left.  \varphi
_{p}\right\vert h_{1}\varphi_{p}\right\rangle .
\end{equation}
The expression eq. $\left(  \ref{EnFuncNew}\right)  $ leads to
\begin{align}
\frac{\delta\mathcal{E}\left(  \mathbb{D}^{2},\bm{\varphi}\right)  }%
{\delta\varphi_{j}^{\ast}}  &  =F\left(  \mathbb{D}^{2},\bm{\varphi }\right)
\left\vert \varphi_{j}\right\rangle =\sum_{1\leq k\leq2s}\left\vert
\varphi_{k}\right\rangle \varepsilon_{kj}\nonumber\\
&  =h_{1}\left\vert \varphi_{j}\right\rangle +\sum_{1\leq p\leq2s}N_{p}\left(
\mathbb{D}^{2},\bm{\varphi}\right)  \frac{\delta\mathcal{E}_{p}\left(
\mathbb{D}^{2},\bm{\varphi}\right)  }{\delta\varphi_{j}^{\ast}};\quad1\leq
j\leq2s,
\end{align}
where $\left\{  \varepsilon_{kj}\right\}  $ are the Lagrange parameters.

\bibliographystyle{achemso}
\bibliography{agpopt,bw,coleman,DMF,DMF1,yngve}

\end{document}
